% =====================================================================
%  R. Nielsen, MATHEMATICS AND DIALECTIC. A PHILOSOPHICAL TREATISE.
%  Copenhagen: Gyldendal (F. Hegel), 1859.  ENGLISH TRANSLATION, from transcription.tex (Danish) as finished
%  2026-10-06 (before the first commit of the book; see its header).
%  GENERATED by ebook-method/mathematik-og-dialektik/tr_build.py --write from tr_template.tex and the
%  translators' parts in texts/nielsen/mathematik-og-dialektik/.parts/tr/ (TRANSLATOR-BRIEF.md).
%  The sequel to "Philosophie og Mathematik" (1857), translated in this repo as "Philosophy and Mathematics".
%  Formulas are carried over from the Danish unchanged, printer's errors included (marked there PRINTED AS IS).
%  Nine translator agents (2026-10-06); their odd spots were read at the image and the Danish mended
%  (finish.py ODD), and the same changes carried into the parts by hand (pp. 8, 14, 22, 23, 28, 47, 52,
%  75, 76, 85, 89, 99-101, 106, 107, 121, 129, 143); misprints the image confirmed carry a "% print:" comment.
%  Conventions: TRANSLATION-PLAYBOOK.md.  Page numbers of the 1859 print in the margin.
% =====================================================================
\documentclass[12pt,a4paper,openany]{book}
\usepackage[utf8]{inputenc}
\usepackage[T1]{fontenc}
\usepackage{amsmath}
\usepackage{amssymb}
\usepackage{libertinus}
\usepackage{libertinust1math}
\usepackage{textalpha}
\usepackage[english]{babel}
\usepackage{enumitem}
\usepackage[protrusion=true,expansion=true]{microtype}
\usepackage{setspace}
\usepackage[margin=1.3in]{geometry}
\usepackage{fancyhdr}
\usepackage{hyperref}
\hypersetup{pdftitle={Mathematics and Dialectic}, pdfauthor={Rasmus Nielsen}, hidelinks}
\setlength{\headheight}{15pt}
\pagestyle{fancy}
\fancyhf{}
\fancyhead[LE,RO]{\thepage}
\fancyhead[C]{\textit{R. Nielsen: Mathematics and Dialectic (1859)}}
\renewcommand{\thefootnote}{\arabic{footnote})}
\setstretch{1.3}
\usepackage{marginnote}
\newcommand{\opage}[1]{\marginnote{\footnotesize\textit{[#1]}}}

\begin{document}

\begin{titlepage}
\begin{center}
\vspace*{3em}
{\Huge Mathematics and Dialectic.}\\[1em]
\rule{4em}{0.4pt}\\[1.5em]
{\large \emph{A Philosophical Treatise}}\\[1.5em]
by\\[1em]
{\large R. Nielsen,}\\[0.3em]
Professor of Philosophy.\\[4em]
{\large Copenhagen.}\\[0.5em]
Published by the Gyldendal Bookshop (F. Hegel).\\[0.5em]
{\small \emph{Thiele's Printing House.}}\\[0.5em]
1859.\\[4em]
{\small Translated from the Danish}
\end{center}
\end{titlepage}
\thispagestyle{empty}
\vspace*{14em}
{\small
\noindent What is the life-principle in science ought never, out of a misunderstood regard
for popularity, to be sacrificed or even merely veiled; whoever lacks the ability and
the will to grasp philosophical subjects in their truth had better leave off occupying
himself with them altogether.\par
\hfill J. L. \emph{Heiberg.}\par}
% the motto, Danish: „Hvad der er Livsprincipet i Videnskaben, bør aldrig, af misforstaaet Hensyn paa Populariteten,
% opoffres eller endog blot tilhylles; hvo som ikke har Evne og Villie til at fatte de philosophiske Gjenstande i deres
% Sandhed, maa hellere reent lade være at beskjæftige sig med dem.“
\clearpage

% ===== PREFACE (Forord, pp. I-VII) =====
% --- p. I ---  (Forord, PDF 11; no folio printed)
\begin{center}{\large Preface.}\end{center}

\opage{I}To bring about a good understanding between mathematics and
philosophy is still not quite so easy, although, going by
Plato's hint, one might almost think that the matter went of itself.

The Platonic view, that mathematics must essentially
be regarded as a preparatory school for dialectic, is developed above all
in the doctrine of the state (books 6 and 7). The mathematical,
which lies on the boundary between the visible and the
invisible, is just the thing to promote abstraction. ``Those
who occupy themselves with geometry and arithmetic and other
mathematical sciences presuppose in every proof
certain things, such as the even and the odd, figures, three kinds of
angles, and other things akin to these. In these
investigations the soul is compelled to make use of presuppositions, not
in order to arrive at the first ground, for it cannot go
beyond these presuppositions, but it uses images borrowed
from earthly and sensible things, and chooses among these
those that must be regarded as having the greatest clarity and must be esteemed
most highly.'' ``If the talk is thus of a square and its diagonal,
they do not conduct the proof for the sake of the one they draw, but they
have the square itself in their thoughts''; in everything
else, too, that they form and draw they make use only of images,
``in that they strive to cognize that very thing itself, which is seen only with
the understanding.'' Number and unity belong, moreover, to that in
objects which particularly summons to thinking. ``For if
% --- p. II ---
\opage{II}unity in and for itself could be sufficiently grasped by
sight or by another sense, then it could hardly
lead us to what is ...; but when it is never seen except
in opposition to its contrary, so that nothing more shows
itself as one than as the opposite, then a judge is surely
necessary, who can decide the matter, and the soul must necessarily
become perplexed, and set all its thoughts in motion in order
through an investigation to find an answer to the question
what unity in and for itself actually is.'' --- ``The art of reckoning
and counting, which is necessary as much for the warrior as for the
philosopher, is therefore quite suited to be introduced by law
into the state, and we ought to persuade those who are to take part in the
highest affairs of state to occupy themselves with the art of reckoning,
not in the common way, but in such a way that through
pure thinking they arrive at a view of the nature of numbers,
and to learn it, not in order, like merchants and peddlers, to
use it in buying and selling, but partly to apply it in
war, and partly to make it easier for the soul to raise
itself from the perishable and changeable to the true and
what is.''

This conception of mathematics on Plato's part is certainly
somewhat one-sided, for mathematics is not for the sake of war
and not for the sake of philosophy; it is originally for its
own sake. But if mathematics, because it is and
must be a science independent of philosophy, should now also
wish to make itself independent of the universally scientific, and
thereby detach itself from the humane orientation of the spirit, the one-sidedness
would surely become still greater, and moreover more
alarming.

How one is to go about it when one really wants to
drive it to the point of becoming an out-and-out mathematician has been indicated, in a pointed characterization, by a French author
of the 18th century (\textit{Saverien: Histoire du
Calcul des Infiniment-Petit}).
The talk is of Maclaurin's objections to regarding the
infinitely small and the infinitely great as magnitudes, since
% --- p. III ---
\opage{III}the transition from the finite to the infinite is utterly
inconceivable. ``Meanwhile'' --- this is a view the author
imagines being introduced --- ``the consequences
one derives from those inconceivable presuppositions are after all real
truths. And what does it matter, then, that the principle set up
is inconceivable, if it only leads us to
discover a truth?'' This view is nevertheless rejected,
and for the following reasons. ``In the first place,'' the author replies,
``when the principle is unknown, we cannot know
whether all the consequences following from it can be admitted:
we go as if blindfolded, and it might happen that after
laborious exertion we finally ran up against something that precisely
revealed the nature of the principle, i.e. against something we could not
conceive. In the second place,'' he asks: ``is it not
a shameful degradation of reason to let it
follow a concatenation of terms, a series, whose beginning
and end it cannot see through? It is a great evil
that is born of this kind of routine: accustomed to operating without
seeing the ground, the mind loses all practice in thinking, and
ends by carrying out everything quite mechanically. \textit{Aussi voit-on
les jeunes Gens les plus incapables de refléxions devenir
facilement habiles Calculateurs ... Comme on croit que
cette aptitude est un fruit du génie on fait sonner bien
haut ce succès; et un enfant se trouve grand Homme,
sans qu’ il se soit encore avisé de penser. Les éloges
qu’on lui croit dûs et qu’on ne cesse de lui prodiguer, le
persuadent. En faut-il davantage pour faire éclore dans
lui l’orgueil, enfant chéri de l’ignorance, et pour le rendre
dans la suite un homme vain presque inutile à la Société
et toujours insupportable}.''

Certainly it is still not so rare that an ``\textit{enfant
chéri de l’ignorance},'' such an unschooled little ``\textit{grand
Homme},'' later works his way up to becoming something more
than a ``\textit{habile Calculateur},'' to becoming a mathematician, even
a capable mathematician and an effective member of society; but
% --- p. IV --- (pp. IV-V are missing in the NYPL copy scanned; read at the image of the second copy, Google Books X3haAAAAcAAJ, PDF 16-17)
\opage{IV}impressions from the boyhood years are lasting, and lack of a good
upbringing is not easily made good. If that \textit{orgueil} in Saverien's
dark picture, under favorable circumstances,
is duly hatched out into mathematical brashness, the capable
of society will nevertheless always have the intelligence somewhat coarse,
and then naturally, when occasion is given, they will charge at
philosophy in a great deal of wind.

In order, then, to withstand mathematical whirlwinds, when
it blows from the raw quarter and the wind rises, philosophy
must see to securing for itself some pillars that are not easily
blown over, i.e. some presuppositions to which no fair
and competent reader will deny his recognition.

The first presupposition is that the mathematician who
wants to raise objections against the way in which one who philosophizes
has conceived and applied mathematical formulas and
expressions should take the trouble to look so much at the train of thought
that he does not read into such expressions and formulas anything else or
more than what the one who philosophizes can be seen, expressly, from the context
in which they occur, to have wished to put into them.
How easy it is, when one takes propositions that, according to
the context, hold only with a definite restriction
as if they were said in pure generality, to make even
a hero of mathematics into an ignoramus in mathematics, is
accordingly not difficult to prove. It says in Plato's
Timaeus: ``But if the body of the universe had
had to become a plane without depth, then one mean would have
sufficed to connect the other two parts and itself
with them. But since it was to have the constitution of the solid,
and the solid is never connected by one but by two means,
God therefore set water and air between fire and earth, and
as far as possible brought them into such a relation to one another
that as fire is related to air, so
air is related to water, and as air to water, so
water to earth.'' Concerning this exposition it is then added in
C. J. Heise's note (to the translation, p. 28), that it
% --- p. V ---
\opage{V}is ``grounded on the following two mathematical propositions: 1) \emph{Two
planes or the areas of such can, as extreme terms
in a proportion, be connected by one plane or
one area as middle term, and 2) Two bodies or the
volumes of the same cannot, as extreme terms, be connected
by one body or one volume, but
require two different such as middle terms}.''
Since one can after all, as Proclus already remarked, always
set two or more geometrical means between two ``planes,'' and likewise insert
one geometrical mean between two ``bodies'': it will --- from the
present standpoint of ``criticism'' --- be easy, even for a mediocre
mathematician quite easy, to catch Plato in ``the
grossest errors,'' if the task is essentially to make wind.
If the task, on the other hand, is not to make wind, but to understand
Plato, then admittedly several restrictions must be made here,
since here, according to the context, it must be presupposed that Plato
did not think of the solution of the problems by construction,
but by calculation, since he presupposed that the sought-for
mean is to be rational, and the extreme terms in question
squares and cubes of prime numbers; but with these restrictions
the propositions are also quite correct.

The next presupposition is that the mathematician, if
in a philosophical treatise he here and there comes upon real
errors, inaccuracies and misunderstandings, nevertheless does not
straightway deny to the one who philosophizes, who could commit
such errors, all preparation and capacity for entering
into a mathematical train of thought. Certainly such
a declaration of incompetence always has its momentary effect when
the task is to play the schoolmaster. But when the task is not
to play the schoolmaster, but to find out the right relation between
insight and error, the sentence of condemnation here complained of
will surely suit a boastful pettifogger far better than
the one who is truly superior, who really sees what
matters. If one may from errors in the treatment of a
% --- p. VI ---
\opage{VI}mathematical formula or from the misunderstanding of a single
expression immediately and without further ado infer ignorance of the
mathematical: then many become ignorant here, then D'Alembert
and Euler are surely ignorant too, indeed, then Cartesius is a
bungler, and the ``strict criticism'' of the present day may brandish the corporal's stick
over Newton and Leibnitz.

A third presupposition is that the mathematician who
wants to undertake to judge the philosophical in mathematics
must first of all take the trouble to acquaint himself
to some extent with the whole philosophical train of thought from which
the treatment in particular proceeds, and moreover know how to keep
the different views of mathematics and philosophy
so apart from one another that they can properly be
set side by side and clarified by one another. Certainly it is
important that inaccuracies be pointed out and errors committed
be corrected; but there is one thing that in the scientific sense is
still more important, and that is that a distinction be expressly drawn between
essential and inessential errors, a separation that in
the investigation concerning the relation between philosophy and
mathematics is all the more unavoidable, as no one of scientific
education will be able to deny that much, when
seen from the side of philosophy, and moreover in a sketch, an
outlined presentation, is more hinted at than developed, and can be
quite inessential, whereas in a mathematical and moreover
detailed textbook it would be even very essential. By
either not heeding this difference at all, or better: by
expressly admitting it in words, in order afterwards in the judgment
itself to deny it in deed, it is an easy
matter to magnify inessential errors, a slight art to make
philosophy's gnats into mathematical elephants and impress
the reader with many elephants.

That such presuppositions, which are unavoidable only with regard to
the scientific, naturally must not be
troublesome to such as otherwise, from their own standpoint,
might feel prompted to favor the present
% --- p. VII ---
\opage{VII}treatise with ``a strict criticism,'' goes without saying. If the fence is given up and the field-gate open,
then that author would truly have to be very naive who would
think that some propositions set up in a preface could
be for him either a protection against the loose or a hedge against the
bellowing, against those on the common pasture of literature who like
to butt; no, the propositions set up are by no means to prescribe
duties for the horned cattle of mathematics, they are only to
be binding for the author himself and for those who with
him want something scientific.

Should an assessment of the present treatise,
keeping to the presuppositions set up, either
really or supposedly be able to point out serious errors
and slips: the author hereby declares himself bound to
answer to the named reviewer, whether he belongs, moreover,
among the younger or among the older, among authorities or
non-authorities, in every way, so that the matter
may in the end be cleared up; whereas --- and this is
hereby likewise expressly laid down --- every attempt at
criticism which, by breaking the said propositions,
confounds the essential with the inessential, and confuses the points of view,
is to be rejected as unwarranted, whether the attempt come for
the rest from a small or from a great authority, and offer whatever
it will, and be as it will, with coarse sand or with
fine sand, with slight standing, or with bluster and bravado.

\begin{flushleft}\hspace*{2em}Copenhagen in December 1859.\end{flushleft}
\begin{flushright}R. Nielsen.\end{flushright}

\clearpage

% --- p. 3 ---
\setcounter{footnote}{0}%
\opage{3}Mathematics and philosophy are sciences of different kinds;
for that reason each must assert its own independence. Mathematics
is an exact science, philosophy a dialectical one.
As an exact science mathematics has no use for
philosophy; as a dialectical science, philosophy, by contrast, has
use for mathematics; the exact science relates
to itself by excluding other sciences;
the dialectical science relates to itself through other sciences.\footnote[1]{Cf. Philos. Propædeut. i Grundtræk 1857, p. 99 § 22.}

From the standpoint of the theory of cognition, philosophy's
relation to itself through the mathematical is therefore in general
to be determined thus: mathematics fixes its
concepts in the exact form; the concepts so fixed
cannot without further ado be transferred into philosophy; philosophical
dialectic, on the other hand, can, by bringing the exact
determinations of relation to consciousness, gain in clarity and
sharpness, and so attain the precision that the
cognition of things necessarily must presuppose. In order to show
how the theory of cognition articulates the dialectical development of concepts
by means of moments of exact analysis, we let
the exposition in the present treatise turn upon three
main points: I. \emph{The concept of function}, II. \emph{The problem of infinity}
and III. The analysis of the infinite.

% --- p. 4 ---
\setcounter{footnote}{0}%
\opage{4}
\begin{center}I.\end{center}
\begin{center}The Concept of Function.\end{center}

That the concept ``function'' does not belong exclusively
to mathematics is well known.\footnote[1]{Cf. Phil. og Mathem. pp. 5--6.} Just as intuitions, according to Kant,
rest on ``affections'', concepts rest on ``functions'':
``Ich verstehe aber unter Function dieEinheit derHandlung verschiedene
Vorstellungen unter einer gemeinschaftlichen zu ordnen.''
(Kritik der r. V.) The unity of action by which Kant
characterizes the activity of the understanding from the general-logical
standpoint is, as is easily seen, essentially akin
to the unity of action by which the understanding apprehends
the mathematical relation between dependent and independent
magnitudes; only the difference must not be overlooked. Outside mathematics,
as far as the category of function is concerned, the emphasis
falls on the law-determined action, in mathematics, by contrast, on
the relation of dependence by which the action (the operation)
is determined; outside mathematics the category is apprehended chiefly
from the active side, in mathematics, if one likes, from
the passive or at any rate neutral side. ``A magnitude that
depends on one or more arbitrary magnitudes, with which
it can therefore be thought of as connected by an equation, is called a
function of this or these. Thus e.g.\ $ax^{3}$ is a
function of $x$, $a$ being thought of as unchanged, the logarithm a
function of the number, the sine (as a line) a function of the radius
and the arc, etc.''\footnote[2]{Chr. Jürgensen. Begyndelsesgrundene af den mathematiske Analyse. Kbhvn 1855. p. 18.}

The special significance that the category has thus acquired
in mathematics shall now contribute to a philosophical
development of that same category, as we investigate the concept's
requisites, its relation to law, and its peculiar
presentation.
% --- p. 5 ---
\setcounter{footnote}{0}%
\opage{5}
\begin{center}a) The original requisites of the function.\end{center}

For the proper determination of the mathematical
function there is required:

\begin{center}α) a \emph{constant} and a \emph{variable} magnitude.\end{center}

In the expression $f(x)=ax$, $a$ is constant, $x$ variable; in
the expression $f(a)=ax$, by contrast, $x$ is constant and $a$ variable. If
now every magnitude can, according to its concept, be now constant,
now variable, as circumstances may be, then it is already evident from this
how easily the mathematical axiom ``every magnitude is equal to itself''
can be attacked by dialectic. That every finite magnitude
must, insofar as it is constant, be equal to itself: $a=a$, is beyond contradiction; but how does it stand with
the variable magnitude? The variable magnitude must surely,
according to its concept, be able incessantly to depart from any
given value whatever, and just thereby, in its equality with itself,
become unequal to itself:
\[ x \overset{>}{\underset{<}{=}} x \]
. Were one, in order to maintain the exact validity of the
principle laid down, to take the variable
magnitude, $x$, as a provisional designation, a collective
expression, which says that $x$, by assuming any values whatever,
resolves itself into a plurality of constant magnitudes: $a=a$, $b=b$, $c=c$ etc.; then one would thereby not only have made the
variable magnitude one with the arbitrary one, which in other respects is mathematically inadmissible, but even
altogether abolished the concept of the variable magnitude. That the determination
\emph{variable} belongs to magnitude just as essentially
as the determination \emph{constant}, is seen from the fact that magnitude
is defined precisely as that which can be increased and diminished.\footnote[1]{Everything that can be thought of as increased or diminished, and that can be resolved into homogeneous parts, is called magnitude (quantum), and is said to possess greatness (quantitas). Ramus: Elementær Geometrie. 1850. Indledn.}
% --- p. 6 ---
\setcounter{footnote}{0}%
\opage{6}Were one, trusting in the supreme principle of formal logic,
to explain $x=x$ thus: the variable must, insofar as
it is not constant, necessarily be equal to itself in this, that
it is variable; then one might with the same right as when one
says: the variable is variable, therefore: $x=x$, also
say: the infinite is infinite, therefore: $\infty=\infty$; but
mathematics must protest against such an explanation. What
does $\infty$ signify? If $\infty$ does not designate a magnitude, then the
equation $\infty=\infty$ is meaningless; in general logical terms it would then have to
read: the infinite is what it is, a negation of magnitude.
If, however, the infinite is to be taken as
a magnitude, then the equation $\infty=\infty$ is incomplete, i.e.\
as invalid as it is valid:
\[ \infty \overset{>}{\underset{<}{=}} \infty \]
is the complete
expression. To lay emphasis on the difference between the
determinate and the as yet indeterminate magnitude does not help.
If the finite magnitude is, according to its concept, determined by opposition
to the infinite, then the infinite is surely at the same time
determined by opposition to the finite; if finite
magnitudes are to be determined by their relation to one another,
then the corresponding must surely hold of the
infinite ones.

The axiom laid down: ``every magnitude is equal to itself'',
has thus now resolved itself into two propositions, of which
the first is a tautology, the second a contradiction. For it
is not magnitude as magnitude, but magnitude
as constant magnitude, that is equal to itself; the emphasis falls
not on being magnitude, but on being a
constant magnitude. And what then is a constant magnitude?
A magnitude that is equal to itself. Thus the axiom says:
a magnitude that is equal to itself is equal to itself; a pure
tautology! If, next, magnitude is taken as that ``which
can be thought of as increased or diminished'', then the emphasis falls
not on being magnitude, but on being variable
% --- p. 7 ---
\setcounter{footnote}{0}%
\opage{7}magnitude; the variable must show itself in varying; the variable
magnitude, which according to the axiom was likewise supposed to be equal to itself,
can therefore be equal to itself only in this, that it constantly
becomes unequal to itself, which is a contradiction.

These difficulties are removed when we look to the function.
The concept of function is, seen dialectically, the true concept of magnitude,
since the two-sided determination of magnitude, to
be constant, i.e.\ equal to itself, and yet variable, i.e.\ unequal to
itself, is here posited in unity. Even if one cannot immediately
say: $x=x$, insofar as $x$ is variable, one can nevertheless
say with truth: $f(x)=ax=f(x)$, since the same
function, as function, must under all changes be
equal to itself.\footnote[1]{Thus also $x=x$ after all, insofar as $x$ is a function of $x$.} In that the tautology is removed, the contradiction
is at once abolished. The unequal no longer
stands as abstract negation over against the equal, nor the variable
over against the constant; the function precisely lets the equal
reflect itself in the unequal, the variable in the constant, the moment of difference
in identity.

To the necessary requisites of the concept of function there belongs further:

\begin{center}β) \emph{Two variable magnitudes, an independent and a dependent one.}\footnote[2]{That a function such as $y=f(x, z, t, u....)$ contains several independent variables is left out of consideration in this context, where the matter concerns merely a development of the concept.}\end{center}

In the equation $f(x)=y=ax$ the magnitudes $x$ and $y$ are changeable, but in such a way that $x$ is independent, $y$ dependent.
The change in $x$ then relates to the change in $y$ as
presupposition to consequent, as ground to consequence. Since the
concept of function shows us precisely the reflected unity of presupposition and consequent,
of the independent and the dependent, the difference between the mere presupposition and
the proper ground can in this context be clarified. The
% --- p. 8 ---
\setcounter{footnote}{0}%
\opage{8}presupposition in its immediacy is a given, an assumed thing,
which is apprehended in representation, and falls under
reflection only insofar as it is seen that, in order to be presupposition,
it must have its consequences. One may thus assume
a presupposition, hold it fast in representation as
something familiar, and yet be greatly surprised by its
consequences. The surprise stems from this, that the essence of the presupposition
was not seen through from the beginning;
the presupposition was indeed made distinct in a representation, and
to that extent apprehended from a certain side, but not clarified on all sides,
not apprehended according to its concept. Only when we see through the
presupposition, and discover the consequences hidden in it,
do we cognize the essence of the presupposition, posit it as
ground, and the consequence as the grounded result. What
difference it makes whether one has the ground merely as an immediate
presupposition, or sees the ground in the essence of the presupposition,
can with regard to $f(x)=ax=y$ be clarified
by a very simple example. A asks B: What will you
take for your game bag? ``3 Rdlr.!'' is the answer. ``3 Rdlr.,''
continues A, ``I will give you at once, and if you will wait
eight days, until a remittance that I am expecting from
home arrives, I will even give you $3 x$ Rdlr.''
``What do you mean by this $x$?'' ``$x$ is,'' answers A, ``a
changeable magnitude, a varying unit, which can assume
different values, e.g.\ $1, 2, 3, 4 ...$; the $3 x$ Rdlr. would thus, as you see, in this way
become $... 3. 1 Rdlr., 3. 2 Rdlr., 3. 3 Rdlr.$ etc.

B reflects: ``$x$ is always a unit; the $3. 1 Rdlr.$ I am
certainly sure of; perhaps what is promised may, when the remittance
arrives, and A is in a good mood, rise to $3. 3$, indeed
perhaps even to $3. 4 Rdlr.$''; so: B relies on the presupposition
lying in the varying unit, accepts
the offer, and waits. The eight days pass; the remittance
fails to arrive, and A now explains to B the meaning of the
variable magnitude: ``$x$,'' he says, ``can, as said, be:
% --- p. 9 ---
\setcounter{footnote}{0}%
\opage{9}$1, 2, 3, 4 \ldots$, but, since it is a varying unit, it can
surely also --- it lies in our presupposition --- be
$\frac{1}{2}, \frac{1}{3}, \frac{1}{4} \ldots$ We can now this time choose to set $x=\frac{1}{288}$; you get, then, according to our express agreement, i.e.\
according to our presupposition, $3x$ Rdlr. $= \frac{1}{288}$ Rdlr., say and % print: og (misprint)
write one whole shilling for your game bag.'' Now a
light dawns on B, a light that is surely worth its 2 Rdlr. 95 ß;
for B makes the acquaintance of the varying unit,
he sees through the presupposition, and discovers the ground.\footnote[1]{It will be understood that what in this example, seen from a purely mathematical standpoint, would be too subjective, acquires from the standpoint of philosophy a significance in that the illusion hidden in the presupposition is psychologically revealed.}

To engage with presuppositions without knowing the fundamental relations
is often as precarious as playing with
loaded pistols; before one well looks around, a presupposition
may so to speak go off, and fire off a consequence that,
even if it does not kill, still wounds. The concept of function
accordingly also demands that

\begin{center}γ)

\emph{The relation between the function's dependent and its independent is determined as a relation of ground.}\end{center}

Notwithstanding that the relation of the consequence to the presupposition always
is a relation of necessity, it is nevertheless recognized as a
relation of ground only when the presupposition shows itself as ground.
The relation of necessity does indeed express that every presupposition
must have its corresponding consequence, and every consequence
its corresponding presupposition; but so long as the relation
remains in indeterminate generality, it is still
undecided what the expression ``corresponding'' really means. In order
that a corresponding consequence shall proceed from the corresponding
presupposition, the condition must come in; only through the condition
do the presupposition and the consequence become more closely determined.
% --- p. 10 ---
\setcounter{footnote}{0}%
\opage{10}But the condition, which, seen from one side, is different
from the presupposition, is, seen from the other side, identical
with the presupposition. In its difference from the conditions
the presupposition is the essential, and according to its essence ground;
but only in its unity with the conditions is the ground sufficient,
and as sufficient ground (\textit{ratio sufficiens}) it is
the comprehensive presupposition: through the presupposition's identity
with the condition the formal relation of necessity becomes a
relation of ground.

So long as the equation $f(x)=y$ merely states that every
change in $y$ is dependent on a corresponding change in $x$, the relation of dependence is still indeterminate and formal,
sounding immediately of presupposition and consequence. It
is, however, so far from the case that the formal relation of dependence
should in and for itself possess any true general validity,
that one is not even able to point out its significance
without introducing conditions;\footnote[1]{Already by passing from $\Phi(x, y)=0$ to $y=f(x)$ we have set a condition.} but in that conditions are introduced,
the general relation is transformed into a particular one. ``One
cannot'', it says in mathematics, ``say that $2$ is a
function of $3$, because $2=3^{2}-7$; for $2$ is also $= 3-1$, or $= 3^{3}-25$ etc., so that $2$ may depend on $3$ in
manifold ways. But one can well say that
$y$ is a function of $x$ by means of the equation $y=x^{2}-7$, and
that therein $y=2$ corresponds to $x=3$''.\footnote[2]{Chr. Jürgensen. Begyndelsesgrund. af den mathem. Anal. p. 19.} Thus: the condition for
$y$ to show itself dependent on $x$ is that different values are assigned to $x$; the condition for the stated dependence
to designate a function is that the manner of dependence
be more closely fixed. But the condition for the closer
determination of the manner of dependence is in turn, when we
look more closely, this, that the general function with inner
% --- p. 11 ---
\setcounter{footnote}{0}%
\opage{11}necessity must divide itself into particular functions: $f(x)=y$ acquires significance only by being set as: $f_{1}(x)=ax$, $f_{2}(x)=x^{a}$,
$f_{3}(x)=a^{x}$, etc. Only in that reflection is able
to develop the general function into particular functions,
and the particular functions in turn into identity with the
general one, can the relation of dependence be recognized as a relation of ground.

If the particular functions: $f_{1}(x)=ax$, $f_{2}(x)=x^{a}$,
$f_{3}(x)=a^{x}$ etc., are set against the general function: $f(x)=y$, then it at once appears that the general function can neither immediately
keep itself outside, nor coincide with, the
particular ones. If the general function acquired a validity of its own over against
the particular functions, it would thereby precisely become
particular, and thus not general. If the general function
coincided immediately with one of the particular ones, it would thereby
have the rest outside itself, and would not be the general one either;
indeed, even if it were assumed to coincide with each single particular one in such a way
that none fell outside its domain,
it would precisely be parcelled out into particularity, could not even
be a function in itself, let alone a general function.

That the general function cannot be held fast in immediate being does not at all mean that the general should
be without true validity, an empty schema, a subjective
abstraction; it proves on the contrary that the general has in itself
an essential and for that very reason reflected validity: in its
reflected validity the general function is determined as
the ground of all particular functions, as the fundamental function.

The manner in which the particular functions, in that they
each according to its peculiarity render the fundamental function,
reflect themselves in one another, shows everywhere that the presuppositions
with their conditions are at once moments of reflection
in the essence of the ground and yet determinations of immediacy
in the posited multiplicity, the grounded
% --- p. 12 ---
\setcounter{footnote}{0}%
\opage{12}existence. For the understanding of the many particular functions
it would therefore, according to the relation of ground, be sufficient
to understand a single one, if the understanding of this one were
exhaustive; but for an exhaustive understanding of one
function it would, likewise according to the relation of ground, be
necessary to understand them all.

\begin{center}b) The Function and \emph{Law}.\end{center}

To grounding there belongs therefore an element that goes
together with the ground, and yet falls outside the ground;
this element is the condition. The condition, which in reflection
is identical with the ground, insofar as the particular is grounded,
falls in its immediacy outside the ground, and gives in
this respect a shade of the \emph{contingent}. If the
contingent is in this way a something necessary in the ground,
the necessary must surely at the same time acquire a tinge of
the contingent, and show itself as conditioned necessity. That
which is \emph{contingent in cases} is the changing singular, the
multiple; that which is necessary in cases is the essentially
One, the general, \emph{law}; between these extremes falls
the intermediate determination of the superficially general, the common,
\emph{the rule}. The way from the multiply given through the conditions
to the unity of ground is a way from the \emph{case} through the rule to law.

\begin{center}α) \emph{The Case}.\end{center}

In the world of experience the case is found as a fact,
and must be taken as it is given; in the a priori
science the knower himself posits the case; but in both
cases the contingent is like itself in this, that, once
it is set, it is set, whether it has shown itself
as fact, or as something assumed \textit{a priori}. It
might indeed at first glance seem as though the case given as fact were otherwise contingent than that which consciousness
itself posits. As objective case the fact is
% --- p. 13 ---
\setcounter{footnote}{0}%
\opage{13}surely, because it is, and as it is, independent of
subjectivity; what is subjectively posited, by contrast, is because consciousness
wills that it shall be, and is precisely so because
it is posited so: the arbitrary thus seems to devour
the contingent. This, however, is only an illusion; that arbitrariness
cannot annihilate contingency is easy to
see. If the subjectivity that posits a case wants to posit
it solely by virtue of arbitrariness, what then is
abstract arbitrariness other than a conscious contingency?
I arbitrarily posit an $a$, and am conscious that I
might just as well have posited a $b$; I see that the arbitrary
is in itself a contingent. The more subjectivity
surrenders itself to arbitrary fancies and representations, the
more it also delivers itself over to contingency. If
the arbitrarily posited is to acquire intellectual significance, it must
be posited as a presupposition whose consequence is acknowledged:
as many special presuppositions, as many consequences;
the peculiarity of the cases must for the sake of the consequence be
respected. To respect the peculiarity of the cases is
to attribute objective significance to them; but, in that the case
acquires objective significance, the contingent precisely shows itself.

Not only in the empirical external world, but also in the
a priori sphere of representation can the multitude of cases
therefore prove overwhelming, and the unforeseen surprise us. In
mathematics the number of numbers is innumerable, and the possibility
of numerical combinations inexhaustible. In that the individual numerical magnitudes
are held fast just as they are immediately had in representation,
every numerical magnitude is a phenomenon, and one grasps involuntarily
that the numerical phenomena in their changing multiplicity
do not fall behind the phenomena of experience. Let the number $9$ be
e.g.\ the given phenomenon; let one now begin to investigate
in what way this magnitude comes about,
and lets itself be resolved by other magnitudes; and one will
soon convince oneself that it is not this or that
% --- p. 14 ---
\setcounter{footnote}{0}%
\opage{14}special case, but millions upon millions of possible cases
that one has to do with.

How then shall we treat the innumerable cases?
This question is in the a priori essentially to be answered
in the same way as in the empirical. As regards the empirical
cases, the point is, by comparative
abstraction, to bring out the common, separate the dissimilar,
and combine the similar; to the a priori cases, insofar
as the contingent is brought out, a similar procedure
is to be applied.

\begin{center}β) \emph{The Rule}.\end{center}

When what is alike in many cases is abstracted and held fast,
the rule emerges. In the four arithmetical operations, as
these are set forth and taught popularly for practical
use, a multitude of rules is seen. One begins with prescriptions
for the treatment of the special problems; it is shown
in individual examples what one has to do, how one
goes about setting up the problem and arriving at the result.
By applying the prescribed procedure to a
series of cases, and seeing how each time it works out
that the correct result comes out, one convinces oneself of the validity of the rule. The rule, which is abstracted from the individual
cases, has in this context interest only insofar as it can be applied to the individual cases: an
investigation of the rule in its generality as rule and of
the mutual relation of rules presupposes a higher standpoint.

So long as consciousness is lost in cases and rules,
thought is ensnared in the finite; scientific insight
is lacking. Acquaintance with the rules and skill in
applying them is not yet science: the master reckoner is
not yet a mathematician. The prejudice that many seem to
harbor against mathematics, that it, without essential relation
to the Idea and the ideal, should be only a dry object
for ``the dry understanding'', stems surely from this, that those who
judge thus know the mathematical only in its
% --- p. 15 ---
\setcounter{footnote}{0}%
\opage{15}finite, most superficial form as a sum total of sums
and rules; for in this poets, women, and
the brilliant are undeniably right, that sums and rules exist for
``the dry understanding.''

By bringing out the rule, the understanding seeks to attain the
generally valid, but the rule as rule is not generally valid;
every rule has, or must at least be able to have, its
exception: \textit{nulla regula sine exceptione.}\footnote[1]{When, therefore, in a freer use of language, one speaks of generally valid rules, one must, insofar as the generally valid is brought out, think of regulative laws. In this sense expressions such as ``General rules of differentiation'' are then taken. Jürgensen. Begyndelsesgrunde af den mathem. Analyse. p. 43.} Should anyone
ask whether one could not then abstract from cases such as this: $2(3+4)=2.3+2.4$ the rule valid in itself: $a(b+c)=ab+ac$? then we should have to answer:
No! What is merely abstracted is not yet cognized as the
generally valid; the in itself generally valid is not a rule.
How precarious it is to want to secure the generally valid by abstraction from individual
cases can, as regards mathematics, be seen among other things
in the doctrine of prime numbers. One
finds: $2^{3}-1=7$, $2^{5}-1=31$, $2^{7}-1=127$, nothing but
prime numbers; but if one were now to draw from these three cases
the conclusion that $2^{2m-1}-1$ must necessarily be a
prime number, one would miscalculate, since already the fourth case: $2^{9}-1=511=7.73$, gives no prime number.
Fermat assumed that $2^{(2^{n})}+1$ was a prime number; but Euler
showed that $2^{(2^{5})}+1=2^{32}+1$ is divisible by $641$, and consequently no prime number. That in such examples the exception to
the rule occurs now with the fourth, now with the fifth
case etc., is in the whole consideration itself like a case.
Had the exception first occurred with the 99th case,
we would yet in the extended rule still not have
found the in itself generally valid, but would on the contrary have been exposed
to being deceived by the multitude of cases.
% --- p. 16 ---
\setcounter{footnote}{0}%
\opage{16}
\begin{center}γ) \emph{The Law}.\end{center}

That which is universally valid, from which there is no exception and can be none, is not a rule but a law. If the equation
$a(b+c)=ab+ac$ is to be conceived as a law, its validity must not rest on an abstraction from the individual cases
but on insight into the essence of the matter. A single case that
leads to insight is of greater worth in this respect than
a thousand cases that merely supply material for a comparative
abstraction. That $\frac{1}{1-x}$ will, by continued division, yield a
series of countless terms according to ascending powers of $x$, namely:
\[ 1+x+x^{2}+.... \]
is not proved by continuing the division
to infinity, but by making clear what the division,
on the given presupposition, necessarily carries with it.
In its immediacy the case stands before reflection as
that on which reflection is exercised; the result of finite
reflection, the general, the relatively universal, is the rule abstracted from
the cases and applicable to the cases;
in reflection's unity with the essence of the matter the law is cognized.
In mathematics, where intelligible magnitudes take the place of
sensuous phenomena, the law can be presented in the pure
exact form. What, then, is the law? According to its concept
the law expresses the necessity with which the consequence enters,
once the whole presupposition, i.e.\ the ground with its conditions,
is given. If $a=b$ and $b=c$, then $a=c$; the
necessity that the law expresses is thus hypothetical: as
the expression of hypothetical necessity the law is the universally valid.
Whether we say: ``it is a law that the magnitudes $a$ and $c$ must be equal to one another, when each of them taken by itself
is equal to the magnitude $b$''; or we say: ``it is necessary
that it must be so''; or: ``this proposition is universally valid'', it means One and the Same.

As the universally valid, the law finds its expression in the \emph{axiom}, in the \emph{definition}, in the \emph{theorem}. In each of these
% --- p. 17 ---
\setcounter{footnote}{0}%
\opage{17}modes of expression we recognize the hypothetically necessary, which is
designated by the relation between presupposition and consequence:
if $a=b$, and $b=c$, then $a=c$; if all
points in a plane curve have the same distance from a
determinate point in the plane, then the curve is a circle; if
the figure $ABC$ is a triangle, then the sum of the angles is $=$ $2R$. Whether the form of expression itself, incidentally, is hypothetical or
thetic makes no difference in this respect; the essential
difference rests on the grounding.

The \emph{axiom} expresses the law in an immediate way.
``Only one line can be the shortest path between two points.''
Were one to paraphrase this fundamental proposition thus:
``It is a law that between two points there can be only
one line that is the shortest,'' such a paraphrase
would admittedly be superfluous and forced. But that the
paraphrase is omitted must nevertheless not be misunderstood, as if
the law and the fundamental proposition fell apart, so that
the fundamental proposition were one thing and the law another; for the truth
is that the law is not specially mentioned as law, precisely because it
coincides immediately with the fundamental proposition: the legal proposition
that is in itself so evident that it does not require
grounding is a fundamental proposition.

The \emph{definition}, which states the presuppositions from
which the consequences of the law proceed under given conditions,
stands in a relation of reflection to the law. The definition does not coincide immediately
with the law, as the fundamental proposition does; on the contrary, it seems even at first glance to
be indifferent to the law; but when the relation between
the law and the fundamental proposition is to be determined more closely, it
turns out that the law rests in the fundamental proposition only insofar as
the latter rests on the definition. ``Why can only one line
be the shortest path between two points?'' It lies in the
nature of the matter, because it lies in the concept ``the shortest'',
but that it lies in the concept means, in other words,
that it lies originally in the definition. If $a^{x}=y$,
% --- p. 18 ---
\setcounter{footnote}{0}%
\opage{18}then the law holds: $x=\log y$, with $a$ as base. On what
does the validity of the law rest, then? On the definition of the logarithm.

In the \emph{theorem} the relation stated by the definition is
mediated in a peculiar way, and the law expressly
posited as law. ``A proposition or theorem (\textit{theorema})
presents a truth that is not immediately evident.
It contains a presupposition (\textit{hypothesis}), consisting of
one or more conditions, and a conclusion (\textit{thesis}). The
series of inferences by which the conclusion is led back
to the presupposition, so that it is thereby grounded with necessity,
form the proof of the proposition (\textit{demonstratio}).''\footnote[1]{C. Ramus. Elementær Geometrie 1850. p.\ 7—8.}

To grasp the essence of the law and its relation to the function
it is necessary to see how the proof and the
definition, the proposition and the fundamental proposition reflect
themselves in one another. Suppose, e.g., $Log$ as logarithm corresponding
to the base $A$, and $log$ corresponding to the base $a$; then
the law according to which the transformation from one of
these systems to the other takes place will show itself both to lie
in the concept of the logarithm and yet to require a proof. From
the proof: $A^{x}=y$, and $x=\mathrm{Log}\,y$; $a^{z}=A$, $z=\log A$, $(a^{z})^{x}=y$,
and $xz=\log y$; hence: $\mathrm{Log}\,y \log A=\log y$, and $\mathrm{Log}\,y=\log y \frac{1}{\log A}$, it is seen that the conduct of the proof is a development of the concept of the logarithm,
and that this development of the concept is a development of the law.
Since the conduct of the proof essentially coincides with
the development of the concept, the dialectical unity of the fundamental proposition and the proposition is seen.
It is the task of science to transform
every fundamental proposition into a proposition; that which is evident in itself
is evident only to a reflection: the immediate
must be mediated, the ungrounded must be grounded. That
exact science itself acknowledges this demand is seen, e.g.,
from its attempt to give a proof of the proposition that
% --- p. 19 ---
\setcounter{footnote}{0}%
\opage{19}``the shortest line between two points is a straight one''.\footnote[1]{See e.g.\ L. Oppermann. Elementær Plangeom. Kbh. 1845. p.\ 43.} That,
conversely, the proposition must be raised to identity with the fundamental proposition
means, in other words, that what is mediate in cognition
must be raised to a higher immediacy. The proof
has precisely the significance that, by making the individual
points of the matter immediately evident, it shows us the whole in a light
of reflected immediacy.

That function and law coincide in one development of the concept
is therefore indisputable in the theory of quantity; in order to
determine what a function is, we must presuppose the law,
in order to determine the law, we must in turn presuppose the function:
as many functions, so many laws. Just as the law, by
having the cases in its power, is raised above the mere
rule, so too the function is raised above the mere
rule; just as the law cannot dispense with the case, because in its hypothetical necessity
it cannot dispense with the immediate
presupposition, so the function in its conditioned
universal validity cannot dispense with the immediacy of the presupposition
and of the condition;\footnote[2]{``Indem die Function über den Begriff des realen continuirlichen Seins im Raume hinausgeht und die einfache Unmittelbarkeit von dessen Begrenzung negirt, so geht sie doch wieder vollständig mit ihm zusammen, weil sie die allgemeine Regel ist, vermöge deren auf die gesammte Folge von Punkten dieser Curve (here reference is made to the way in which a plane curve is determined, in that every point ($x, y$) can be determined by means of the equation $=$ $f(x)$) zurückgekommen werden kann, weil sie in der unveränderlichen Folge ihrer Rechnungsoperationen das allgemeine Gesetz, das Bleibende darstellt, wodurch die aus ihrem Flusse herausgerissenen isolirten vielen Punkte wieder zu einer Einheit zusammengeschlossen werden, und weil sie mithin dieselbe befassende Einheit ist, als welche die ebene Curve in der Anschauung vorliegt.'' Versuch einer Philosophie der Mathematik verbunden mit einer Kritik der Aufstellungen Hegel's über den Zweck und die Natur der höheren Analysis von Hermann Schwarz. Halle 1853, p.\ 26.} just as the one fundamental
% --- p. 20 ---
\setcounter{footnote}{0}%
\opage{20}function is, by virtue of the diversity of the cases and yet with
inner necessity, divided into a system of distinct
functions, so the one fundamental law is, on the same
terms, differentiated into a system of distinct laws.% NOTE: no mark printed in the print; the note stands at the foot of the page and fits best here (its last words: "et System af særskilte Love")
\footnote[1]{How function and law coincide in the fundamental concept is seen from countless examples. That $F(x, y, z) = 0$ is the equation of a surface means, in other words, that this function expresses the law for the changes of coordinates, corresponding to a given system, by which the points of the surface in question are determined. The number of points is infinite, an infinite number of points cannot be determined one by one; in the choice of the points that one will determine singly there is thus something arbitrary and contingent, but in virtue of the function the determination of the countless points is precisely reflected in the way in which the single point is determined: thus everywhere the finite is subsumed under the infinite, the changeable under the constant, the phenomena under the law, the contingent under the necessary, the thing under the concept, the distinct under the universal.\par
How many particular determinations of essence a fundamental function, and with it a fundamental law, can contain in the concept, without there yet being any question of the distinctness that individualizes itself for an immediate intuition, is likewise evident from the example cited. The general form $F(x, y, z) = 0$ states not only the essence of the surface as surface, but from this form even a plurality of properties, namely the fundamental properties of the surface: its curvature at a certain point, its position relative to other surfaces at a certain point, the equation of its tangent plane at any point whatever, etc., can be derived. Only as the general equation of the surface is articulated into the equation of distinct surfaces does the content of the fundamental law appear in a system of distinct laws.}

\begin{center}c) The Peculiar Presentation of the Concept of Function.\end{center}

For stating the requisites of the concept of function the
mathematical mode of notation has already been used; but what is peculiar
in this mode of notation can become the object
of a special reflection only when attention fixes
on the significance of the pasigraphic language of mathematics,
% --- p. 21 ---
\setcounter{footnote}{0}%
\opage{21}on the thought corresponding to the language, and on the insight resulting from the union of both, % print: Tænkuing (misprint)
the punctual comprehension.

\begin{center}α) \emph{Language.}\end{center}

A system of signs by which now feelings, now representations
and thoughts are expressed is a language. Among the
many kinds of differences of language, the difference between
the universally logical spoken language and the pasigraphic sign language is without
doubt the most thoroughgoing. ``For a philosophical presentation
of the concept the word is used, and the word sounds different
in the different national languages; but mathematics is not content
with the word; independently of the diversity of the national languages
the exact science has its own language: $f(x + y) = f(x). f(y)$ is neither Danish nor French, but mathematical.''\footnote[1]{Cf.\ Phil. og Mathem. p.\ 9—15.}

The mathematical language, in whose development men
such as Vieta, Cartesius and Leibnitz have worked with so much
success, is for the understanding what the language of musical notation is for feeling;
except that, as a language of the understanding, it stands in a more essential
relation to logic. % print: Logikeu (misprint)

The pasigraphic language of mathematics is in a certain respect
overvalued by Leibnitz, but undervalued by Hegel.

Leibnitz\footnote[2]{Cf.\ Kuno Fischer. Geschichte der neuern Philosophie, zweiter Band. Mannheim 1855. p.\ 10—16.} who, with so immense a learning that ``he
alone by himself represented an academy'', intervened in all the scientific
affairs of his age, was the founder of the academies
in Berlin, Dresden, Vienna and Petersburg, conceived
even in Rome the bold plan of introducing natural-scientific
studies into the Italian monasteries, and everywhere aimed
at establishing scientific states, republics of learning on an
extraordinary scale. For a universal monarchy of the sciences
the diversity of the national languages was, however, a
palpable hindrance; the thought then lay near, whether it were not
possible, by the introduction of a direct mode of notation,
% --- p. 22 ---
\setcounter{footnote}{0}%
\opage{22}an immediate \textit{signatura rerum}, to raise itself above this
linguistic particularism. Conceived in general terms and viewed
from a distance, such an idea had to have something exceedingly
attractive for a genius like Leibnitz. It was, after all,
precisely such thought-hieroglyphs that the ancients, in particular the Pythagoreans,
sought in the symbolic-mystical numbers, and of which the Kabbalists
had dreamed so much; such concept-characters,
such a signative language, the most exact sciences had indeed possessed,
in particular mathematics in its arithmetical and algebraic
signs; why should it now not be possible to extend
and express the world-wisdom that constitutes the whole
content of the sciences in a world-script, a pasigraphic
language? ``Der Mathematiker Leibnitz hat, scheint es, den
Philosophen Leibnitz zu dem Probleme der Pasigraphie überredet,
und der nach Universalität strebende Geist des letzteren
begünstigte gern ein Projekt, welches der Weltaufklärung so
viele Schwierigkeiten ersparen und die Idee auf dem kurzesten
Wege befördern wollte.'' Kuno Fischer.

Hegel has gone to the opposite extreme in his Logic.
In the doctrine of ``the Concept'' he censures Euler and Lambert, who
sought to reduce determinations of concepts to designations by
lines, figures and the like. ``Such objects,'' he says,
``have this peculiarity about them, that they are external to
one another, fixed determinations. Taken in this way, concepts
cease to be concepts.'' ... ``Ihre Bestimmungen
sind nicht so ein Todtliegendes, wie Zahlen und Linien, deren
ihre Beziehung nicht selbst angehört; sie sind lebendige Bewegungen''.

Here there is misunderstanding upon misunderstanding. Whereas
Leibnitz wants the whole cycle of scientific cognitions
expressed by a mode of notation analogous to the mathematical,
Hegel in his objection to this goes so far
that he even regards the mathematical language of signs as so
little corresponding to the demand of logic that only the spoken language
alone can be the language of concepts. ``Da der Mensch die
% --- p. 23 ---
\setcounter{footnote}{0}%
\opage{23}Sprache hat, als das der Vernunft eigenthümliche Bezeichnungsmittel,
so ist es ein müssiger Einfall, sich nach einer unvollkomnern
Darstellungsweise umsehen und damit qvälen zu
wollen. Der Begriff kann als solcher wesentlich nur mit
dem Geiste aufgefaßt werden dessen Eigenthum nicht nur,
sondern dessen reines Selbst er ist. Es ist vergeblich, ihn
durch Raumfiguren und algebraische Zeichen zum Behufe
des äußerlichen Auges und einer \emph{begrifflosen} mechanischen
Behandlungsweise, eines Calculs, festhalten
zu wollen.''\footnote[1]{Wissenschaft der Logik, zweiter Theil, Berlin 1834. Side 59.}

The one-sidedness of this view is conspicuous.
Admittedly the spoken language is a medium indispensable for the development of the concept;
but since ``the concept as such can essentially be
grasped only with the spirit'', it follows manifestly from this
that what matters is not the means as means, but the spiritual
use of the means. Nor is it
figures of space and algebraic signs alone that can be grasped
in a spiritless way; the same can happen with the
words and sentences of language. He who sees in the sentence only a
sum of words, in the word a sum of syllables, in the syllable
a sum of letters, does not see the concept either. Since it
is then not the means in and for itself, but the right use
of the means that is at issue, there arises in this
matter the decisive question whether, for an all-sided and complete
development of the concept, other means besides those the spoken language
offers might not be needed. To answer
this question we must look at:

\begin{center}β) \emph{Thought.}\end{center}

When Hegel shows the unfruitfulness of wanting, like Leibnitz,
to make the syllogistic figures the object of a combinatorial
calculus, and censures a Ploucquet, who by mathematical
determinations of the logical believes he has brought it
so far: \textit{posse etiam rudes mechanice totam logicam}

% --- p. 24 ---
\setcounter{footnote}{0}%
\opage{24}\textit{doceri, ut pueri mathematicam docentur,} then he is
in the right: to transform logic into mechanics of the understanding is
spiritless. Likewise, what Hegel remarks under ``Der Schluß des
Daseyns'' concerning ``the fourth figure: $A$---$A$---$A$ or
the mathematical inference'', namely that the middle term,
\textit{terminus medius,} which here was to be the mediating element,
has no determinateness of its own over against the extremes,
and that it is precisely this tautological character that makes
the inference coincide with a mathematical axiom,
--- is certainly entirely justified. But with all this Hegel has, in
his judgment of the mathematical mode of expression, overlooked
a side which, as regards thought and the concept, is the
actual main side, namely \emph{the punctual.}

To an all-sided development of the categories of quantity
it belongs that the punctual must go together as one with the universal:
here the question is not merely of a great and a greater, here it is asked
also how great? how much greater? The spoken language, which
cannot by the word point out the single in its immediate
existence, cannot either by the word make clear the punctual
determinateness; the spoken language therefore, when thought
demands the punctual, has a deficiency that must be made good by
the mathematical language of signs. Such fundamental equations
as: $f(x + y) = f(x) + f(y)$, and $f(x + y) = f(x).f(y)$
are at the same time propositions that pronounce a concept, hence logical
judgments. What is to be understood by $f(x + y)$ is in
these equations essentially determined in the same way as
when, in the judgment, it is determined by the predicate what
is to be understood by the subject. Insofar as the equation expresses a
law, it is a judgment that asserts the universal, and the language of signs can in this respect be said to be congruent with the spoken language;
but insofar as the equation is a formula, it has
in itself what the spoken language does not have, a possibility of exact
determinations to infinity. What velocity is one can
indeed in a way express in words, but the concept of velocity,
when it is to be
% --- p. 25 ---
\setcounter{footnote}{0}%
\opage{25}clarified with its exact determinations, demands exact analysis. It makes a change of sphere
in cognition when, after having explained in words what
velocity is, one applies the mathematical notation, calls space $s$, time $t$, velocity $v$, and sets $v=\frac{s}{t}$ as the definition of velocity; indeed, in order that the definition
may fit not only uniform but also
non-uniform motion, goes further, sets in place of space $s$ the element of space $ds$, in place of time $t$ the element of time $dt$, and
thus defines velocity by $v=\frac{ds}{dt}$. He who has only
the concept of velocity insofar as it can be presented and
made clear in words has it only in a hovering generality,
and hence in a superficial way. The category ``velocity''
belongs to a circle of concepts that have this peculiarity
about them, that they cannot be made clear by ordinary reflection,
but for their essential and punctual development demand the
intellectual act, the mathematical operation: all
concepts of function are concepts of operation.

\begin{center}γ) \emph{Language and Thought.}\end{center}

If one develops, as Hegel has done in his direction, the
categories of quantity and of measure in logical-dialectical form
in such a way that the moments of operation are touched on only fleetingly, the
development of the concept must become one-sided and in several points, even
as regards the universal, untenable. Concepts that
have their scientific roots in the soil of the exact analyses
are stunted when the exact moments must be missed in
their dialectic. That a dialectical development of the concepts of quantity
does not make up for the lack of insight into the fundamental relations of operation is seen from this, that, whereas there is
indeed a transition from exact analysis to the dialectical, there is
not conversely a transition from the generally logical and
dialectical analysis to the exact: ``\textit{Non datur metaphysica
in mathesin via}''.\footnote[1]{C. Ramus: ``Differential= og Integralregning''. Fortale.}

% --- p. 26 ---
\setcounter{footnote}{0}%
\opage{26}If one sets up, in order to escape further trouble with
the moments of operation, the assertion that the concepts of function corresponding to quantity are to be conceived from two such different main sides that the one concerns exclusively
mathematics, the other philosophy: what follows
from this? Either one must then presuppose that the side that
belongs to mathematics is, as a side of the concept, unessential,
and thus fail to recognize the exact science; or else admit
that the concepts in question do indeed have in mathematics
a side essential for comprehension, yet inaccessible to philosophy.
But if one admits that here in
science there is a side of the concept inaccessible to philosophy, a \textit{cognoscendi genus}, whose essence philosophy
is unable to see through: then one has at once
admitted that philosophy, even if it can
reason a good deal about thinking in general and
inquiry in particular, nevertheless, when it comes
to the point, can furnish neither a general i.e.\ thoroughgoing
theory of cognition nor a comprehensive theory of science.

The way in which the punctual in the concept of function
unites into one with the universal shows
that the presentation cannot dispense with the exact formula.
When Plato's Timaeus wants to explain how it came about that
the god who formed the world ``filled up the double and triple
intervals by cutting off parts of the whole and
placing them between these, so that in every interval
there were two means, of which the one exceeded the first extreme
and was exceeded by the other by the same fraction of
each of them'' etc.; and how he (the god) made
``the circle of the Same move along the side of a parallelogram
from left to right, but that of the Other from right to
left along the diagonal'': then one feels the burdensomeness, % print: Diogonalen (misprint)
not to say the insuperability, of having to describe in words, without the help of
the language of signs, what words cannot describe.
% --- p. 27 ---
\setcounter{footnote}{0}%
\opage{27}If one now sees, on the other hand, the necessity of making
use of the language of signs, and of pointing up by formulas, one sees
at the same time that the presentation ought to be kept so that it is neither
entangled in mathematical operations, nor sets down the formulas
as loose, unintelligible aphorisms. What is meant by this
an example must make clear. In his Logic, in the section
of the doctrine of the Idea that treats of cognition,\footnote[1]{Wissenschaft der Logik, zweiter Theil, Berlin 1834. p.\ 286.} Hegel wants to establish
that it is only in the higher analysis that propositions
become synthetic, and says thus: ``So ist die Aufløsung % print: Aufløsung (ø for ö) (misprint)
der Gleichung $X^{m}-1=0$ mit Hülfe der Sinus ....

\dots{} eine synthetische Auflösung, weil die zu Hülfe genomenen
Bestimmungen, die Sinus\dots{}. nicht eine Bestimmung
der Aufgabe selbst ist.''

But what does it mean that sines must be brought in
when one is to solve the equation: $X^{m}-1=0$? A detailed
application of Euler's and Moivre's formulas, a complete
deduction of the equation:
\[ X^{m}-1=(X-1)(X-\cos\frac{2\pi}{m}-\sin\frac{2\pi}{m}\sqrt{-1})(X-\cos\frac{4\pi}{m}-\sin\frac{4\pi}{m}\sqrt{-1})\ldots, \]
would here transform the train of thought from a philosophical into a
mathematical one, whereas a complete lack of elucidation
would, as one can certainly object in this case to
the Hegelian presentation, prevent the reader who is not especially
occupied with the mathematical from grasping the point.

To what extent the mathematical elucidation must be imparted
must therefore depend essentially on the emphasis that is placed on the development in a
given context:\footnote[2]{If by the cited example Hegel had really wanted to inform the reader what it means that a mathematical proposition becomes synthetic, it would have been easy for him to simplify the presentation. Setting in the equation $X^{m}-1=0$, $m=2$, one gets $X^{2}-1=0$, i.e.\ $X=\pm\sqrt{1}=\pm 1$; but $\cos 0+\sin 0\sqrt{-1}=+1$, and $\cos\frac{2\pi}{2}+$ $\sin\frac{2\pi}{2}\sqrt{-1}=-1$; setting further $m=3$, one gets $X=\sqrt[3]{1}=1^{\frac{1}{3}}=+1$, $\frac{-1+\sqrt{-3}}{2}$, $\frac{-1-\sqrt{-3}}{2}$; but $\cos 0+\sin 0\sqrt{-1}=+1$, $\cos\frac{2\pi}{3}+\sin\frac{2\pi}{3}\sqrt{-1}=\frac{-1+\sqrt{3}\ \sqrt{-1}}{2}$, and $\cos\frac{4\pi}{3}+\sin\frac{4\pi}{3}\sqrt{-1}=\frac{-1-\sqrt{-3}}{2}$. The reader therefore needs only general preliminary knowledge in order to see how it comes about that sine and cosine are introduced in the solution of the stated equations. Had the philosophical task, on the other hand, been, by the expression $X^{m}-1=0$, to make clear by way of example that the roots of an algebraic equation can be reduced to the form $a\pm b\sqrt{-1}$, which for $b=0$ also includes the real roots, and thereby to elucidate what it means that the imaginary roots occur in pairs, so that each pair can be regarded as roots of an algebraic equation of the 2nd degree ($(x-a)^{2}+b^{2}=0$), and the equation thus show itself as composed of nothing but real factors of at most the 2nd degree: then the chosen expression would have had to be illuminated from another side.} as the thought
is, so must the language be.
% --- p. 28 ---
\setcounter{footnote}{0}%
\opage{28}
\begin{center}II.\end{center}

\begin{center}The Problem of Infinity.\end{center}

In the function the unity of the constant and the
variable magnitude is posited in such a way that the possibility of a development
of the relation between the finite and the infinite
is thereby given. In that the magnitude is posited as constant, it is posited
with a being limit, and conceived as finite; in that the magnitude
is posited as variable, it is posited with a becoming limit,
and indicates a transition to the infinite. The changes
whose possibility the function indicates are countless, and
show us the phenomenal side of infinity; the law to which the changes
are subject in accordance with the function is the immutable in itself,
and presents the infinite from the essential
side. ``The problem of infinity is the great common problem of mathematics and
% the note 2) of p. 27 runs over onto the foot of this page: set in full at its mark
% --- p. 29 ---
\setcounter{footnote}{0}%
\opage{29}philosophy. But before one can properly think
of the solution of a problem, one must first of all see to it that
the problem comes to be. The present problem is now of
such a nature that it can stand clear before consciousness only
when mathematics will help philosophy to
set it forth.''\footnote[1]{Phil. og Mathem. p.\ 20.}

In order that the problem of infinity may become clear, and
the conditions of its solution be seen, it will chiefly
be a matter of clarifying the following points: the difference of the
finite from the infinite, the approximation of the finite to the
infinite, and the way in which this approximation
lets itself be articulated.

\begin{center}a) The Finite in Its Immediate Difference from the Infinite.\end{center}

For immediate consideration the finite shows itself
as the being, the positive, the infinite as the non-being,
the negative; there is in this respect a qualitative
limit between the finite and the infinite. By holding fast
to the merely qualitative limit, however, it would
be impossible to develop the concept of the infinite; robbed of
every reflected image of finitude, the infinite coincides
with empty negation: the problem cannot come to be.
If, on the other hand, the finite is conceived not merely as a being,
but as a being magnitude, as the magnitude whose
limits can be exceeded, then the infinite is no
longer an absolute negation of magnitude, but a negation
of what is finite in magnitude, a negation of the
finite magnitude, a negation of magnitude that essentially
belongs to the concept: magnitude. In that the qualitative
limit thus shows itself as quantitating, the
problem arises.

For if the infinite is conceived not as absolute negation of magnitude,
but as negation of the finite magnitude,
% --- p. 30 ---
\setcounter{footnote}{0}%
\opage{30}this infinite must show itself to immediate intuition
from entirely opposite sides. As that which in greatness and
smallness surpasses the finite, it falls outside finitude;
as that into which every finite magnitude can be
resolved, it falls inside, indeed coincides with, finitude.
By seeing through the opposed views, reflection discovers
the dialectical unity of the finite with the infinite.

\begin{center}α) \emph{The Infinite outside the Finite.}\end{center}

How the infinite, in that it surpasses every finite magnitude in greatness and
smallness, comes to
fall outside finitude, and to maintain itself as a beyond,
is made intuitively evident by series whose terms are functions of the
place that each of them occupies in its determinate series. The equations for
the arithmetic, the harmonic and the geometric series:

\[ u_{n}=f_{1}(n)=a+nx, u_{n}=f_{2}(n)=\frac{1}{a+nx}, u_{n}=f_{3}(n)=ax^{n} \]
show that between the greatest possible finite term and
the infinitely great, as also between the smallest possible
finite term and the infinitely small, there is an infinite distance.
In that the limit between the finite and the infinite is
at once both qualitative and quantitative, a transition forms
here, which, however, so long as the two determinations
are held immediately against each other, is no transition.

In the equation $u_{n}=f(n)=ax^{n}$ it admittedly makes in one respect
a difference whether one excludes between $ax^{0}$ and $ax^{\infty}$
every finite term, or places millions of times
millions of intermediate terms; but this difference is nonetheless
entirely vanishing. Insofar as the limit between the
finite and the infinite is quantitative, it seems indeed
at first glance as if the many intermediate terms filled up the
chasm, and brought the finite nearer to the
infinite; but when the qualitative in the limit:
``either a finite or an infinite'', is then stressed, the
millions of terms that were to fill up the chasm lose their significance
as transitional terms. Instead of forming a transition from
% --- p. 31 ---
\setcounter{footnote}{0}%
\opage{31}the finite to the infinite, the many terms denote only
a continuation of the finite; when $ax^{\infty}$, $ax^{-\infty}$ are posited, it is posited expressly by a breaking off,
a negation, a leap. The qualitative limit that
separates the one finite from the other, white from black,
day from night, land from sea etc.\ can, so long as here in
quantitative respect only a finite is reflected on, be
now nearer now farther: He who is only a mile from the
coast is naturally nearer the limit than he who is
many miles out at sea; but when the limit is between
finite and infinite, hence infinitely quantitating, it is
as qualitative limit everywhere equally near and equally far:
the negation by which the finite is broken off, and the
infinite posited, is for each finite term equally unconditioned,
equally negative. That there is here a transition which is no transition
shows that the stated view is unquiet, unsatisfying.
The unsatisfying character is seen, among other things, from this, that
one, instead of letting the infinite fall outside the finite,
with just as much right can choose the opposite view,
and seek:

\begin{center}β) \emph{The Infinite within the Finite.}\end{center}

How the transition from finite to infinite, i.e.\
from finite to $0$, and from finite to $\infty$ coincides
with a corresponding transition from finite to finite,
the trigonometric functions can make intuitively evident. If, e.g., $x=90^{\circ}$, then $\cos x=0$, $\sin x=1$, $\mathrm{tg}\,x=\infty$ etc.
If now, as regards $\angle x$, a transition is made from $89^{\circ}$ to $90^{\circ}$, then $\sin x$ makes a transition from finite to finite, $\cos x$ from finite to $0$, and $\mathrm{tg}\,x$ from finite to $\infty$. Indeed,
the transition from finite to finite is itself a transition in
an infinite. Between $\angle x=89^{\circ}$ and $\angle x=90^{\circ}$ lie
indeed infinitely many infinitely small angles, all of which must
be traversed if the transition is to take place: the finite angle
is, like every finite magnitude in general, divisible to
infinity. By virtue of infinite divisibility the
% --- p. 32 ---
\setcounter{footnote}{0}%
\opage{32}infinite becomes enclosed by finite limits; in that the parts become
infinitely small, their number becomes infinitely great: the finite
limits enclose the infinitely great, precisely because they
enclose the infinitely small. What, then, is the finite?
A product of infinity, a synthesis of the infinitely
small with the infinitely great. And the infinite? A
product of finitude, a carried-through unfolding and thereby
a carried-through dissolution of the finite. How
the finite, in that it shows itself to be without end, can
bring forth the infinite is seen, e.g., from a quotient series such as

\[ \frac{1}{1-1}=1+1+1+.... \]
for even when the series
is thought continued to infinity, we still have
\[ \frac{1}{1-1} \]
as remainder
and must begin again from the start. The infinite arises through
the unceasing repetition of finite terms: in that the finite by means of
divisibility is dissolved into the infinite, the
infinite is at the same time posited by means of the finite.
From this follows:
\begin{center}γ) \emph{The Dialectical Unity of the Finite and the Infinite.}\end{center}

``Formal logic holds fast to the principle
that the finite is not infinite, and the infinite
not finite. But however irrefutable these negative judgments, guaranteed
by the principle of contradiction, may seem to
be, just as poor and empty of content are the corresponding positive ones:
``the finite is the finite, the infinite the infinite'',
for their so-called identity is a pure tautology. As
dialectical science philosophy has not been able to remain
standing at such tautologies. That the infinite is not
finite means, in other words, that it is a limit
for the finite, hence the not-finite. The finite,
which is the not-infinite, is then at the same time a limit
for the infinite. But when the finite and the infinite
in this way are each other's limits,
% --- p. 33 ---
\setcounter{footnote}{0}%
\opage{33}the determination of the limit must become dialectical. For if the infinite has
its limit in the finite, then it has a limit of finitude,
a finite limit, and is itself finite.
If, conversely, the finite reaches its limit only in the infinite,
then it has no finite limit, and is consequently infinite.''\footnote[1]{Philos. og Mathem. p.\ 21—22.}

Insofar as the infinite surrounds the finite, the
limit is determined by the infinite; insofar as the finite
surrounds the infinite, the limit is determined by the finite;
that this two-sided determination necessarily belongs
to the thoroughly reflected limit is easy to show. An
angle of $90^{\circ}$ can be divided into infinitely many infinitely small
angles, an angle of $1^{\circ}$ likewise; a square mile can
be divided into infinitely many infinitely small squares, a square inch
likewise: what follows from this? If the emphasis is laid
on the dissolution of the finite into the infinite, then
the consequence indeed seems to be that there is here no longer any difference between
the angle of $90^{\circ}$ and the angle of $1^{\circ}$, that the part becomes
equal to the whole, that the limit, in other words, is effaced.
But if the emphasis is then laid again on the existence of the infinite
in and through the finite, then there is not merely a difference
between finite and finite, between the part and the
whole, but there shows itself at the same time a determinate difference
between infinite and infinite, with respect to the
infinitely small as well as to the infinitely great. The one infinite
number of infinitely small angles can then be as different
from the other as $1^{\circ}$ from $90^{\circ}$; the one infinite
number of infinitely small squares as different from the
other as a square inch from a square mile. That
the limit in the finite is canceled means that it is posited in
the infinite, that the infinite takes the limit up into itself
means that the finite is posited. The dialectical unity shows
itself in this, that the limit belongs just as much to the infinite
as to the finite.
% --- p. 34 ---
\setcounter{footnote}{0}%
\opage{34}
\begin{center}b) The \emph{transition} of the finite into the infinite: \emph{approximation.}\end{center}

If the limit between the finite and the infinite were merely qualitative, no transition could be given here, but only a breaking off; were the limit merely quantitative, the breaking off posited by negation would evaporate, and transition would lose its significance. The category \emph{transition} presupposes a dialectical unity of the qualitative and the quantitative, which finds its development in \emph{approximation.} Let $A$ be, e.g., separated from $B$ by a qualitative limit; $A$ is thus not $B$, $B$ not $A$; insofar as the limit is exclusively qualitative, there is here nothing intermediate: where $A$ ceases, $B$ begins, here there is no transition, no approximation. If the limit between $A$ and $B$ becomes at the same time quantitative, on the other hand, something intermediate shows itself; that is to say: one can, if the quantifying element is a difference of degree, strengthen $B$ by weakening $A$, and conversely; or, if the quantifying element is a relation of distance, approach $B$ by moving away from $A$, and conversely. Since here there is a determination of quality, there is thus a determinate opposition; since here in the qualitative there is a quantifying, there is between the opposite terms something intermediate; since here there is something intermediate, there is a transition, and in the transition an approximation.

Since there is now given here not only a qualitative difference between finite and infinite, but within the bounds of finitude also a qualitative difference among the finite magnitudes themselves, approximation can, as a consequence, be partly an approximation of the finite by the finite to the finite, and thereby to an attainable; partly an approximation either of the finite by the infinite to the finite, or of the finite by the finite to the infinite, and thereby to an unattainable. By the infinite approximation, however, the unattainable is attained.

\begin{center}α) An attainable: approximation of the finite to the finite.\end{center}

When Diogenes, disputing with the Eleatics
% --- p. 35 ---
\setcounter{footnote}{0}%
\opage{35}about motion, walked, so it is said, back and forth, in order thus to prove in deed the reality of motion, that was a superficial manner of proof. By walking back and forth he showed only, in the sensible, a transition from finite to finite, an approximation to an attainable, which was not really in dispute. Here, on such a path, there are e.g. 12 steps from the point of departure to the goal; when 1 step has been taken, only 11 steps remain; when 11 steps have been taken, only 1 step remains; 1 step is a finite length, 12 steps is a finite distance; he who has already taken 11 steps is naturally 10 steps nearer the goal than he who has taken only 1 step: motion in the finite, the finite approximation, is something so much a matter of course that there can as yet be no question here of any problem. The problem shows itself only when infinity enters. In the finite an approximation can undeniably begin, if one will but take the first step; but precisely in relation to infinity it holds that \textit{c'est le premier pas qui coute,} for in the first step lies the whole problem. In order to take a whole step, one must first take a half step; in order to take a half step, one must first take the half of a half step; before the half of a half step comes the half of the half of the half --- ad infinitum --- of a half step. It is thus not subtlety but logical consequence when the problem of motion is traced back to an infinitely small part of the first step, the limit between the zero point of the beginning and the smallest finite stretch of the way.

\begin{center}β) An unattainable: approximation of the finite to the infinite.\end{center}

The infinitely large is too large for finitude, for finitude contains only the finitely large; the infinitely small is too small for finitude: the infinite is something unattainable for finitude. Can there then be an approximation to an unattainable? That there is an approximation to
% --- p. 36 ---
\setcounter{footnote}{0}%
\opage{36}an attainable goes without saying. ``Approximation to an attainable involves no contradiction; an unattainable without approximation is likewise no contradiction. There is thus nothing contradictory in the proposition that parallel lines, however far they are extended, will never intersect. When the mathematician sometimes uses the expression of parallel lines that, extended to infinity, they will intersect: one cannot thereby really think of any approximation; it is an impossibility, here presented under the form of possibility. Approximation to an unattainable, on the other hand, seems to present the hardest of all contradictions. By approximation one must surely in one sense or another come nearer to the goal; but to come nearer to the goal must surely be the same as to come nearer to attainment; but from the attainment of the unattainable one is surely always equally far away.''\footnote[1]{Phil. og Mathem. p. 25—26}

Should the concepts ``unattainable'' and ``approximation'' be brought into logical connection, one might, in accordance with ``\textit{diversus respectus},'' perhaps begin by conceiving the matter thus: the category ``approximation'' must always be seen from two standpoints: 1) the point from which one proceeds, the point of beginning, and 2) the point that is to be reached, the goal, the end point. If there is any motion at all from a determinate point of departure toward a goal, then, if the goal is infinitely far away, one may in a sense readily concede that every approximation is as a nothing against the infinite remoteness of the goal, and yet, when one looks back, regard the continued removal from the point of departure as an approximation to the goal. Here there is thus in one respect, namely insofar as one looks forward, no approximation, but yet in another respect, when one, mark well, looks back, a real approximation.'' According to this explanation one should thus see approximation as a
% --- p. 37 ---
\setcounter{footnote}{0}%
\opage{37}removal, and look forward by looking back; which is a contradiction.

\begin{center}γ) The infinite approximation.\end{center}

Approximation to an unattainable is a contradiction, namely when one presupposes that it should be able to stop at this or that finite point. When the goal of the approximation is attainable, it is completed by the attainment of the goal, and thus gets an end; but the approximation that is directed toward the unattainable can never end, its completion is its uninterrupted continuation: it is and must be \emph{infinite.} That the transition from the finite to the infinite can be determined by infinite approximation has its ground in the essence of finitude. For the finite can continue itself in a twofold manner such that precisely by this its continuation it sublates itself, and passes over into the infinite.

In the series $1+1+1+\ldots$ corresponding to the quotient $\frac{1}{1-1}$ the finite is seen to continue itself in finite terms without end; if one stops here at an arbitrary term, the number of terms may be as large as one likes, it is nevertheless only a finite number: the series has only finite value, here there is no approximation to the infinite. If the series is continued, on the other hand, without breaking off, then the continuation is an unceasing exceeding of the finite limit, but a continued exceeding of the finite limit is an approximation to the infinite, an \emph{infinite} approximation.\footnote[1]{By this infinite approximation no regard is yet taken to the difference between \textit{convergent} and \textit{divergent} series, but solely and alone to the finite's own infinitization.} In the series $\frac{1}{1-\frac{1}{2}}=1+\frac{1}{2}+\left(\frac{1}{2}\right)^{2}+\ldots$ the whole magnitude is enclosed between finite limits, since $\frac{1}{1-\frac{1}{2}}=2$; but for the diminutions of the terms there is here no existing limit, as little as there is here
% --- p. 38 ---
\setcounter{footnote}{0}%
\opage{38}in the finite any limit for the number of terms. That the limit of infinity is unattainable is signified thereby that every finite limit must be exceeded; but, in that the limit is exceeded, there is here yet an approximation to the infinite, insofar as the approximation itself is infinite.

In geometric form the infinite approximation can be made intuitable by the relation between a branch of a hyperbola and its asymptote. To the hyperbola's equation: $y=\frac{b}{a}\sqrt{x^{2}-a^{2}}$ corresponds the asymptote's equation: $y_{1}=\frac{b}{a}x$, and the equation: % print: Assymptotens (misprint)
\[ y_{1}-y=\frac{ba}{x+\sqrt{x^{2}-a^{2}}} \]
to the distance between the two, vanishing by degrees.. In such an example one sees the peculiar relation between intuiting and thinking, a relation which is decisive for the problem of infinity and its solution. Insofar as the curve with its asymptote % print: Assymptote (misprint)
is presented to intuition, the assertion that these lines can never, under finite continuation, possibly touch each other must appear as a paradox; but the equation cited for $y_{1}-y$ proves that this paradox is an incontrovertible truth. Here there is a continued approximation, for the greater $x$ becomes, the smaller $y_{1}-y$ becomes; but since $y_{1}-y$ can become $0$ only when $x$ becomes $\infty$, the goal is unattainable, and the approximation infinite.

\begin{center}c) Approximation in punctual development.\end{center}

The finite is in itself something incomplete, its completion a transition into the infinite; as the negation of the finite, as abstraction, the infinite is likewise something incomplete, its completion a progression of the finite; the reflected unity of the finite and the infinite is thus not a being but a \emph{becoming}: in \emph{becoming} the infinite approximation has its categorical expression. But a general becoming without an existence in which the becoming takes place is
% --- p. 39 ---
\setcounter{footnote}{0}%
\opage{39}meaningless, and a general existence without existing terms impossible; if therefore the becoming in which the finite and the infinite are united as in infinite approximation is to be made clear, then the concept of this becoming must show itself as a concept of punctuality.

The concept of punctuality can, to be sure, as a concept not demand an express enumeration of all single terms and relations in all single cases, for, since such an enumeration is impossible, the universal would be lost, and the concept vanish; on the other hand punctuality can and must demand that the possibility of a determination of the single terms in their peculiar relation reflect itself in the general expression as in an exact formula. The transitions in which the becoming of the finite and the infinite is punctualized are then to be apprehended both from the mathematical and from the logical side, and their lawful multiplicity pointed out, not merely in functions, but in systems of functions.

\begin{center}α) Mathematical-logical transitions.\end{center}

The transitions of conceptual development can in general be characterized thereby that from a given presupposition a consequence is drawn out, and that in and with the consequence a new determination is posited. If the emphasis falls on the development of the consequence from the presupposition, the transition is analytic; if the emphasis falls on the new, specially added determination, the transition is synthetic: the analytic stands to the identity prevailing in difference as the synthetic to the difference prevailing in identity. With identity and difference all thinking is occupied; but in logic otherwise than in mathematics. In logic thinking is \emph{reflection}, in mathematics \emph{operation}; if now the analytic has preponderance over the synthetic, then thinking becomes in the logical sense \emph{objective} reflection, insofar as the mathematical action becomes a transformation; if conversely the synthetic has preponderance over the analytic, then
% --- p. 40 ---
\setcounter{footnote}{0}%
\opage{40}thinking becomes, logically viewed, \emph{teleological} reflection, insofar as the mathematical action becomes a combination.

\begin{center}αα) The exact application of analysis and synthesis: operation.\end{center}

According to the Kantian critique of reason pure mathematics rests on ``synthetic judgments a priori.'' Propositions such as: ``$5+7=12$''; ``the straight line is the shortest way between two points,'' etc., are grounded on a ``synthesis a priori.'' According to the Hegelian logic the arithmetical theorems are, on the other hand, with few exceptions, formally analytic. The truth is that mathematical cognition must, as follows necessarily from the essence of operation, be analytic-synthetic % print: analythisk (misprint)
from first to last.\footnote[1]{The question of the relation between the analytic and the synthetic judgment does not, to be sure, coincide immediately with the question of the relation of the analytic method to the synthetic; nevertheless the dialectical holds in equal degree for both relations. For closer elucidation we let \emph{Kant} and Hegel speak more fully. Kant (Kritik der rein. V. Gesammtausgab. Leipzig 1838. Einleitung p. 46—48): „Man sollte anfänglich zwar denken: daß der Satz $7+5=12$ ein bloß analytischer Satz sei, der aus dem Begriffe einer Summe von Sieben und Fünf nach dem Satze des Widerspruchs erfolge. Allein, wenn man es näher betrachtet, so findet man, daß der Begriff der Summe von 7 und 5 nichts weiter enthalte, als die Vereinigung beider Zahlen in eine einzige, wodurch ganz und gar nicht gedacht wird, welches diese einzige Zahl sei, die beide zusammengefaßt. Der Begriff von Zwölf ist keineswegs dadurch schon gedacht, daß ich mir jene Vereinigung von Sieben und Fünf denke, und ich mag meinen Begriff von einer solchen möglichen Summe noch so lange zergliedern, so werde ich doch darin die Zwölf nicht antreffen. Man muß über diese Begriffe hinausgehen, indem man die Anschauung zu Hülfe nimmt, die einem von beiden correspondirt, etwa seine fünf Finger, oder (wie Segner in seiner Arithmetik) fünf Puncte, und so nach und nach die Einheiten, der in der Anschauung gegebenen Fünf zu dem Begriffe der Sieben hinzuthun. (Cf. Phil. Propæd. p. 180—181). Hegel, on the other hand (Wissensch. der Log. 2te Theil. Berlin 1834. p. 282—83 ff.): „Bekanntlich wird die Arithmetik und die allgemeineren Wissenschaften der discreten Größe vorzugsweise analytische Wissenschaft und Analysis genannt .... der Fortschritt ist die Reduktion des Ungleichen auf immer größere Gleichheit. Um ein Beispiel an den ersten Elementen zu geben, so ist die Addition das Zusammenfassen ganz zufällig ungleicher Zahlen, die Multiplikation dagegen von gleichen, worauf noch das Verhältniß der Gleichheit von der Anzahl und der Einheit folgt, und das Potenzen-Verhältniß eintritt.“ „Weil nun die Bestimmtheit des Gegenstandes und der Verhältnisse eine gesetzte ist, so ist die weitere Operation mit ihnen auch ganz analytisch, und die analytische Wissenschaft hat daher nicht sowohl Lehrsätze, als Aufgaben. Der analytische Lehrsatz enthält die Aufgabe schon für sich selbst als gelöst, und der ganz äußerliche Unterschied, der den beiden Seiten, die er gleich setzt, zukommt, ist so unwesentlich, daß ein solcher Lehrsatz als eine triviale Identität erscheinen würde. Kant hat zwar den Satz $5+7=12$ für einen synthetischen Satz erklärt .... Allein wenn das Analytische nicht das ganz abstract Identische und Tautologische $12=12$ bedeuten, und ein Fortgang in demselben überhaupt seyn soll, so muß irgend ein Unterschied vorhanden seyn, jedoch ein solcher, der sich auf keine Qualität, keine Bestimmtheit der Reflexion und noch weniger des Begriffs gründet. $5+7$ und 12 sind durchaus ganz derselbe Inhalt .... Hier ist im Geringsten kein Uebergang zu einem Andern; es ist ein bloßes Fortsetzen, d. h. Wiederholen derselben Operation, durch welche 5 und 7 entstanden ist.“ „Der Beweis eines solchen Lehrsatzes — einen solchen erforderte er, wenn er ein synthetischer Satz wäre — würde nur in der Operation des durch 7 bestimmten Fortzählens von 5 an, und in dem Erkennen der Uebereinstimmung dieses Fortgezählten mit dem bestehen, was man sonst 12 nennt, und was wieder weiter nichts, als eben jenes bestimmte Fortzählen selbst ist“ .... „Es ist daher ein höchst überflüssiges Gerüste, hier die Form der geometrischen Methode, welche sich auf synthetische Sätze bezieht, anzuwenden und der Aufgabe außer der Auflösung auch noch einen Beweis folgen zu lassen. Er kann nichts als die Tautologie ausdrücken, daß die Auflösung richtig ist, weil man operirt hat, wie aufgegeben war. Wenn die Aufgabe ist, man soll mehrere Zahlen addiren; so ist die Auflösung: man addire sie; der Beweis zeigt, daß die Auflösung richtig ist, darum weil aufgegeben, war zu addiren, und man addirt hat.“ — As regards synthetic cognition, Hegel (work cited, p. 288—313) lays particular emphasis on the relation between the definition, the division and the theorem. „Die Definition enthält nur Eine Bestimmtheit, die Eintheilung die Bestimmtheit gegen andere; in der Vereinzelung ist der Gegenstand in sich selbst aus einander gegangen. Insofern die Defintion beim allgemeinen Begriffe stehen bleibt, so ist dagegen in den Lehrsätzen der Gegenstand in seiner Realität, in den Bedingungen und Formen seines reellen Daseyns erkannt“ ... „Der Lehrsatz nun nach der angegebenen Bestimmung ist das eigentlich Synthetische eines Gegenstandes, insofern die Verhältnisse seines Bestimmtheiten nothwendig, das ist, in der innern Indentität des Begriffes gegründet sind. Das Synthetische in der Definition und Eintheilung ist eine äußerlich aufgenommene Verknüpfung; das Vorgefundene wird in die Form des Begriffes gebracht, aber als vorgefunden wird der ganze Inhalt nur monstrirt; der Lehrsatz aber soll demonstrirt werden.“ .... „Das glänzende Beispiel der synthetischen Methode ist die geometrische Wissenschaft (p. 313). .... „Wenn man die Lehrsätze“ — it says on p. 307 — „einer synthetischen Wissenschaft, und namentlich der Geometrie, näher vergleicht, so wird sich dieser Unterschied zeigen, daß einige ihrer Lehrsätze nur einzelne Verhältnisse des Gegenstandes enthalten, andere aber solche Verhältnisse, in welchen die vollständige Bestimmtheit des Gegenstandes ausgedrückt ist. Es ist eine sehr oberflächliche Ansicht, wenn die sämmtlichen Sätze an Werth einander gleichgeachtet werden, weil überhaupt jeder eine Wahrheit enthalte, und im formellen Gange, im Zusammenhange des Beweisens gleich wesentlich sey. That this conception of the synthetic method, resting on the separation of definition, division and theorem, applies, however, just as well to arithmetic as to geometry, is among other things obvious from the fact that arithmetic needs definitions, divisions and theorems just as much as geometry. In Elementær Algebra ved C. Ramus, Kbhvn 1855, in particular in ``Second Section. Fundamental propositions on whole numbers,'' and in ``Third Section. Irrationality and approximation; decimal fractions,'' arithmetical theorems occur, and the method is thus to that extent synthetic. If now the synthetic method is applicable in arithmetic, then the analytic method, which modern science owes chiefly to Cartesius, is conversely also applicable to geometry. What on the one side is set up as a theorem, and treated synthetically, can then on the other side be presented as a problem, and treated analytically. ``The line which bisects the line AB and stands perpendicular to it is the geometrical locus of the points which have equal distance from the two fixed points A and B.'' That one, in proving this theorem from the congruence of right-angled triangles, should arrive at a more content-rich cognition than by transforming the theorem into a problem and solving the problem analytically, is not to be seen. On the contrary! when the problem is treated analytically, when one chooses a rectangular coordinate system, e.g. that which has AB for X-axis and A for origin, and lets the point (x, y) be one of the sought points, one has, for AB=a, $x^{2}+y^{2}=(a-x)^{2}+y^{2}$, and in the equation $x=\pm(a-x)$ a double solution, namely $x=\frac{a}{2}$, which designates a straight line through the midpoint of AB and perpendicular to it, and $x=\infty$, i.e. a straight line likewise perpendicular to AB, but at an infinite distance from the points A and B. While now the synthetic theorem stated only one line as the geometrical locus, analysis has on the other hand shown that there can here be two lines; the analytic method has thus in this case given a richer cognition than the synthetic. It will then without doubt, the more carefully one examines the matter, become ever more evident that, insofar as one does not let the separation of analytic and synthetic method hold by tradition and convention, but looks to the essence of cognition, the two methods cannot hold themselves closed off against each other, but the analytic must, even where it predominates, everywhere presuppose the synthetic, and conversely, the synthetic the analytic.} The equation: $a+a=2a$ % print: retvinklene (misprint) % print: Defintion, Indentität (Hegel as printed) (misprint)
% --- p. 41 ---
\setcounter{footnote}{0}%
\opage{41}is at bottom an analytic judgment, which says that multiplication is only an addition; but since multiplication
% the note 1) of p. 40 runs over onto the foot of this page: set in full at its mark
% --- p. 42 ---
\setcounter{footnote}{0}%
\opage{42}is nevertheless a peculiar action, a distinctive function, the judgment is at the same time synthetic. The same holds of
% the note 1) of p. 40 runs over onto the foot of this page: set in full at its mark
% --- p. 43 ---
\setcounter{footnote}{0}%
\opage{43}$a\cdot a=a^{2}$ etc. How the synthetic can grow together with the analytic under a progressing operation can be expressly shown. In an equation such as: $a^{x}=A_{0}+A_{1}x+A_{2}x^{2}+\ldots$ the relation between the coefficients can be
% the note 1) of p. 40 runs over onto the foot of this page: set in full at its mark
% --- p. 44 ---
\setcounter{footnote}{0}%
\opage{44}found out only when one places beside it a \emph{second} equation of the same form. $a^{z}=A_{0}+A_{1}z+A_{2}z^{2}+\ldots$ If the two equations are now multiplied by each other in such a way that the product receives a double form, namely: $a^{x}a^{z}=A_{0}+A_{1}^{2}xz+A_{2}^{2}x^{2}z^{2}+\ldots.$ and $a^{x+z}=A_{0}+A_{1}(x+z)+A_{2}(x+z)^{2}+\ldots.$, then one sees already from this how the synthetic during the operation has been strengthened by the continued analysis. The cooperation of synthesis becomes especially noticeable when, after having developed the one term $A_{2}(x+z)^{2}$, one compares $2A_{2}xz$ with $A_{1}^{2}xz$ in the first series, for $2A_{2}$ and $A_{1}^{2}$ then present themselves as coefficients which belong each to its own series, with such an immediate independence over against each other that, in order to prove their identity, one must apply the method of undetermined coefficients.

As regards the transition of the finite into the infinite itself, the punctual in approximation rests essentially on the relation between analysis and synthesis. In the quotient $\frac{1}{1-x}$ the finite terms $1$ and $x$ are so combined (synthesis) that an infinite series can be developed from them (analysis). The intellectual action of development, operation, is ``a synthesis \textit{a priori}''; but a synthesis \textit{a priori} is an analyzing synthesis: the a priori is analytic.

\begin{center}ββ) Analysis on a synthetic basis: transformation.\end{center}

If now it is the same content that is presented in different form, analysis has preponderance over synthesis. In the equation: $\frac{a+b}{c+d}=\frac{1}{c+d}(a+b)$ the one expression is only a rewriting of the other; even in the equation $\frac{a+b}{c+d}=\frac{1}{c}(a+b)\left(1-\frac{d}{c}\ldots\right)$ the analytic prevails. Insofar as analysis is applied to the solution of a determinate problem, an interchange of the different modes of expression is nevertheless more than an
% --- p. 45 ---
\setcounter{footnote}{0}%
\opage{45}arbitrary rewriting; it is a reshaping that conditions the development of the content, an essential remodeling, a transformation. The quotient $\frac{a+b}{c+d}$ expresses a determinate relation between the magnitudes $a, b, c$ and $d$; the product $\frac{1}{c}(a+b)\left(1-\frac{d}{c}\ldots\right)$, in which an infinite approximation is indicated, expresses the relation between the same magnitudes in another manner, and thereby contributes to giving the relation a complete transparency; the relation cleared to complete transparency is \emph{objective} reflection.

How transformation serves objective reflection, and conditions insight, a closer treatment of the series: $a^{x}=A_{0}+A_{1}x+A_{2}x^{2}+\ldots$ can show. That the series proceeds by increasing powers of $x$ is so evident that one can in this respect immediately give the $n$th term as $A_{n-1}x^{n-1}$, and in this one expression find the countless powers of $x$ of the series reflected. Such a reflection does not, however, yet take place as regards the coefficients $A_{0}, A_{1}, A_{2}, \ldots$ That there must here be a relation between these coefficients and $a$ in the expression: $a^{x}$ may be presupposed; but in order that this relation may be shown and cleared in reflection, transformation must be applied. By setting $x=0$, and finding $a^{0}=A_{0}=1$, one can then first transform the series to $a^{x}=1+A_{1}x+A_{2}x^{2}+\ldots$; by next, in consequence of our preceding $A_{2}=\frac{A_{1}^{2}}{1.2}$ setting $A_{3}=\frac{A_{1}^{3}}{1.2.3}$, $A_{m}=\frac{A_{1}^{m}}{1.2.3\ldots m}$, one can transform the series to $a^{x}=1+A_{1}\frac{x}{1}+A_{1}^{2}\frac{x^{2}}{1.2}+A_{1}^{3}\frac{x^{3}}{1.2.3}+\ldots$, and in this manner see the countless coefficients of the series reflected in $A$; if one now sets, further, $x=\frac{1}{a}$, and whereby % print: x=\frac{1}{a} (misprint)
the series changes to: $a^{\frac{1}{A_{1}}}=1+\frac{A_{1}}{1}\left(\frac{1}{A_{1}}\right)+\frac{A_{1}^{2}}{1.2}\left(\frac{1}{A_{1}}\right)^{2}+\ldots$
% --- p. 46 ---
\setcounter{footnote}{0}%
\opage{46}$=1+\frac{1}{1}+\frac{1}{1.2}+\frac{1}{1.2.3}+\ldots=e$, and $A_{1}=\frac{\log a}{\log e}$, or, in the natural system, $A_{1}=l.a$, then the series is finally transformed to $a^{x}=1+l.a\frac{x}{1}+(l.a)^{2}\frac{x^{2}}{1.2}+(l.a)^{3}\frac{x^{3}}{1.2.3}+\ldots$; and one then has a form which raises the relation to all-sided reflection.\footnote[1]{That the reflection is now carried through is seen also from the fact that all particular terms with their peculiar determinacy are mirrored in the $n$th or general term: $l(a)^{n-1}\frac{x^{n-1}}{1.2\ldots(n-1)}$.}

Since the relation between the finite and the infinite is punctually determined by approximation, it hovers dialectically between representation and reflection. The particular terms of the infinite series exist for representation, and in representation the finite prevails; but when the all-sided relation has been demonstrated, and the countless terms of the series are seen to reflect themselves both in $a^{x}$ and in one another, infinity comes to light: the infinite does not exist for representation, it is only in reflection for pure insight.

\begin{center}γγ) Synthesis on an analytic basis: combination.\end{center}

If analysis now opens up for insight such a multitude of possibilities that synthesis goes hand in hand with the choice of those cases that are just needed for the solution of the problem, reflection becomes \emph{teleological.} If one has thus assured oneself by analysis that the coefficients in the series: $a^{x}=A_{0}+A_{1}x+A_{2}x^{2}+\ldots.$ are all contained in the formula: $A_{m}=\frac{A_{1}^{m}}{1.2.3\ldots m}$, then a determination of the relation in which $A_{1}$ stands to $a^{x}$ is the nearest problem. The solution of this problem is thus an end in view; but for the attainment of the end it holds here, among millions of possible cases, to choose just the case that yields the means; the comprehensive problem
% --- p. 47 ---
\setcounter{footnote}{0}%
\opage{47}entails a multiplicity of smaller problems: the expedient of setting $x=\frac{1}{A_{1}}$ is chosen from among countless possibilities, and presupposes a consciousness hovering above them.

Since teleological reflection thus determines the exact action, and is absorbed in it, operation is combining. If the heterogeneous series: $e^{x}=1+\frac{x}{1}+\frac{x^{2}}{1.2}+\ldots.$ $\sin x=x-\frac{x^{3}}{1.2.3}+\frac{x^{5}}{1.2.3.4.5}-\ldots$,
$\cos x=1-\frac{x^{2}}{1.2}+\frac{x^{4}}{1.2.3.4}-\ldots$, are each developed and determined by an analysis predominating over synthesis; then these series, as Euler has shown, can in turn be combined by a synthesis predominating over analysis. In combination the synthetic prevails; but the similarities, possibilities and intermediate terms on which the transitions rest must be tested and justified by analysis. The intermediate term which combination here requires,\footnote[1]{While the starting points peculiar to the development are indicated and fixed by combining synthesis, the development itself rests, on the other hand, on analysis. Under combination the living standpoints can be fixed in such a way that synthesis gets not merely one but several, even very different, intermediate terms. If, e.g., $e^{\pm x\sqrt{-1}}$ is developed in a series, one gets, by a separation of the real and the imaginary part: $\left\{1-\frac{x^{2}}{1.2}+\frac{x^{4}}{1.2.3.4}-\ldots\right\}\pm\left\{x-\frac{x^{3}}{1.2.3}+\frac{x^{5}}{1.2.3.4.5}-\ldots\right\}\sqrt{-1}$, which can be denoted by $C(x)\pm S(x)\sqrt{-1}$. If it now shows itself, by appropriate treatment of this expression, that namely: 1) $C^{2}(x)+S^{2}(x)=1$, 2) $S(x\pm y)=S(x)C(y)\pm C(x)S(y)$ and $C(x\pm y)=C(x)C(y)\mp S(x)S(y)$, 3) $C(4pq+x)=C(x)$ and $S(4pq+x)=S(x)$: then we conclude that if $q$„ is found to have the value $\frac{\pi}{2}$, for which $\cos=0$, $\sin=1$, then $C(x)$ and cos x, $S(x)$ and sin (x) run through the same values for the same $x$, and thus are identical.“ $\pi$ is thus here, in that we conclude thus, an intermediate term of synthesis. Cf. A. Steen. Reen Mathematik. Kbhvn 1856. p. 318—26.} namely $\sqrt{-1}$
% --- p. 48 ---
\setcounter{footnote}{0}%
\opage{48}denotes a periodicity, for continued exponentiation gives $\sqrt{-1}:\ \sqrt{-1},-1,-\sqrt{-1},+1,\sqrt{-1}$ etc.

If therefore $e^{x\sqrt{-1}}$ is put instead of $e^{x}$, the following results are obtained:
$e^{x\sqrt{-1}}=\cos x+\sin x\sqrt{-1}$, whereby the exponential function is combined with the trigonometric ones; next $x\sqrt{-1}=l(\cos x+\sin x\sqrt{-1})$, whereby the concept of the logarithm can be developed to the point that $l(+a)=l^{1}a+2p\pi\sqrt{-1}$ proves real for $p=0$, whereas $l(-a)=l^{1}a+(2p+1)\pi\sqrt{-1}$ can never become real for any whole value of $p$ whatever. According to combination, therefore, the transitions from the finite into the infinite become not merely of different kinds, but also systematic, and the systematic transitions countless.

\begin{center}β) Systems of functions: the inner infinity.\end{center}

The mathematical function, which, according to its concept, expresses the dialectical unity of the changeable with the unchangeable, expresses at the same time the transition of the finite into the infinite: as many functions, so many transitions, so many forms of the punctual becoming, so many paths of approximation. These paths do not merely run alongside one another; they cross one another, they wind into one another, they form a net in which the finite, so to speak, catches the infinite. For operation this net is now a system of functions, indeed, of systems of functions; for reflection it is the unfolding of the content of the infinite, the inner infinity. The fundamental logical movements in unity with the exact actions are then: a referring of the many functions under a fundamental formula, a development of the particular functions, insofar as these can each for itself be derived from the common fundamental formula, and finally a
% the note 1) of p. 47 runs over onto the foot of this page: set in full at its mark
% --- p. 49 ---
\setcounter{footnote}{0}%
\opage{49}demonstration of the significance which the punctual increments and remainders acquire in each particular series.

\begin{center}αα) The fundamental formula of the systems of functions: reduction.\end{center}

Just as the many finite cases can be referred under a particular function, so the many particular functions can be referred under a fundamental formula. By this reduction the single and particular is, to be sure, dissolved in the universal, and goes insofar, logically speaking, to the ground. But to go to the ground means, as Schelling and Hegel have shown in a metaphysical direction, a twofold thing: to be annihilated, that is to say: \emph{dissolved}, and to be sublated, i.e. \emph{preserved} in its ground. In order to reduce, one must transform: every reduction is a transformation, but not conversely. As regards the exact development, all possible functions of $x+h$ can be referred to the so-called Taylor series: $f(x+h)=f(x)+f'(x)\frac{h}{1}+f''(x)\frac{h^{2}}{1.2}+\ldots$ This series, which has the peculiarity that the coefficients of the increasing whole and positive powers of $h$ are functions of $x$, shows how a function of $x$ can unfold into a system of functions, when $x$ receives the increment $h$. Insofar as $f(x)$ is not yet more closely determined, the whole system of functions proceeding from $f(x+h)$ is likewise indeterminate; but this indeterminacy is, as lies in the essence of the ground, nevertheless a negative determinacy, a reflection of functions dissolved and preserved in the dissolution. The formula set up is not an immediate first; it is, by virtue of the reduction, a result transformed into an essential presupposition, and as such a fundamental formula.

What a difference it makes whether in cognition one arrives at the universal by abstraction alone, or by reduction, can be seen in the present connection. According to abstraction the universal is empty, the series set up a dead schema, which falls outside the content with which it can be
% --- p. 50 ---
\setcounter{footnote}{0}%
\opage{50}arbitrarily filled; according to reduction the universal is that which is rich in content in the ground, and which confirms its universal validity by the grounding of the particular. That the particular series of functions are reduced to one fundamental formula means, in other words, that the particular laws have their subsistence in one fundamental law.\footnote[1]{From the determinacy of this spherical surface given for intuition: $x^{2}+y^{2}+z^{2}=3^{2}$ one goes by reduction to a determination of the kind of the surface as sphere: $x^{2}+y^{2}+z^{2}=r^{2}$, from this by a new reduction to the sphere as surface of revolution: $z=f(x^{2}+y^{2})$ and from this particular kind to the surface as mathematical genus: $F(x,y,z)=0$. As regards deduction, the inverse procedure, it must be remarked that the mathematical individuals can be deduced, the empirical ones, on the other hand, not.} % print: impiriske (misprint)

\begin{center}ββ) Particular systems of functions: deduction.\end{center}

Only what is contained in the ground can under given conditions be derived from it: deduction presupposes reduction. The condition for the emergence of a particular system from the fundamental formula is an express fixing of the original function. For if one knows the law according to which $f'(x)$, developed in a series by the fundamental formula, arises from the series development for $f(x)$, then one knows at the same time the law according to which $f''(x)$ is derived from $f'(x)$, $f'''(x)$ from $f''(x)$, etc. % print: særeget (misprint)

If, e.g., $f(x)=ax$ is set, the series for $a(x+h)$ becomes $ax+ah$; if $f(x)=x^{a}$ and $f(x+h)=(x+h)^{a}$ is set, one gets the system: $x^{a}+ax^{a-1}\frac{h}{1}+a(a-1)x^{a-2}\frac{h^{2}}{1.2}+\ldots.$; to $f(x)=a^{x}$ and $f(x+h)=a^{x+h}$ corresponds the system: $a^{x}\left(1+la\frac{h}{1}+(l.a)^{2}\frac{h^{2}}{1.2}+\ldots\right)$; to $f(x)=l.x$, and $f(x+h)=l.(x+h)$ the system $lx+\frac{1}{1}\frac{h}{x}\div\frac{1}{2}\frac{h^{2}}{x^{2}}+\ldots$ etc. In relation to reduction possibility is ground, in relation to deduction the ground is possibility. The identity of ground and possibility is seen from the different
% --- p. 51 ---
\setcounter{footnote}{0}%
\opage{51}forms of the disjunctive judgment: $A$ is neither $B$ nor $C$ nor $D\ldots,$ $f(x)$ is neither $ax$ nor $x^{a}$ nor $a^{x}\ldots;$ the particular determinations are dissolved in the universal, and thereby gone to the ground; $A$ is, in the ground, both $B$ and $C$ and $D\ldots,$ $f(x)$ is, in possibility, both $ax, x^{a}, a^{x}\ldots$ Insofar as the particular in the content of the ground lets itself appear, it is reflected in possibility; where one case is possible, at least two cases must be possible, for if the opposite of the possible case were not also possible, the case itself would be necessary: possibility has alternating terms. If now one of the terms is to be posited as existing, the disjunctive judgment in its usual form shows: $A$ is either $B$ or $C$ or $D\ldots,$ $f(x)$ is either $ax$ or $x^{a}$ or $a^{x}\ldots$ the tension between the alternating terms. How the existing is derived by an express positing from the ground, the actual from possibility, the particular from the concrete universality, can be seen from the double \textit{modus} of the disjunctive inference: $f(x)$ is either $ax$ or $x^{a}$ or $a^{x}\ldots,$ now $f(x(=ax,$ thus neither $x^{a}$ nor % print: f(x(=ax (misprint)
$a^{x}\ldots,$ (\textit{modus ponendo tollens}), and: $f(x)$ is either $ax$ or
$x^{a}$ or $a^{x};$ now $f(x)$ is neither $x^{a}$ nor $a^{x},$ thus it must, presupposing that there are here only three possibilities, or rather that the possibility at issue has only three alternating terms, necessarily be $ax$ (\textit{modus tollendo
ponens}). That the functions: $ax, x^{a}, a^{x}\ldots\ldots$ exclude one another (\textit{principium exclusi medii}) has, according to deduction, its significance in this, that they in the ground (\textit{principium inclusi medii}) presuppose one another. The same thing that makes $f(x+h)=(x+h)^{a}$ give the series: $x^{a}+ax^{a-1}\frac{h}{1}+\ldots,$ makes also $a(x+h)$ give $ax+ah,$ and $a^{x+h}$ give $a^{x}+a^{x}\,l.\,a\,\frac{h}{1}+\ldots$ etc., for it is precisely the fundamental law itself, the fundamental thought of the law.

The particular systems are expressions of particular laws. If now reduction shows that the particular laws must have their
% --- p. 52 ---
\setcounter{footnote}{0}%
\opage{52}essential presupposition in the one fundamental law, then deduction shows conversely that the one fundamental law has its subsistence, its finite existence, in the particular laws.

\begin{center}γγ)

Increments and remainders: \emph{Articulation.}\end{center}

As that in which finitude is dissolved, the infinite is, essentially viewed, the ground and possibility of the finite; a carried-through development of the forms of finitude is then at the same time a development of the content of infinity. The demand is that the development become punctual. If one casts punctuality aside, and holds one-sidedly to the universal, one easily has done with the difference between the infinite in representation and in the concept, or, to use Spinoza's expression, between \textit{infinitum imaginationis} and \textit{infinitum actu;} but if the demands of punctuality are acknowledged, it becomes clear that the terms of representation belong to the concept itself.

As regards the quotient $\frac{1}{1-1}$, the series $1+1+1+\ldots$ is necessary to a determination of the infinite; as regards the series, the quotient is necessary. If we hold to the series $1+1+1+\ldots$ alone, then we see only the one side of the infinite, namely the unattainable, that whose finite presentation can never be completed, \textit{infinitum imaginationis}.
If we hold exclusively to the notation $\frac{1}{1-1}$
$=\frac{1}{0}=\infty,$ without considering how $\infty$ first arises by an infinite increase of a finite quotient, in that $0$ arises by an infinite diminution of a finite divisor, we get formally, to be sure, an \textit{infinitum actu,} a concept of infinity, but a dead concept. To insight into the essence of infinity, the inner infinity, it thus belongs that infinitum imaginationis and \textit{infinitum actu} must be posited in dialectical unity, that a series like: $1+1+1+\ldots$ and a quotient like: $\frac{1}{1-1}$ must reflect themselves in one another.
% --- p. 53 ---
\setcounter{footnote}{0}%
\opage{53}Now if the representation is incommensurable with the concept, the articulation of the content of infinity must become an articulation of the infinite approximation; such an articulation is conditioned on the demonstration of the different significance of the increments in the different systems of functions, and on the exact specification of the convergence or divergence of infinite series by the determination of remainders and limits.

That a finite increment $(h),$ by which $x$ in $f(x)$ is increased, must entail an increment of the function itself corresponding to the nature of the function, and that this increment $(k)$ must therefore be different for the different functions, is a matter of course. If one puts $f(x)=y,$ and $f(x+h)=y+k,$ one sees that, whereas $k,$ for $f(x+h)=a(x+h),$ becomes $ah,$ and makes up only one term, and for $f(x+h)=(x+h)^{a},$ for $a$ a positive integer, becomes $ax^{a-1}\frac{h}{1}+a(a-1)x^{a-2}\frac{h^{2}}{1.2}+\ldots h^{a},$ and yields a plurality of terms, namely $a$ terms, $k$ must, e.g., for $f(x+h)=a^{x+h}$ become $a^{x}+a^{x}(la)\frac{h}{1}+a^{x}(la)^{2}\frac{h^{2}}{1.2}+\ldots$ and develop into countless terms. But if one now looks at the series with the countless terms, that is, at the infinite series, then an exact development of their values seems at first glance impossible: if an infinite series is to be unfolded, it must surely be carried to its end, but that is impossible; if the infinite series can never be carried to its end, then a remainder must surely be left over that is not unfolded; but if a remainder is left over that is not unfolded, then the approximation, so it seems, must lose itself in the approximate.

In the series $1+x+x^{2}+\ldots.\ x^{n-1}$ corresponding to $\frac{1}{1-x}$ one may of course gladly add the remainder $\frac{x^{n}}{1-x};$ yet, why stop at $x^{n-1}?$ one might just as well go to $x^{2n-1}$ and then add the remainder $\frac{x^{2n}}{1-x};$ but, when the breaking off
% --- p. 54 ---
\setcounter{footnote}{0}%
\opage{54}is so arbitrary, how can the result then rise above the approximate? The remainder that is thus specified, whether it be given as $\frac{x^{n}}{1-x}$ or as $\frac{x^{n+1}}{1-x}$ or as $\frac{x^{2n}}{1-x}$ contains within itself a something that is not yet developed. If what is contained in the remainder needs no development, then the quotient $\frac{1}{1-x}$ presumably needs no development either, and the whole series becomes superfluous; if the remainder, however far one goes, constantly needs a development, then there remains here a something that can never be developed. If one wants to determine the value of the remainder by noting the difference between such infinite series in which each following term is greater than its predecessor, and such in which the reverse takes place: then this procedure, too, does not seem able to lead in every case to an exact result. If one puts, e.g., $x=2$ in $\frac{1}{1-x},$ so that the series becomes $1+2+2^{2}+2^{3}+\ldots.;$ and next puts $x=\frac{1}{2},$ so that the series becomes $1+\frac{1}{2}+\left(\frac{1}{2}\right)^{2}+\left(\frac{1}{2}\right)^{3}+\ldots,$ then one easily sees that the remainder in the first of these series will stand in a quite different relation to the sum of the terms than the remainder in the second. In the first series the remainder, however long one continues the division, will retain an infinitely great value, whereas in the last series it will decrease. For the series $1+\frac{1}{2}+\left(\frac{1}{2}\right)^{2}+\left(\frac{1}{2}\right)^{3}+\ldots.$ the remainder, however, is not determined by our continuing the division, but by our being informed by another way of the value corresponding to the whole series. How little one can rely upon a steady diminution of the remainder, because one has a series with steadily decreasing terms, can be seen from another case, in which we likewise are informed by a special way of
% --- p. 55 ---
\setcounter{footnote}{0}%
\opage{55}the true value of the series. For since $a^{-\infty}=\left(\frac{1}{a}\right)^{\infty}=0,$ it follows that $\mp\infty=\mathrm{Log}\,0;$ but $\mathrm{Log}\,0=\mathrm{Log}\,(1-1)=-\mathrm{Log}\,e\left(1+\frac{1}{2}+\frac{1}{3}+\ldots.\frac{1}{n}+\ldots\right),$ and, in the natural % print: Da nemlig (misprint)
system, $l(0)=-\left(1+\frac{1}{2}+\frac{1}{3}+\ldots.\frac{1}{n}\ldots.\right)=-\infty;$ from this it is seen that not every series with decreasing terms gives a steadily decreasing remainder. If to this we add that there are, after all, manifold series whose remainders, as, e.g., the remainder for the logarithmic series cited, or as the remainder for the series corresponding to $e$: $1+\frac{1}{1}+\frac{1}{1.2}+\ldots$ cannot even be presented in finite form: then it must surely appear to the reflection hovering over punctuality, which allowed approximateness a questionable latitude everywhere: by the tension between the approximate in the approximation and the exact, i.e.\ by the contradiction hidden in the ``articulation,'' the problem of infinity is characterized.

To what degree the problem of infinity has been able to astonish even the mathematician familiar with analysis, when he began in earnest to want to think it through, is shown in quite a striking way in the following guise, taken from \textit{les Elémens de Mathématique du P. Lami}: \textit{``Le
mauvais Riche brûlé de soif, prie, dit on, Abraham de
lui laisser distiller une goute d'eau. Abraham obtient
ce qu'il demande. Il laisse tomber une goute d'eau,
qui fait 100 lieuës à la premiere minute, 99 à la seconde,
ainsi de suite, toujours en même proportion à
chaque minute de 100 à 99. Or la distance des lieux
où sont Abraham \& le mauvais Riche étant supposée
infinie, on demande en combien de tems cette goute
arrivera au mauvais Riche? On répond à cela, jamais.
En effet, l'espace étant infini, la goute d'eau ne peut
pas parvenir à un terme. On conçoit bien que la chose
ne peut être autrement. Mais ce qui a droit de sur}%
% italic French quotation continues on next page (set as \textit there)
% --- p. 56 ---
\setcounter{footnote}{0}%
\opage{56}\textit{prendre, c'est qu'on détermine le nombre des lieues que
feroit la goute en tombant pendant une éternité; \& ce
nombre est 100,000 lieues. Voila donc une quantité
finie déterminée, quoiqu'il n'y ait point de terme auquel
elle soit limitée, le tems de la chûte d'eau étant infini.
Cela est assurement inconcevable. Afin de rendre la
chose sensible, on fait observer, que la chûte de la
goute d'eau est rétardée, suivant l'hypothese, à chaque
minute. Ce rétardement augmentant toujours, il devient
à la fin si peu considérable, qu'on peut le regarder
comme nul. Alors la goute d'eau ne se meut pas;
puisqu'on néglige son mouvement, qui est infiniment petit.
La question est de savoir, si on peut négliger cet Infiniment-Petit, \& si dans un tems infini, ce rétardement,
quoique toûjours plus petit, n'a pas une valeur réelle.
Notre esprit se perd dans ce raisonnement. Nous concevons
que la progression decroissant toûjours, la valeur
des termes doit diminuer tellement qu'elle ne soit pas
sensible, \& cependant nous voulons que le nombre des
termes, qui composent cette progression, soit infini.
Ceci est tout-à-fait contradictoire.''}
\textit{Application de la
Géométrie ordinaire et des Calculs différentiel et intégral
à la résolution de plusieurs Problêmes. Par feu
M. Robillard. Paris MDCCLIII. (Histoire critique du Calcul
des Infiniment-Petits contenant la Métaphysique et la
Théorie de ce Calcul. P. XVII--XVIII.)}

\begin{center}III.\end{center}

\begin{center}The Analysis of the Infinite.\end{center}

That an infinite approximation can nevertheless end, the unattainable be attained, and the approximate dissolve into the exact, is shown by the analysis of the infinite. In order to make the problem of infinity recognizable, reflection must lean upon
% --- p. 57 ---
\setcounter{footnote}{0}%
\opage{57}the operation; in order to solve the problem, it must appropriate the intelligence expressed in the operation. That between infinite magnitudes there can be a finite relation, that a sum of infinitely many terms can in some cases make up a finite, in others an infinite magnitude, that a series whose terms are so related that the following one divided by the preceding one gives the quotient $1,$ can be either convergent or divergent, etc., such things will altogether escape the thinking that does not, at the hint of the operation, catch sight of the exact determinations. That these determinations in philosophical form must become dialectical, whereas in mathematical form they neither are nor can be so, shows a characteristic difference between mathematics and philosophy, but by no means any conflict. No concept can, however exactly it is defined, maintain itself in a being indifferent to the other concepts: the validity of concepts is seen in their transitions; but the transitions that the exact analysis makes precisely through the operation, the thinking that penetrates the operational concepts makes through dialectic.

The exact determinateness in which the finite touches the infinite: \emph{the limit of the function}; the exact expression of punctuality by which the finite is dissolved into the infinite as into its original moment: \emph{the differential}; the exact form under which the elements of the infinite discretion are united into a continuum completed within possible limits, so that the finite can thereby be restored from the infinite: \emph{the integral,} are then the fundamental concepts that dialectic in this connection has to elucidate.

\begin{center}a) The Limit of the Function: the Dialectical Contact of the Finite with the Infinite.\end{center}

The limit to which a function approaches, as its independent variable either grows or decreases to infinity, is by reflection posited at once both as becoming and as being, and may moreover be now $0,$ now $\infty,$ now finite. If one thus denotes the limit of $y$ in $f(x)=y,$ when $x=\infty,$ by
% --- p. 58 ---
\setcounter{footnote}{0}%
\opage{58}$\mathrm{Lim}\,f(x)=f(\infty),$ and, when $x=0,$ by $\mathrm{lim}\,f(x)=f(0);$ then one sees that here to $f(x)=ax=y$ there must correspond $\mathrm{Lim}\,f(x)=\infty$ and $\mathrm{lim}\,f(x)=0,$ to $f(x)=\frac{a}{x}=y,$
$\mathrm{Lim}\,f(x)=0$ and $\mathrm{lim}\,f(x)=\infty,$ to $f(x)=\frac{x}{x}=y$
on the other hand $\mathrm{Lim}\,f(x)=\mathrm{lim}\,f(x)=1.$ How it comes about that the exact action cannot take notice of the dialectical, notwithstanding that the thoroughgoing thinking must of necessity discover the dialectical, can be shown in the present connection. If one has, e.g., $\mathrm{Lim}\,f(x)=\mathrm{Lim}\frac{a_{0}+a_{1}x+a_{2}x^{2}\ldots+a_{n}x^{n}}{b_{0}+b_{1}x+b_{2}x^{2}\ldots+b_{n}x^{n}},$ and puts without further ado $x=\infty,$ one gets in $\frac{a_{0}+\infty+\infty\ldots}{b_{0}+\infty+\infty\ldots}=\frac{\infty}{\infty}$ a wholly indeterminate result; likewise, if in $\mathrm{lim}\,f(x)=\mathrm{lim}\frac{a_{1}x+a_{2}x^{2}\ldots+a_{n}x^{n}}{b_{1}x+b_{2}x^{2}\ldots+b_{n}x^{n}}$ one sets $x=0$ without any preparation, the result $\frac{0+0+\ldots}{0+0+\ldots}=\frac{0}{0}$ will remain indeterminate.

How wholly determinateness is dissolved, how hovering the limit is, is seen in the cases cited from the fact that the results obtained, though different from one another, can be transformed into a deceptive mutual likeness, namely $\mathrm{Lim}\,f(x)=\frac{\infty}{\infty}=\frac{1}{\infty}\infty.=0\infty,$ and $\mathrm{lim}\,f(x)=\frac{0}{0}=0.\frac{1}{0}=0.\infty.$ By a preparatory action, a determinate reduction, one comes, on the other hand, to a wholly determinate limit.
\[
\mathrm{Lim}\,f(x)=\mathrm{Lim}\frac{\frac{a_{0}}{x^{n}}+\frac{a_{1}}{x^{n-1}}\ldots+a_{n}}{\frac{b_{0}}{x^{n}}+\frac{b_{1}}{x^{n-1}}\ldots+b_{n}}=\frac{a_{n}}{b_{n}}=c,
\]
and
% --- p. 59 ---
\setcounter{footnote}{0}%
\opage{59}\[
\mathrm{lim}\,f(x)=\frac{a_{1}+a_{2}x+\ldots a_{n}x^{n-1}}{b_{1}+b_{2}x+\ldots b_{n}x^{n-1}}=\frac{a_{1}}{b_{1}}=d,
\]
which cannot be transformed into $c.$ If one puts, to cite yet another example, in $\mathrm{Lim}\left(\frac{1}{x^{x}}\right)^{\frac{1}{x}}$ without further ado $x=\infty,$ one gets only $\left(\frac{1}{\infty}\right)^{\frac{1}{\infty}}=0^{0},$ which in this case is indeterminate; if one, on the other hand, first undertakes the requisite reduction, one obtains $\mathrm{Lim}\left(\frac{1}{x^{x}}\right)^{\frac{1}{x}}=\mathrm{Lim}\frac{1}{x}=0,$ which is a determinate limit.

The exact actions thus follow one another in a certain order: % print: hianden (misprint)
 they are, so to speak, intellectual undertakings carried out in time: first the one happens, then the other; first the one way of viewing is applied, then the other; first one treats $x$ as a finite magnitude, afterwards it is posited as $\infty$ or as $0;$ first one assumes the limits $\infty$ and $0$ to be unattainable, afterwards, when by proper reduction of the still finite terms one has made the law of the origin recognizable, one regards the matter as if the limit that is unattainable for $x$ in $f(x)=y$ had nevertheless been attained after all, and then sees what influence this must have on $y.$ It is this separation of the exact actions in time that lets the opposed ways of viewing appear beside one another, alternate with one another, fall apart from one another, and thus prevents the dialectical from coming to light.

In the thoroughgoing reflection the separation in time falls away; the finite partitioning vanishes in the unity of the concept, the opposed determinations shine through one another, and the dialectical shows itself. The equations: $\mathrm{Lim}\,f(x)=f(\infty)=c,$ $\mathrm{lim}\,f(x)=f(0)=d$ express, when they are grasped in the concept as logical judgments, that the unattainable is attainable. This contradiction, which, torn loose from its presupposition, is
% --- p. 60 ---
\setcounter{footnote}{0}%
\opage{60}just as intolerable as the contradiction that the square is round, has in its essential connection the significance that it compresses two opposed, mutually exclusive determinations into one fundamental determination, not in order to make them straightforwardly into one, but to show their common origin. That the limit for $x,$ and consequently also the limit for $f(x)$ in $\mathrm{Lim}\,f(x)=f(\infty)=c$ is unattainable, is of course precisely grounded in $f(x),$ which demands that the countless intermediate terms by which $x$ as a finite magnitude differs from $\infty$ be run through; but the equation itself, in which there is not even a shadow of the approximate, indicates by the exact value $c$ at the same time the attained, thus after all attainable limit. The explanation that the limit is reached by a leap, in that one has leapt over countless intermediate terms, is one-sided and misleading; in equations such as: $\mathrm{Lim}\,f(x)=c$ and $\mathrm{lim}\,f(x)=d,$ no intermediate term is, essentially seen, leapt over; the same law that in $\mathrm{Lim}\,f(x)=c$ holds for $x=\infty,$ holds also for $x$ as finite, and the law that in $\mathrm{lim}\,f(x)=d$ holds for $x$ as finite, holds in consequence also for $x=0.$ The proposition: the unattainable is attainable, has its ground in the essence of the function; when this grounding is seen into, the contradiction vanishes, for the contradictorily opposed determinations are united only in order to be separated again, and in finite separation to exclude one another.

In what sense the unattainable excludes attainment can be seen from the breaking off of the infinite series; in what sense the unattainable nevertheless becomes attained can be seen from the essential closure of these series; how intimately the opposed determinations of limitation interpenetrate one another can finally be seen in the infinite magnitudes.

\begin{center}α)

The Breaking Off of Infinite Series: the Limit as Unattainable.\end{center}

Since an infinite series cannot be carried to its end, it must be broken off; the remainder that thereby arises is an
% --- p. 61 ---
\setcounter{footnote}{0}%
\opage{61}unknown magnitude. Here, then, we meet the difficulty touched on above, that in order to know the complete sum one would have to know the remainder beforehand, and conversely, in order to know the remainder, one would first have to know the complete sum. This difficulty is removed when the relation of the terms under a finite continuation shows to what limit the infinite continuation would lead. It is the difference between convergent and divergent series that here furnishes the determinate point of departure: ``if $s_{n},$ in the sum $s_{n}=u_{0}+u_{1}\ldots+u_{n},$ approaches, for $n$ constantly growing, a determinate finite limit, the series itself is said to be convergent; if a series is not convergent, it is said to be divergent.'' The exact conception is now briefly as follows:

``In every convergent series $\mathrm{Lim}\,u_{n}=0,$ but this is only a necessary, not a sufficient condition for convergence, for an infinite number of positive terms converging to $0$ can give a sum $=\infty.$ On the other hand, when $r_{n}$ denotes the sum of all the terms $u_{n+1},u_{n+2},u_{n+3},\ldots.$ following $u_{n}$ to infinity, which is called the remainder of the series, then $\mathrm{Lim}\,r_{n}=0$ will express both a necessary and sufficient condition for the convergence of the series, for one can then always make $n$ so great that $r_{n}$ is smaller than any given magnitude whatever, whereby $s_{n}$ is brought as near as one will to a certain limit.''\footnote[1]{Ramus. Algebra og Functionslære. p. 80.}
If one has now, by applying the expressions found for this, $\mathrm{Lim}\,\frac{u_{n+1}}{u_{n}}$ and $\mathrm{Lim}\,n\left(\frac{u_{n}}{u_{n+1}}\div1\right)$ assured oneself of the convergence of a series, then a demonstration of the convergence can in turn lead to a closer determination of the remainder. As regards the quotient $\frac{1}{1-x}$, the expression $\frac{1}{1-x}=1+x+x^{2}+x^{3}+\ldots.\left\{\begin{array}{l}x=-1\\x=+1\end{array}\right.$ indicates ``that this form
% --- p. 62 ---
\setcounter{footnote}{0}%
\opage{62}satisfies the condition $\mathrm{Lim}\,r_{n}=0,$ when $x$ has a value enclosed between the limits: $\pm1$ exclusive. If now $u_{n}=x^{n}$ gives as remainder $r_{n}=\frac{x^{n+1}}{1-x}=u_{n}\frac{x}{1-x},$ then one sees, ``that as soon as in a series the terms decrease so strongly that $\frac{u_{n+1}}{u_{n}}<x$ for a certain index $n$ and for all higher ones, then also $r_{n}<u_{n}\frac{x}{1-x}.$'' ``This,'' says Ramus, ``can serve to determine the higher limit for the error committed by an infinite series being broken off at $u_{n}$.''

In the exact analysis one thus commits an error by breaking off an infinite series; the index $n$ in $u_{n}$ may be as great as one will: as long as $n$ is finite, the breaking off entails an error. Thus the above series gives for $x=\frac{1}{2},$ $r_{n}=u_{n},$ for $x=\frac{1}{3},$ $r_{n}=\frac{1}{2}u_{n}\ldots.$ for $x=\frac{1}{p},$ $r_{n}=\frac{1}{p-1}u_{n}.$ If $\frac{u_{n+1}}{u_{n}}$ is not constant but ``for a certain index $n$ and for all higher ones'' steadily decreasing, the remainder will naturally become still smaller. In such cases one can determine the error insofar as one can determine the remainder; but that there always remains a remainder here proves that the breaking off of the series always gives an error. For the series: $e=1+\frac{1}{1}+\frac{1}{1.2}+\ldots,$ e.g., it has now presumably been proved, both as regards the sum, that $2<e<3,$ and that the remainder $r_{n}<\frac{1}{1.2.3\ldots.n.n'},$ when the series is broken off at $u_{n};$ but one has then also at the same time proved the impossibility of avoiding an error, in that one has namely proved that the whole magnitude and consequently also the remainder is incommensurable with the unit. In this point the opposition between the operational concept and the concept of reflection is quite striking. The operational concept, which
% --- p. 63 ---
\setcounter{footnote}{0}%
\opage{63}is designed for calculation, is satisfied exactly by the limits of the error being determined; the concept of reflection, which is designed for logical insight, demands a dialectical satisfaction. Insofar as the exact analysis must confine itself to setting limits for the error without being able to remove the error, it itself concedes the imperfect, gives up its exact character, and contents itself with the approximate; but when the analysis in turn specifies the relation between greater and smaller error, it becomes distinctly conscious of the error, and knows how to set it as narrow limits as is wanted, it has precisely the error in its power, and can thus correct the error\footnote[1]{The series for the Newtonian binomial formula, $(a+x)^{n}=a^{n}+\frac{n}{1}a^{n-1}x+\ldots.,$ which, when n is negative or fractional, becomes infinite, is convergent when $\mathrm{Lim}\frac{A_{r+1}}{A_{r}}x=\mathrm{Lim}\left(\frac{n-r}{r+1}\right)\frac{x}{a}=\frac{x}{a}<1,$ thus when $a>x>-a.$ If the series is used for ``approximate calculation, in that one takes e.g. the sum of p terms, the remainder of the series, or the error that is committed, becomes expressed by: $r_{p}=A_{p}a^{n-p}x^{p}\left\{1+\frac{n-p}{p+1}\frac{x}{a}+\frac{(n-p)(n-p-1)}{(p+1)(p+2)}\left(\frac{x}{a}\right)^{2}+\ldots\right\}<A_{p}a^{n-p}x^{p}\left(1+\frac{x}{a}+\left(\frac{x}{a}\right)^{2}+\ldots\right)=A_{p}a^{n-p+1}\frac{x^{p}}{a-x}\left(=A_{p}a^{n-p}x^{p}\frac{a}{a-x}\right)$ ``in that one merely observes that $n-p<p+1$ or $n-1<2p,$ which can always be achieved by continuing the series far enough.'' A. Steen. Reen Mathm. p. 307--308. % print: bruden (misprint)

``A necessary preparation for the approximate calculation of the real roots in numerical equations is their separation, i.e.\ the determination of the different intervals in which they are singly situated. If $f(x)=0$ is an equation with merely odd roots, the separation can be achieved by substituting in $f(x)$ for x the different real magnitudes situated between the extreme limits, which are terms in an arithmetic series whose difference is smaller than the difference of any two roots in $f(x)=0$.. (Ramus. Algebra og Functionslære. p. 129.) If one has, then, in the equation $f(x)=0$ (I) a root completely separated, e.g.\ between $a_{0}$ and $a_{0}+1$, then the equation must be satisfied by $x=a_{0}+\frac{1}{x_{1}}$ $(\alpha)$. Substituting this value in (I), one gets at least $f(x_{1})=0$ (II); herein there must be a real root $>1$, but also only one, since we of course proceeded from the supposition that between $a_{0}$ and $a_{0}+1$ a root was completely separated: $x_{1}$ can, in other words, have only one value $>1$. If this root is now completely separated, we find it situated between $a_{1}$ and $a_{1}+1$, i.e.\ $a_{1}<x_{1}<a_{1}+1$, so that $x_{1}=a_{1}+\frac{1}{x_{2}}$ $(\beta)$. Substituting $(\beta)$ in $(\alpha)$, one gets: $x_{0}=a_{0}+\cfrac{1}{a_{1}+\cfrac{1}{x_{2}}}$ $(\gamma)$. Substituting the expression $(\beta)$ in (II), one gets $f(x_{2})=0$ (III), which likewise has only one real root, which is assumed to be completely separated between $a_{2}$ and $a_{2}+1$, so that $a_{2}<x_{2}<a_{2}+1$, and $x_{2}=a_{2}+\frac{1}{x_{3}}$. By substituting the last expression in $(\gamma)$ one obtains $x=a_{0}+\cfrac{1}{a_{1}+\cfrac{1}{a_{2}+\cfrac{1}{x_{3}}}}$ etc. Here, then, as one sees, a continued fraction forms, which can be continued to infinity; thus one has $\sqrt{2}=1+\cfrac{1}{2+\cfrac{1}{2+\ldots}}$ If one now breaks off such a continued fraction, and assumes two successive convergents to the magnitude x sought by the approximation to be $\frac{y_{r-1}}{z_{r-1}}$ and $\frac{y_{r}}{z_{r}}$, then one has $\frac{y_{r-1}>\ \ >y_{r}}{z_{r-1}<x<z_{r}}$. One does undeniably commit an error wherever one breaks off the infinite continued fraction, but the error one commits, by e.g.\ taking $\frac{y_{r-1}}{z_{r-1}}$ for x, is in numerical value smaller than $\frac{y_{r}}{z_{r}}-\frac{y_{r-1}}{z_{r-1}}=\frac{1}{z_{r-1}z_{r}}$, and to that extent exactly determined. (Cf.\ Ramus. Elementær Algebra, p. 103--112. and Algebr. Functionsl. p. 25--35).}; the exact
% --- p. 64 ---
\setcounter{footnote}{0}%
\opage{64}analysis finds itself thus, when seen from the standpoint of reflection, in the dialectical case that, by breaking off the infinite series, it commits error without being in error, that it errs without erring.
% the note 1) of p. 63 runs over onto the foot of this page: set in full at its mark
% --- p. 65 ---
\setcounter{footnote}{0}%
\opage{65}
\begin{center}β)

The Essential Closure of the Series: the Limit \emph{as Attainable.}\end{center}

That in the essential closure of an infinite series everything depends on insight and comprehension is as incontrovertible as that the concept of punctuality, which is at issue here, is different from any concept of general logic whatever. Designations such as: $u_{n},$ the nth term, $\mathrm{Lim}\frac{u_{n+1}}{u_{n}},$
$\mathrm{Lim}\,r_{n}$ and the like are, essentially understood, more than agreed-upon expressions of calculation; they are the stamped forms of the act of thought, operational categories, in which the finite is shone through by the infinite in so peculiar a way that the prolixity of the endless is removed, and the series comes to closure. How the nth term in the generality of the concept mirrors the infinite swarm of the individual terms that are represented has already been brought out in the determination of the relation between the operation and the objective reflection. In order to grasp the essential in the closure of an infinite series it will, however, be necessary here to consider the relation between the operational concept and the concept of reflection somewhat more closely.

For judging the convergence of a series there is set up, as is well known, the doctrine that the series is convergent when $\mathrm{Lim}\frac{u_{n+1}}{u_{n}}<1,$ and on the other hand divergent when $\mathrm{Lim}\frac{u_{n+1}}{u_{n}}>1.$ Although one may well possess a certain skill in using the formulas set up, and to that extent have the operational concepts expressed in them, without yet seeing through the essence of the matter, the concept of reflection nevertheless lies so near in the present case that under continued exercise of the operation it must in the end come as of itself. That a series whose terms grow in such a way that $\mathrm{Lim}\frac{u_{n+1}}{u_{n}}>1,$ must
diverge is obvious; nor can the understanding be remote that a series whose terms decrease in such a way that
% --- p. 66 ---
\setcounter{footnote}{0}%
\opage{66}$\mathrm{Lim}\frac{u_{n+1}}{u_{n}}<1,$ must converge. To what degree the operational categories contribute to clarifying the objective reflection can be seen here by juxtaposing such series as: $1+\frac{1}{2}+\left(\frac{1}{2}\right)^{2}+\ldots$ and $1+\frac{1}{2}+\frac{1}{3}+\ldots$ Since the terms in both these series are decreasing, one would at first glance suppose that both must be convergent, the latter just as well as the former. By the formulas of the operation, however, we are informed not merely of the contrary, but we are also led to see the ground. If we consider the general terms $\left(\frac{1}{2}\right)^{n}$ and $\frac{1}{n}$ somewhat more closely, we are convinced that $\left(\frac{1}{2}\right)^{n+1}$ must, when $n$ is exceedingly great, be, in relation to $\left(\frac{1}{2}\right)^{n}$, quite differently vanishing than $\frac{1}{n+1}$ in relation to $\frac{1}{n}.$ In the series $1+\frac{1}{2}+\left(\frac{1}{2}\right)^{2}\ldots+\left(\frac{1}{2}\right)^{n}+\ldots$ the diminution of each following term thus stands in such a relation to the increase of the number of terms that the first cancels the last in the infinite, and the whole sum falls within finite limits. In the series: $1+\frac{1}{2}+\frac{1}{3}\ldots+\frac{1}{n}\ldots$ the diminution of each following term becomes, on the contrary, the more insignificant the more the number of terms grows; however great $n$ becomes, $r_{n}$ will on these terms, by its countless terms, always be able to give a sum $=\infty,$ and the series thus be divergent. It thus becomes evident in this way that a series in which $\mathrm{Lim}\frac{u_{n+1}}{u_{n}}=1,$ can, as in the given example, diverge, whereas a series in which $\mathrm{Lim}\frac{u_{n+1}}{u_{n}}<1$ must, according to the nature of the case,
% --- p. 67 ---
\setcounter{footnote}{0}%
\opage{67}be convergent. Here one can speak, in the strictest sense, of insight and comprehension.

A congruence of reflection with operation so complete does not, on the other hand, take place as regards the expression: $\mathrm{Lim}\ n\left(\frac{u_{n}}{u_{n+1}}-1\right)>1$. A series in which the terms are so related that $\mathrm{Lim}\ \frac{u_{n+1}}{u_{n}}=1$, can, as we have seen, be divergent, but under a given condition it can also be convergent. When the convergence of such a series is to be investigated, one chooses as norm another series whose convergence can be demonstrated specially, independently of general theories, and in which e.g. $\mathrm{Lim}\ \frac{t_{n+1}}{t_{n}}=1$, $t_{n}$ being the general term. Let the mth term of such a series be now: % print: nte Led (misprint)
\[ \frac{1.2.3\ldots m}{a(a+1)..(a+m-1)}\cdot\frac{a-1}{a+m} \]
\[ =\left(\frac{1.2.3\ldots m}{a(a+1)..(a+m-1)}-\frac{1.2.3..(m+1)}{a(a+1)\ldots(a+m)}\right) ; \]
putting herein $m$ successively $=1,2,3\ldots n$, one gets the sum of the series expressed, on the one side by the individual terms, on the other side by a reduction, namely:
\[ \frac{a-1}{a}\left(\frac{1}{a+1}+\frac{1.2}{(a+\cdot 1)(a+2)}\ldots+\frac{1.2.3..n}{(a+1)\ldots(a+n)}+\ldots\right) \]
\[ =\frac{1}{a}-\frac{1.2.3..(n+1)}{a(a+1)\ldots(a+n)} ; \]
by comparing sums
with more and more terms one convinces oneself that the series is
convergent, for $a>1$, and the whole sum a proper fraction.
In this series one then has $\frac{t_{n+1}}{t_{n}}=\frac{n+1}{a+n+1}$, $\frac{t_{n}}{t_{n+1}}=\frac{a}{n+1}+1$, and thus, in conformity with the requirement set up, $\mathrm{Lim}\ \frac{t_{n+1}}{t_{n}}=1$. When now $\frac{u_{n}}{u_{n+1}}$ in a
% --- p. 68 ---
\setcounter{footnote}{0}%
\opage{68}series whose general term is $u_{n}$, is greater than, or equal to $\frac{t_{n}}{t_{n+1}}=\frac{a}{n+1}+1$ (in which $a>1$), then the $u$-series is convergent, since the $t$-series is so. This condition of convergence
is expressed by $\frac{u_{n}}{u_{n+1}}\geq\frac{a}{n+1}+1$; thus:
\[ \mathrm{Lim}\ n\left(\frac{u_{n}}{u_{n+1}}-1\right)+\mathrm{Lim}\left(\frac{u_{n}}{u_{n+1}}-1\right)\geq a, \]
or
\[ \mathrm{Lim}\ n\left(\frac{u_{n}}{u_{n+1}}-1\right)\geq a. \]
Since here as regards $a$ only $a>1$ is required, the expression $\mathrm{Lim}\ n\left(\frac{u_{n}}{u_{n+1}}-1\right)>1$ can in complete
generality specify the condition for convergence in the case asked about. Herewith the operational concept is
essentially closed off; if reflection tries to go further,
a barrier shows itself. From
$\frac{t_{n}}{t_{n+1}}=\frac{a}{n+1}+1$ one can undeniably very easily derive $\mathrm{Lim}\ n\left(\frac{t_{n}}{t_{n+1}}-1\right)$, and when this
expression is found, its application to the $u$-series falls, according to the comparison above, as of itself; but
if one then raises the question how it is actually
grounded in $\frac{t_{n}}{t_{n+1}}$ that $\mathrm{Lim}\ n\left(\frac{t_{n}}{t_{n+1}}-1\right)$ must give
$a>1$, in order that the series can be convergent, the operator, who sees himself able only by a specially practical way to
establish the convergence of the $t$-series, will without doubt confine himself
to pointing to the formula in question, and answer: ``I
can prove it, but I cannot comprehend it.''\footnote[1]{Whether the difficulty here indicated is insurmountable for mathematics itself or not is left undecided in this connection; only this must be expressly emphasized, that the synthetic method just asserts its conclusive independence by the exclusion of the analytic in such cases where the conduct of the proof cannot be raised to development of the concept, and that it to this extent, instead of denoting the more perfect, on the contrary rests on an imperfection.}

% --- p. 69 ---
\setcounter{footnote}{0}%
\opage{69}According to the concept of reflection, the infinite series is essentially closed by the limit of infinity showing itself, at all points, to touch the finite. In equations such as:
\[ \frac{1}{1-x}=\Sigma x^{n}+\frac{x^{n+1}}{1-x} \]
\[ e=1+\Sigma\frac{1}{1.2.3....n}+r_{n}, \]
where $r_{n}<\frac{1}{1.2.3...nn}$, etc., it is a matter of indifference, with respect to the essence of the matter, whether two terms or myriads of terms are put down here in the sums. The limit of infinity is, for all that, equally remote and equally near, equally unattainable, equally attainable. If the matter were not so, an infinite series could indeed be broken off, but never closed. If one sets, e.g., $n=\infty$ in the series $1+x+x^{2}+..+x^{n}+\frac{x^{n+1}}{1-x}$, then surely not only $x^{n+1}$ and $x^{n}$ but likewise $x^{n-1}, x^{n-2}, x^{n-3}\ldots$ must be transformed into $x^{\infty}$; the infinite would then in this way exert so monstrous a retroaction that the series, with myriads of different terms, would in the end have only terms of a single form, namely $x^{\infty-\infty}$; and what then is $x^{\infty-\infty}$?
From $\frac{1}{1-x}$ one obtains, for $x=\frac{3}{2}$, the series: $1+\frac{3}{2}+\left(\frac{3}{2}\right)^{2}....$
\[ +\left(\frac{3}{2}\right)^{n}+\frac{\left(\frac{3}{2}\right)^{n+1}}{1-\frac{3}{2}} ; \]
if now $n=\infty$ is set, then surely $\left(\frac{3}{2}\right)^{n}$
\[ =\infty, \left(\frac{3}{2}\right)^{n-1}=\infty, \left(\frac{3}{2}\right)^{n-2}=\infty \]
etc., the whole sum up to and including $\left(\frac{3}{2}\right)^{n}$ $=\infty$, the remainder $\frac{\left(\frac{3}{2}\right)^{n+1}}{1-\frac{3}{2}}=-\infty$; what
% the note 1) of p. 68 runs over onto the foot of this page: set in full at its mark
% --- p. 70 ---
\setcounter{footnote}{0}%
\opage{70}does this signify? From $\frac{1}{1-x}$ one obtains, for $x=\frac{1}{2}$, the series $1+\frac{1}{2}$
\[ +\left(\frac{1}{2}\right)^{2}\ldots+\left(\frac{1}{2}\right)^{n}+\frac{\left(\frac{1}{2}\right)^{n+1}}{1-\frac{1}{2}} ; \]
if $n=\infty$ is set in this, then not only $\frac{\left(\frac{1}{2}\right)^{n+1}}{1-\frac{1}{2}}=0$ and $\left(\frac{1}{2}\right)^{n}=0$, but $\left(\frac{1}{2}\right)^{n-1}$
\[ \left(\frac{1}{2}\right)^{n-2}, \left(\frac{1}{2}\right)^{n-3}.... \]
$=0$ to infinity; what then is $0$, what is $\infty$?

\begin{center}γ) \emph{Infinite Magnitudes: The Dialectic of Limitation.}\end{center}

By a sharp separation of the magnitude itself and its limit, the designations $0$ and $\infty$ can be retained, not as magnitudes but as limits of magnitude, in such a way that $0$ excludes and $\infty$ swallows up every possible magnitude. In an expression such as
\[ \frac{a}{0}=\infty \]
$0$ then signifies, not that there is, but that there has been a divisor here, which has now vanished; nor can $\infty$ be regarded as a proper quotient, but as a sign of a quotient dissolved into the infinite, a quotient that ceased to be a magnitude just at the moment when the corresponding divisor dwindled into nothing. But in that $0$ and $\infty$ are conceived in this way as sharp negations, as pure limits, the concept of limit itself is conceived one-sidedly; for the limit is, according to its concept, not merely negative but also positive: the negative is a limit by the positive's being so. For it is not merely $0$ and $\infty$ that are limits; magnitudes such as $+1$ and $-1$ are surely limits too, the former negative, the latter positive and negative. A magnitude that falls between the limits $0$ and $\infty$ will therefore, according to the stated point of view, be annihilated when it has reached % print: angivne .Synsmaade (a full stop) (misprint)
% --- p. 71 ---
\setcounter{footnote}{0}%
\opage{71}one of these limits; but a magnitude that falls between the limits $+1$ and $-1$ will, precisely by attaining one of these limits, show itself as a determinately existing magnitude, for $-1$ is surely just as much a magnitude as $+1$. Between the limits $+1$ and $-1$, $0$ is the limit at which $+$ and $-$ touch one another; to hold fast the separation of two kinds of limits, negative and positive, is therefore impossible, since it is just the positive which, in the limit itself, borders on the negative. But if it is impossible to separate what is negative in the limit from its positive, then it is also impossible to hold the limit of the magnitude finally apart from the magnitude itself.

In the cognition that the limit of a magnitude necessarily belongs to its determination as a magnitude, one could nevertheless, without yet conceiving $0$ and $\infty$ as magnitudes, try to go a little step further, and conceive them as transitions of magnitudes.\footnote[1]{``Since no number raised to the 0th power produces a, $\sqrt[0]{a}$ cannot be given by any previously known magnitude. The root cannot be rational, because every power of the $0$th degree is $1$, and it cannot be irrational either, because no number can be comprised between two successive powers of $n$th degree. But $\sqrt[0]{a}$ can be determined by applying the equation $\sqrt[q]{a^{p}}=a^{\frac{p}{q}}$, setting q=0, p=1, which gives $\sqrt[0]{a}=a^{\frac{1}{0}}=a^{\infty}$, which is either $0$ or $\infty$, according as $a<1$ or $a>1$. One sees from this that $0$ and $\infty$ belong neither to the rational nor to the irrational magnitudes; moreover they are in fact not magnitudes at all, but merely indicate the cessation of the magnitude, its forms of transition.'' A. Steen. Reen Mathem. Kbhvn 1856. p. 115--116.} Since what is dealt with here are different functions of $0$, $\sin 0$, $\cos 0$, $\mathrm{tg}\,0$ etc., the meaning is not that $0$ in itself is something, nor that $0$ is plain nothing, for $0$ is a transition; but if $0$ in $\mathrm{tg}\,0$ is a transition, then $\infty$ in $\cot 0=\infty$ is surely a transition too.

% --- p. 72 ---
\setcounter{footnote}{0}%
\opage{72}Now since a transition, as regards the opposition of positive and negative, acquires significance only through the positive's passing over into the negative, and conversely the negative into the positive, hence through the positive and the negative both being in the transition, and for that reason belonging to the concept of transition: such equations as $\sin 0=0$, $\cot 0=\infty$ ought, strictly speaking, to be written: $\sin(\pm 0)=\pm 0$, $\cot(\pm 0)=\pm\infty$, etc. By this separation of $+$ and $-$ in $0$ and $\infty$ there arises in these designations of limit themselves an inner difference, which is in the highest degree dialectical. For the designation $\pm 0$ expresses, when we look more closely, that the one zero has divided into three different zeros, a $+0$, a $-0$, and the pure, neutral $0$, the sharp negation of $+$ and $\div$; the same separation repeats itself as regards $\pm\infty$; here too there is a $+\infty$, a $-\infty$, and a neutral $\infty$, in which $+$ and $-$ have vanished. When it is therefore said that $0$ and $\infty$ are designations of transition, this cannot hold of the pure $0$ or the neutral $\infty$, but only of $\pm 0$ and $\pm\infty$

The way in which the so-called indeterminate forms:
\[ \frac{0}{0}, 0.\infty, \frac{\infty}{\infty}, 0^{0}, \infty^{0} \]
and $1^{\infty}$ are treated in mathematics shows clearly enough what forms of transition mean. These forms, which, being limits of functions, are thus not immediate limits but determinations of limit reflected into themselves, may be thought to arise through $x$ in $f(x)$ being set = $a$, and then determined in accordance with the fundamental equation $f(a)=\lim f(a+h)$; they divide essentially into two groups, of which the first: $\frac{0}{0}, 0.\infty, \frac{\infty}{\infty}$ and $\infty-\infty$ falls under the formula:
\[ \frac{\Phi(a)}{\Psi(a)}=\frac{0}{0} \]
on the condition that $\Phi(a)=0$ and $\Psi(a)=0$;
the second: $0^{0}, \infty^{0}$ and $1^{\infty}$ falls, on the same condition, under the formula $\Phi(a)^{\Psi(a)}$.

% --- p. 73 ---
\setcounter{footnote}{0}%
\opage{73}How the forms belonging under the first group can be determined is seen from the following series:
\[ \frac{\Phi(a+h)}{\Psi(a+h)} \]
\[ =\frac{A_{0}h^{\alpha_{0}}+A_{1}h^{\alpha_{1}}+..}{B_{0}h^{\beta_{0}}+B_{1}h^{\beta_{1}}+..}=\frac{A_{0}}{B_{0}}h^{\alpha_{0}-\beta_{0}}\frac{1+\frac{A_{1}}{A_{0}}h^{\alpha_{1}-\alpha_{0}}+\ldots}{1+\frac{B_{1}}{B_{0}}h^{\beta_{1}-\beta_{0}}+\ldots}, \]
where the terms, be it well noted, are ordered by increasing powers of $h$, so that $\alpha_{0}$ is the smallest exponent in the upper series, $\beta_{0}$ the smallest exponent in the lower series. According as $\alpha_{0}-\beta_{0}\gtreqless 0$, one obtains
\[ \frac{\Phi(a)}{\Psi(a)}=\frac{0}{0}=\begin{cases}0\\ \frac{A_{0}}{B_{0}}\\ \infty\end{cases}. \]

As regards the second group, what essentially matters is to determine $\lim\Psi(a+h)\,l.\,\Phi(a+h)$, since
\[ \Phi(a)^{\Psi(a)}=\lim\Phi(a+h)^{\Psi(a+h)}=e^{\lim\Psi(a+h)\,l\,\Phi(a+h)}. \]
If this exponent is developed
in the series:
\[ (B_{0}h^{\beta_{0}}+B_{1}h^{\beta_{1}}+\ldots)\,l.A_{0}h^{\alpha_{0}}+A_{1}h^{\alpha_{1}}+\ldots) \]
, it will be seen, after
due transformation, that both $0^{0}$ and $\infty^{0}$ always
% NOTE: no mark printed; the note stands at the foot of the page
\footnote[1]{The forms $\infty^{0}$ and $0^{0}$ are namely determined thus: $\Phi(x)^{\Psi(x)}=e^{\Psi(x)\,l\Phi(x)}$, for $e^{l.\Phi(x)}=\Phi(x)$, $\Phi(a)^{\Psi(a)}=\lim\Phi(a+h)^{\Psi(a+h)}=\lim e^{\Psi(a+h)\,l(\Phi(a+h)}=e^{\lim\Psi(a+h)\,l.\Phi(a+h)}$; it thus comes down to determining $\Psi(a+h)\,l\,\Phi(a+h)$. $\Psi(a+h)\,l\,\Phi(a+h)=(B_{0}h^{\beta_{0}}+B_{1}h^{\beta_{1}}+\ldots..)\,l.(A_{0}h^{\alpha_{0}}+A_{1}h^{\alpha_{1}}+\ldots..)=(B_{0}h^{\beta_{0}}+B_{1}h^{\beta_{1}}+\ldots).(l.A_{0}h^{\alpha_{0}}+l(1+\frac{A_{1}}{A_{0}}h^{\alpha_{1}-\alpha_{0}}+\ldots..)=(B_{0}h^{\beta_{0}}+B_{1}h^{\beta_{1}}+..)l.A_{0}\,(\mathrm{I})+\alpha_{0}(B_{0}h^{\beta_{0}}+B_{1}h^{\beta_{1}}+\ldots)l.h\,(\mathrm{II})+(B_{0}h^{\beta_{0}}+B_{1}h^{\beta_{1}}+\ldots)\left(\frac{\frac{A_{1}}{A_{0}}h^{\alpha_{1}-\alpha_{0}}+..}{1}-\ldots\right)(\mathrm{III})$. Since $\Psi(a)=0$, one sees that all the $\beta$ are positive, consequently lim (I) $=0$; since $\lim x^{n}\,l\,x=0$, lim (II) $=0$; since finally $\beta_{0}$ is positive and $\alpha_{1}-\alpha_{0}$ positive, the series being ordered by increasing powers (thus: $\alpha_{1}>\alpha_{0}$), (III) will contain nothing but positive exponents of h, thus: lim (III) $=0$. Thus one has: $\Phi(a)^{\Psi(a)}=e^{\lim\Psi(a+h)\,l(\Phi(a+h))}=e^{0}=1$. Example: $f(x)=(1+x)^{1+x^{2}}$ gives, for $x=-1$, $0^{0}$; thus: $f(-1+h)=e^{(2h-h^{2})l.h}=e^{2h\,l\,h-h^{2}l.h}$, and $f(-1)=e^{0}=1$. It might thus seem as if $0^{0}$ and $\infty^{0}$ were not indeterminate forms at all; nor would they be, if $\Phi(a+h)$ and $\Psi(a+h)$ could always be developed in series by increasing powers of h. But there are, as is well known, some cases where a series development is impossible, and where one can therefore, by special determinations, obtain for $0^{0}$ or for $\infty^{0}$ a value other than 1. Thus $\left(\frac{1}{(\frac{1}{x})^{x}}\right)^{\frac{1}{x}}$, for $x=\infty$, gives $\infty^{0}$; but Lim $\left(\frac{1}{(\frac{1}{x})^{x}}\right)^{\frac{1}{x}}=\mathrm{Lim}\frac{1}{((\frac{1}{x})^{x})^{\frac{1}{x}}}=\frac{1}{\mathrm{Lim}\frac{1}{x}}=\infty$.} % print: l(Φ(a+h) with one closing parenthesis (misprint)
% --- p. 74 ---
\setcounter{footnote}{0}%
\opage{74}have the determinate value 1, provided only that, according to the presupposition, the above series development of $\Phi(a+h)$ and $\Psi(a+h)$ is possible, and that $1^{\infty}$ must, according as $\alpha_{1}+\beta_{0}\gtreqless 0$, give
\[ \Psi(a)\,l\,\Phi(a)=\begin{cases}0\\ A_{1}B_{0}\\ \infty\end{cases} ; \]
thus:
\[ \varphi(a)^{\Psi(a)}=1^{\infty}=\begin{cases}e^{0}=1\\ e^{A_{1}B_{0}}\\ e^{\infty}=\infty\end{cases}. \]\footnote[1]{Examples: 1) $\left(1+\left(\frac{x}{a}\right)^{3}\right)^{\frac{a}{x}}=f(x)$ gives, for $x=0$, $1^{\infty}$; $f(0+h)=e^{\frac{a}{h}l(1+(\frac{h}{a})^{3})}=e^{\frac{a}{h}(\frac{h^{3}}{a^{3}}-\frac{h^{6}}{2a^{6}}....)}=e^{\frac{h^{2}}{a^{2}}-\frac{h^{5}}{2a^{5}}..}$, and $\lim f(0+h)=e^{0}=1$. 2) $f(x)=(2-x)^{\frac{1}{1-x}}$, which, for $x=1$, gives $1^{\infty}$. If $1+h$ is put instead of $x$, one has $f(1+h)=(1-h)^{-\frac{1}{h}}=e^{-h^{-1}l.(1-h)}=e^{h^{-1}(\frac{h}{1}+\frac{h^{2}}{2}....)}=e^{1+\frac{h}{2}+\frac{h^{2}}{3}+....}$; $f(1)=\lim f(1+h)=e$. 3) $f(x)=(1+x)^{\frac{1}{x^{2}}}=1^{\infty}$, for $x=0$. $f(0+h)=e^{\frac{1}{h^{2}}l(1+h)}=e^{h^{-2}(\frac{h}{1}-\frac{h^{2}}{2}+\frac{h^{3}}{3}...)}=e^{h^{-1}-\frac{1}{2}+\frac{h}{3},..}$ and $\lim f(0+h)=f(0)=\infty$.}
% printed mark 1) stands after e^{A_1B_0} inside the cases; footnote set after the display

% the note 1) of p. 73 runs over onto the foot of this page: set in full at its mark

% --- p. 75 ---
\setcounter{footnote}{0}%
\opage{75}Since the limit of the function, in showing the transition of the finite into the infinite, and in vanishedness still mirroring what has vanished, is as much in becoming as in being, the increment $h$ plays the real leading part: in $f(a)=\lim f(a\pm h)=f(a\pm 0)$ it depends, so to speak, on the movement in $h$'s passage through $0$. A couple of
examples can illuminate the matter:
\[ f(x)=\left(\frac{1}{a-x}\right):\frac{1}{\sqrt{x^{2}-a^{2}}} \]

gives, for $a=x$, the form $\frac{\infty}{\infty}$, and since thus $f(a+h)$
\[ =-\sqrt{\frac{2a}{h}+1}, \]
we get $f(a)=\lim f(a+h)=\infty$, thus
in this case $\frac{\infty}{\infty}=\infty$. If, on the other hand, one has $f(x)=\frac{\mathrm{tg}\,x}{\sec x}$
then for $x=\frac{\pi}{2}$ we get $f(x)=\frac{\sin x}{\cos x}:\frac{1}{\cos x}=\sin x=1$, and consequently $\frac{\infty}{\infty}=1$. If one has $f(x)=(a-x)\,l^{\frac{1}{a-x}}$, which for
\[ x=a, \]
gives the form $0.\infty$, then $f(a+h)=-he^{-\frac{1}{h}}$
and $f(a)=\lim f(a+h)=0$; thus $0.\infty=0$; if, on the other hand,
$f(x)=\sin x.\cot x$, which, for $x=0$, likewise gives us the form
\[ 0.\infty, \]
then $f(0)=\lim f(0+h)=\lim\sin(0+h)\frac{\cos(0+h)}{\sin(0+h)}$
\[ =\lim\cos h=1, \]
thus $0.\infty=1$. If, in order also to
cite this third possibility, $f(x)=\frac{1}{(1-x)^{2}}(1-x)$,
which, for $x=1$, likewise gives the form $0.\infty$, one obtains
by reduction
\[ \frac{(1-x)^{2}}{(1-x)^{2}(1-x)}=\frac{1}{1-x}, \]
and consequently $0.\infty=\infty$

Here, then, in $f(a)=\lim f(a\pm h)=f(a\pm 0)$, whichever
of the particular forms one may choose, and however
the conditions may be changed, there are in the one limit three limits,
on which it essentially depends. Insofar as $h>0$, there lies
% --- p. 76 ---
\setcounter{footnote}{0}%
\opage{76}here between the finite positive $h$ and $+0$, on the one side,
and between the finite negative $h$ and $-0$ on the other
side, a magnitude which, from the positive as from the negative
side, is in equal degree finite and infinite, since it is the
smallest possible finite magnitude. The positive $+0$ thus denotes neither $0$ nor a positive magnitude, but the limit
between the positive magnitude and its negation, $0$; likewise the negative $-0$ denotes the limit between the
negative magnitude and its negation. If now, e.g., in $\lim f(0\pm h)=\lim\cot(0\pm h)=f(0\pm 0)=\pm\infty$ there is on
both sides of $0$ a dialectical magnitude of transition, a magnitude
that is infinite-finite, since, as the smallest
possible finite magnitude, it is infinitely small; then here too there is
on both sides of $\infty$ a dialectical magnitude of transition, a
magnitude that is finite-infinite, since it, as the largest
possible finite magnitude, is infinitely large. In this way
every transition of magnitude is transformed into a dialectical
magnitude of transition: infinite in its smallest possible finiteness
the magnitude of transition is a \emph{differential}; infinite in its largest
possible finiteness the magnitude of transition is, insofar as it is
conceived as a sum of differentials, an \emph{integral}.

\begin{center}

b) The Differential: The Dialectic of the Moment.\end{center}

``Before we go further, it will be expedient
to determine more closely the concept of infinitely small magnitudes.
An \emph{infinitely small or vanishing} magnitude
is one that is smaller than any magnitude, however small,
of the same kind. Such a magnitude is not zero
(nothing), for then it would be no magnitude, but in comparison
with a finite magnitude it is \emph{to be reckoned as
zero}; for otherwise it would not be smaller than every assignable
magnitude. If, then, $a$ is a finite magnitude and $\alpha$ an infinitely small one, then $a\pm\alpha=a$, $a\alpha=0$. Every uninterrupted
motion gives an example of infinitely small magnitudes, in that
the path covered grows every moment by such a
% --- p. 77 ---
\setcounter{footnote}{0}%
\opage{77}magnitude. The ratio between two infinitely small magnitudes $\alpha$ and $\beta$ can be finite, and they are then said to be of the same
\emph{order}. If $\alpha$ is an infinitely small magnitude, then $\alpha^{2}$ will be
infinitely smaller and vanishing in comparison with $\alpha$. $\alpha^{2}$ and every magnitude whose ratio to this is finite
is said to be infinitely small of the second order, one thinking
of $\alpha$ as being of the first order. Likewise $\alpha^{3}$ and
every magnitude whose ratio to this is finite is said to
be infinitely small of the third order, and so on.''

``If $a$ is a finite magnitude, then $\frac{a}{\alpha}$ and every magnitude
whose ratio to this is finite is said to be an infinitely
large magnitude, $\frac{a}{\alpha^{2}}$ to be infinitely large of the second order,
etc. It goes without saying that one can also conceive
both infinitely small and infinitely large magnitudes of orders
whose exponents are not whole numbers. An infinitely
large magnitude is not the same as that which is indicated by the
sign $\frac{a}{0}$ or $\infty$, which one might call the absolute
infinite, the limit that lies outside all greatness.
One often writes, however, $0$ and $\infty$ where what is properly
meant is an infinitely small and an infinitely large magnitude.''\footnote[1]{Chr. Jürgensen. Begyndelsesgruudene af den mathematiske Analyse og dens Anvendelser. Kbhvn 1855. p. 41--43.} % print: Begyndelsesgruudene (misprint)

To convince oneself that, in assuming
infinite magnitudes, no untenable concepts are smuggled in,
one should note that every determination of greatness is essentially
a determination of relation, and that a determination of relation
can be constant although the terms increase or decrease
to infinity. Only in the doctrine of infinite magnitudes
can the concept of magnitude be fully developed, for only in this
doctrine does it become sufficiently clear that magnitude is relation
and relation magnitude. In the differential coefficient the % print: Differentialcoifficienten (misprint)
magnitude is taken up into the designating relation as
% --- p. 78 ---
\setcounter{footnote}{0}%
\opage{78}relational magnitude; the different orders of the differential coefficients rest % print: Differentialcoifficienternes (misprint)
on a system of relational magnitudes, and the principles of differentiation
specify the method determined by the peculiar relational consequences.

\begin{center}α)

\emph{Differential and Differential Coefficient: Relational Magnitude.}\end{center}

A general greatness, a general smallness is
as yet no magnitude; only when the greatness and the smallness
are more closely determined does magnitude emerge. Since now
immediate determinateness coincides with limitation,
and the immediately limited is a finite, the view suggests itself
that only the limited, that is, finite magnitude
is magnitude in the proper sense. As
finite magnitude, magnitude can be held fast in representation
and determined for intuition; just as, as regards denominate numbers,
an ell can be determined for intuition partly immediately,
partly by comparison with a quarter, so the finite numbers of the number series can be apprehended partly immediately,
partly by comparison with one another;
each number then has in this its finiteness a name of its own, a designation of its own, % print: et (misprint)
and can to that extent appear as a whole existing
for itself. Now in that the magnitudes thus conceived
enter into connections with one another, it shows that, the
more the finite representation has the upper hand over thinking,
the more emphasis falls on the given numbers,
the existing terms of the relation; only gradually as comprehending
thinking gains power over immediate representation
are the terms transformed into moments, so that the emphasis
comes in the end to rest, not on the immediate
terms, but on the relation. In expressions such as: $1$, $\frac{2}{2}$, $\frac{a}{a}$, $\mathrm{Lim}\frac{x}{x}$
and $\lim\frac{x}{x}$ unity comes to be seen from essentially different
points of view.

% --- p. 79 ---
\setcounter{footnote}{0}%
\opage{79}So long as representation involuntarily clings to the sign: $1$, it looks as if unity were an immediate
something, an existent, which one could behold and point out directly;
this immediacy of representation vanishes when
one begins to conceive unity as a relation: $\frac{2}{2}=\frac{3}{3}$
\[ =\frac{4}{4}\ldots \]
With cognition of relation comes
insight, but
the terms still have a special significance over against the
relation whose terms they are: cognition of relation is still
superficial. With the designation $\frac{a}{a}$ a considerable advance occurs,
in that the general expression replaces the particular ones,
and it is brought to consciousness that magnitudes can not merely
enter into different relations, but that the relation necessarily
belongs to the essence of the magnitude. The algebraic
designations are, however, burdened with the accompanying notion
that every magnitude, however arbitrary, must be finite.
The equations: $\frac{a}{b}=3$, $\frac{b}{c}=4$, $\frac{c}{d}=e$ have indeed their
significance in the relations of $a$ and $b$, $b$ and $c$, $c$ and $d$; but
the illusion that these determinate relations would dissolve
into indeterminacy, when $a$, $b$, $c$ and $d$ ceased to be finite
magnitudes, has not yet been removed. This illusion vanishes:
when it is seen what lies in the designation:
\[ \mathrm{Lim}\frac{x}{x}=\lim\frac{x}{x}=\frac{\infty}{\infty}=1, \]
for however one may explain $\frac{\infty}{\infty}$ and $\frac{0}{0}$ to oneself, this at any rate stands firm, that unity here comes to be
conceived as the pure expression of relation. That with such an expression of relation it does not at all depend on the finite
terms, but solely and alone on the relation determined by the terms,
can hardly be made known in a more emphatic
way than by the determinate relation's remaining
unchanged after $\infty$ and $0$ have taken the place of the terms.

% --- p. 80 ---
\setcounter{footnote}{0}%
\opage{80}That the concept of magnitude must be completed in the concept of function
is evident from the law of function, as stated above,
embracing at once the finite and the infinite.
From $f(x)=y$ and $f(x+h)=y+k$ follows $\frac{f(x+h)-f(x)}{h}$
\[ =\frac{k}{h}, \]
and lim $\frac{f(x+h)-f(x)}{h}=\lim\frac{k}{h}=\frac{dy}{dx}$.
In $f(x+h)-f(x)=k$, $h$ is a magnitude for itself, $k$ likewise; $h$ and $k$ have in finite comparison immediate validity; such
validity cannot, on the other hand, be ascribed to $dx$ and $dy$ in $f(x+dx)=y+dy$, wherefore exact analysis also
constantly resists their being conceived without further ado as
magnitudes; in pure immediacy they are indeed not
magnitudes either: as relational magnitudes they are dialectical. If one overlooks
the dialectical, one can describe the differential
as one will; one never hits the matter. The differential $dy$ is, according to $f(x+dx)-f(x)$, an infinitely small difference;
but the infinitely small difference has in itself the contradiction
of being like $0$ and yet greater than $0$. If the differential
is conceived not as magnitude but as disposition to magnitude, a
disposition to magnitude must surely in principle itself be magnitude.
If one determines the differential not as an arbitrarily
small but as an infinitely vanishing magnitude, how
can one then in $f(x+dx)=y+dy$ set $dx=dx$, while $dy$ varies? How can a constantly becoming, an
incessantly vanishing, i.e. not-being magnitude be
constant? If, on the other hand, one acknowledges the dialectical,
the differential, as infinitely small difference, as element of magnitude,
as moment of relation, will illuminate the concept of magnitude from all sides.
If one has $f(x+dx)=4(x+dx)=4x+4dx=y+dy$, then it is undeniable that $dy=4dx$, and to that extent
we seem in $dy$ and $dx$ to have immediate magnitudes; but in that the emphasis falls, not on $dy$ and $dx$ for themselves, as if
by $dx+dx+dx+dx$ we could obtain $dy$, but on $\frac{dy}{dx}$
% --- p. 81 ---
\setcounter{footnote}{0}%
\opage{81}thus on the relation; the differentials become hereby expressly % print: bive (misprint)
determined as relational moments. In the equation $dy=4dx$, $4$ is the coefficient by which $dx$ is to be multiplied to
give $dy$, but as an immediate magnitude the coefficient
does not yet correspond to its concept: differential coefficient.
To be recognized as differential coefficient, $4$ must from an % print: Differentialcoefficent (misprint)
immediate magnitude be transformed into an expression of relation reflected into itself:
$\frac{dy}{dx}$. That $dy=\frac{dy}{dx}dx$ must not be understood tautologically,
as if here stood $a=\frac{a}{b}.b$; for $\frac{dy}{dx}=4$ denotes
expressly a magnitude transformed into a relation, a relational magnitude:
\emph{The differential coefficient is essentially a quotient.} The differential coefficient, which in the exact % print: Differentialcoefficenten (misprint)
language is the limit of the relation between the corresponding
and simultaneously vanishing increments of the dependent
and independent variable, thus gives the differential its
significance, not conversely.

The intuitive image of dialectic is motion, in particular
motion in space. Newton, who, in order to solve phoronomic
problems, takes geometry to his aid, presents the
uniformly growing time as abscissa, the space traversed in time
as ordinate. The coordinates of the curve, which
represent the ever-changing, decreasing and increasing magnitudes,
time and space, he calls ``Fluents'', and
designates by $x$, $y$ etc., but the velocities with which
they decrease or increase he calls ``Fluxions'', and designates
them by $\dot{x}$, $\dot{y}$ etc. It goes without saying that the fluxion of the abscissa $\dot{x}$, since it presents the velocity for the time,
which it after all represents, and time is thought to grow uniformly
(``\textit{uniformiter fluens}''), must be regarded as constant. It
follows from this that the second, the third and all following
fluxions of $x$ are equal to zero, further, that the fluxion of $x$ is the measure by which the fluxion of $y$ is measured, and
can thus be assumed as unity. Now since time and space and
% --- p. 82 ---
\setcounter{footnote}{0}%
\opage{82}consequently the lines by which they are represented, $x$ and $y$, have in the instant, i.e. in time $0$, the one the velocity $\dot{x}$, the other $\dot{y}$, the increase and decrease of these lines, since this
is directly proportional to the velocity, can be expressed
by $\dot{x}0$ and $\dot{y}0$. This increase or decrease of the fluents, proportional to the velocity, which must be infinitely small, since
the velocity can be regarded as constant only with respect to the single instant,
Newton designates by the expression:
\emph{Moment.}\footnote[1]{Cf. Weissenborn. Die Principien der höheren Analysis. p. 21--57.}

Der Gedanke kann nicht richtiger bestimmt werden,
als Newton ihn gegeben hat.\footnote[2]{Here too we have an example showing how a concept which, so long as it is held in logical generality, proves entirely satisfactory, can nevertheless, when it is a matter of an articulating discretion, an unfolding of punctual determinations, be deficient and unsatisfactory. „Man hat öfter die Fluxionsmethode der jetzt vielfach gebräuchlichen Grenzmethode nahe gestellt, und dies grossentheils mit gutem Rechte, indem Newton stets auf Verhältniss der Fluxionen $\dot{y}$ und $\dot{x}$ kommt. Auf diese Weise umgeht er die Schwierigkeit, die den damaligen Mathematikern ein unüberwindliches Hinderniss schien, nämlich das Rechnen mit unendlich kleinen Grössen oder Null. Denn so lange die Incremente endlich sind und eine angebbare Grösse besitzen, hat ihre Anwendung im Calcul kein Bedenken, Newton lässt sie auch bis zu Null abrechnen, und nennt sie denn Momente oder auch („\textit{Princ. phil. nat.}« \textit{Lib. II. Sect II. Lemma II. pag. 251} der ersten Ausgabe und „\textit{Tractatus de quadratura curvarum.}“ \textit{Prop. I. Problem I. Demonstratio.)} »\textit{incrementa momentanea}“ oder „\textit{momenta nascentia sive evanescentia}“, er hütet sich aber wohl, mit diesen zu rechnen, sondern lässt im Zustande des Verschwindens die Fluxionen, endliche und bestimmte Grössen, gewissermassen als Reserve an die Stelle der zum weiteren Dienste untauglich gewordenen Momente einrücken, auf diese Weise das Rechnen mit unendlich kleinen Grössen, welche Null sein und doch noch einen von Null verschiedenen Werth besitzen sollen, auf das Glücklichste vermeiden; und wenn er (»Tractatus de quadratura curvarum.“ Introductio) sagt: „Volui ostendere, quod in methodo fluxionum non opus sit figuras infinite parvas in Geometriam introducere«, und ebendaselbst: Quantitates mathematicas non ut ex partibus quam minimis constantes, sed ut motu continuo descriptas hic desidero«, so muß man unbedingt zugeben, daß er nicht mehr verspricht als er in der That gehalten hat. Unstreitig würde der große Vortheil, den seine Methode durch die Vermeidung des zu Widersprüchen führenden Begriffs des Uendlichkleinen, wodurch zugleich der unnatürlichen Auffassung der Curven als Polygone vorgebeugt wurde, über die Differenzialmethode, wie sie lange Zeit hindurch dargestellt ward, gewonnen hatte, — unstreitig würde dieser große Vortheil ihr schon frühe den Vorrang vor der Leibnizischen Theorie verschafft haben, wenn ihn nicht Newton dadurch erkauft hätte, daß er die Rechnung aus dem Gebiete der Geometrie weg in das der Phoronomie hinüberspielte. Dazu kam aber noch als zweites sehr wichtiges Moment, daß die Fluxionsrechnung Newtons eines hinlänglich passenden und ausgebildeten Algorithmus entbehrte, was sich besonders beim Uebergehen von den Fluxionen zu den Fluenten als ein Nachtheil sehr bemerklich macht. Denn es war ja hauptsächlich die Auffindung eines solchen Algorithmus, um die sich damals Alles drehte. So kann es denn nicht wunderbar sein, wenn die Leibnizische Differenzialrechnung, welche alle Forderungen, die in dieser Beziehung gestellt werden können, auf das Vollkommenste befriedigt, der Fluxionsrechnung trotz ihrer sichereren Begründung sehr bald den Rang ablief.“ Weissenborn. Anfrt. Skr. p. 57--58.} Ich trenne dabei die % print: Udviling (misprint)
% --- p. 83 ---
\setcounter{footnote}{0}%
\opage{83}Bestimmungen ab, die der Vorstellung der Bewegung und
der Geschwindigkeit angehören, (von welcher er vornehmlich
den Namen Fluxionen nahm), weil der Gedanke hierin
nicht in der gehörigen Abstraction, sondern konkret, vermischt
mit ausserwesentlichen Formen erscheint. Diese Fluxionen
erklärt Newton (\textit{Princ. mathem. phil. nat. L. I. Lemma XI.
Schol.}) dahin, daß er nicht \emph{untheilbare} --- eine Form,
deren sich frühere Mathematiker, Cavalleri und andere, bedienten,
und welche den Begriff eines an sich bestimmten
Quantums enthält, verstehe, sondern verschwindende Theilbare.
Ferner nicht Summen und Verhältnisse bestimmter
% the note 2) of p. 82 runs over onto the foot of this page: set in full at its mark
% --- p. 84 ---
\setcounter{footnote}{0}%
\opage{84}Theile, sondern die \emph{Grenzen} (limites) der \emph{Summen} und
Verhältnisse. Es werde die Einwendung gemacht, daß verschwindende
Grössen kein leztes Verhältniß haben, weil
es, ehe sie verschwunden, nicht das Lezte, und wenn sie verschwunden,
keines mehr ist. Aber unter dem Verhältnisse
verschwindender Grössen sey das Verhältniß zu verstehen,
nicht ehe sie verschwinden, und nicht \emph{nachher}, sondern
mit dem sie verschwinden (quacum evanescunt). Ebenso
ist das erste Verhältniß werdender Grössen, das mit dem
sie werden.“\footnote[1]{Hegel. Wissensch. der Logik. 1ster Th. Berlin 1833. p. 302--303.}


\begin{center}β) The Different Orders of the Differential Coefficients: A System of Relational Magnitudes.\end{center}

Although it is, essentially seen, the differential coefficient
that gives the differentials of the dependent and the independent
their significance, it must nevertheless, as a relation of the differentials,
necessarily presuppose the differentials themselves, and comes
to that extent to rest on their peculiar validity.
In mathematical practice it may perhaps have its
significance to keep exactly to what is negative in the limit,
and to set the differential coefficient $\frac{dy}{dx}=\frac{0}{0}$; but thereby
the thinking that works at insight and comprehension is
by no means satisfied. In the expression $\frac{0}{0}=4$, $\frac{0}{0}$ can impossibly refer to something present, but only to something past
and something to come. Were one really to assume $\frac{0}{0}=4$, one would have to assume a relation without terms of relation, a relation
of nothing to nothing; $\frac{0}{0}$ is thus not $4$, nor does it denote that $4$ still is, after the corresponding expression of relation has vanished, but merely makes known that $4$ has been,
% --- p. 85 ---
\setcounter{footnote}{0}%
\opage{85}in that here a now vanished relation has been. If $\frac{0}{0}$ is, strictly
speaking, to denote the differential coefficient, then the differential coefficient
is no magnitude; if the emphasis is to fall on $4$, then $4$ as immediate magnitude is no differential coefficient.
Mathematics also does not simply set $\frac{0}{0}=4$; by the designation for what has vanished it is still mindful of
what has vanished, recalls what it really was that vanished, and how it vanished: $\lim\frac{4(x+h)-4x}{h}=\frac{0}{0}=4$. Insight is here in the strictest sense a recollection;
in recollection the past is seen as a present:
what has been negated reflects its being in the negation. In that $\frac{dy}{dx}$ is put in place of $\frac{0}{0}$, the vanishing terms can constitute a
still existing relation, a relational magnitude of determinate
value.

As determinate relational magnitude the differential coefficient can
be either finite or infinite, either constant or
variable. Thus in $f(x+dx)=a(x+dx)=ax+adx=y+dy$ the differential coefficient $\frac{dy}{dx}=a$ is constant; in $f(x+dx)=(x+dx)^{2}=x^{2}+2xdx+dx^{2}=y+dy$, on the other hand, $\frac{dy}{dx}=2x$ is variable. As variable magnitude the
differential coefficient $f'(x)=\frac{dy}{dx}=2x$ can be differentiated again;
whereas $f'(x+dx)=2(x+dx)=f'(x)+f''(x)dx+\ldots=y_{1}+dy_{1}$ gives a differential coefficient: $\frac{dy_{1}}{dx}=2$, which, as a constant
magnitude, cannot be differentiated. Since the differential coefficient is already as such a relational magnitude, the
differentiated differential coefficient is, as a relational magnitude determined by a relational
% --- p. 86 ---
\setcounter{footnote}{0}%
\opage{86}magnitude, $\frac{dy_{1}}{dx}=\frac{d\frac{dy}{dx}}{dx}$, a relational magnitude
of second order $\frac{d^{2}y}{dx^{2}}$; in this way the
relational magnitudes can develop into a system of lower and
higher orders. From $f(x)=y=x^{n}$, e.g., one obtains:
\[ f(x+h)=f(x)+f'(x)\frac{h}{1}+f''(x)\frac{h^{2}}{1.2}+f'''(x)\frac{h^{3}}{1.2.3}+\ldots \]
\[ (x+h)^{n}=x^{n}+nx^{n-1}\frac{h}{1}+n(n-1)x^{n-2}\frac{h^{2}}{1.2}+n(n-1)(n-2)x^{n-3}\frac{h^{3}}{1.2.3}+\ldots \]
\[ =y+\frac{dy}{dx}\frac{h}{1}+\frac{d^{2}y}{dx^{2}}\frac{h^{2}}{1.2}+\frac{d^{3}y}{dx^{3}}\frac{h^{3}}{1.2.3}+\ldots \]
What inner infinity such a system of relational magnitudes
can embrace is seen when we choose a function of
several mutually independent variables, e.g. $f(x_{1},x_{2},\ldots x_{t})=u$. In the series development corresponding to $f(x_{1}+h_{1},x_{2}+h_{2}\ldots x_{t}+h_{t})\sim e^{\frac{du}{dx_{1}}h_{1}+\frac{du}{dx_{2}}h_{2}+\ldots\frac{du}{dx_{t}}h_{t})}$, where % print: a stray ) in the exponent, no + after the dots (misprint)
$\frac{1}{[r_{1}+r_{2}+\ldots r_{t}]}\left(\frac{du}{dx_{1}}h_{1}+\frac{du}{dx_{2}}+\ldots\frac{du}{dx_{t}}h_{t}\right)^{(r_{1}+r_{2}+\ldots r_{t})}$ is the
general term, there is, since the changes corresponding to the symbolism,
by which exponents of powers are transformed into indices of differentiation,
have been made, $\frac{d^{r_{1}+r_{2}\ldots r_{t}}u}{dx_{1}^{r_{1}}dx_{2}^{r_{2}}\ldots dx_{t}^{r_{t}}}\frac{h_{1}^{r_{1}}}{[r_{1}]}\cdot\frac{h_{2}^{r_{2}}}{[r_{2}]}\ldots\frac{h_{t}^{r_{t}}}{[r_{t}]}$ again the general term of this general term, in which
the differential coefficient can thus present relations of relations
to infinity.\footnote[1]{That the inner infinity can likewise present itself to intuition from infinitely many sides is shown by analytic geometry. If, to choose an example from the generation of surfaces, a surface is to form itself by the motion of a certain line $\{F_{1}(x,y,z,\alpha,\beta,\gamma\ldots..)=0, F_{2}(x,y,z,\alpha,\beta,\gamma\ldots..)=0$ in which $\alpha,\beta,\gamma\ldots.$ are constants, which nevertheless vary from one position of the line to another: then, for the law of this variation to be determined, certain conditions must be presupposed, e.g. that the generator (Generatrix) during its motion shall intersect a number of guide lines. Supposing that n constants enter into the above equations, the conditions must make it possible to get these n constants eliminated, for the surface that is to be formed is the sum total of all the positions of the generator, and can therefore not depend on any single one in particular, i.e. the equation of the surface must be independent of $\alpha,\beta,\gamma\ldots.$ We must thus on the whole be able to form $(n+1)$ equations between $x,y,z,\alpha,\beta,\gamma\ldots.$ To make quite clear what infinity of surfaces can be formed by a given generator, let us for a moment imagine that for the first n of the $(n+1)$ equations we have an entirely determinate form of function, while the form of the $(n+1)$th equation is indeterminate; here, then, an infinity of forms is possible for this $(n+1)$th equation; for each single form we choose we get a surface, thus an infinity of surfaces. If we consider, moreover, that we can let each single one of the $(n+1)$ equations in turn have an indeterminate form, then it becomes clear what infinity of surfaces the one system of $(n+1)$ equations can give occasion to by the elimination of the n constants. In this infinite swarm a systematic division is moreover possible. To undertake such a division, we can take into consideration in particular 1) the generator and 2) the law of generation. All the surfaces that have the same generator can then be referred to the same class; each single class can further be divided into families, which all have the characteristic of the class, namely the same kind of curves as generators; the characteristic of the family is seen from the law according to which the generatrix moves. Within each family, on the other hand, different species arise, in that the same law appears in different modifications, and within the same species the different individuals, in that determinate magnitudes are inserted in the particular form. Cylinders thus constitute by themselves a large family, which belongs under the class of rectilinear surfaces; the characteristic of the family is that all positions of the generator are parallel. The generator can namely, during its motion, continue to be parallel to itself, and yet produce very different species of cylinders: cylinders of revolution, hyperbolic cylinders etc. Cylinders of revolution, e.g., yield an infinity of individuals, for for each different radius we give the circle perpendicular to the generator, we get a different individual. (The categories: class, family, species, individual are cited here after a lecture by Prof. A. Steen).}

% --- p. 87 ---
\setcounter{footnote}{0}%
\opage{87}How differential coefficients of different order acquire
a different significance in their application can be seen
% the note 1) of p. 86 runs over onto the foot of this page: set in full at its mark
% spaced words Klasse, Familier, Arter, Individer stand inside footnote 1 (p. 86), left unchanged
% --- p. 88 ---
\setcounter{footnote}{0}%
\opage{88}from examples drawn from geometry as well as from mechanics.
If we choose from geometry the problem of contact, and let $f(x)=y$, $\Phi(\xi)=\eta$ be the equations of two curves which
have the point $(x,y)$ in common, then the difference of the consecutive ordinates corresponding to the point $(x,y)$ is to be expressed by:

\[ f(x)-\Phi(x)+(f'(x)-\Phi'(x))\frac{h}{1}+\ldots(f^{n}(x)-\Phi^{n}(x))\frac{h^{n}}{[n]}\ldots, \]

\[ y-\eta+\left(\frac{dy}{dx}-\frac{d\eta}{d\xi}\right)\frac{h}{1}+\ldots\left(\frac{d^{n}y}{dx^{n}}-\frac{d^{n}\eta}{d\xi^{n}}\right)\frac{h^{n}}{[n]}\ldots, \]

in which $y=\eta, x=\xi$. The magnitude of this difference depends,
since $h$ has the same infinitely small value, on which
term is the first that is not $0$, i.e.\ on the order
of the contact; if e.g.\ $y=\eta, \frac{dy}{dx}=\frac{d\eta}{d\xi}\ldots.\frac{d^{n}y}{dx^{n}}=\frac{d^{n}\eta}{d\xi^{n}}$, but $\frac{d^{n+1}y}{dx^{n+1}}\gtrless\frac{d^{n+1}\eta}{d\xi^{n+1}}$,
so that the contact
is of the $n$th order; then $\alpha)$ the difference is the smaller, the greater $n$ is, and $\beta)$ it changes sign with $h$ when $n$ is even, but not when
$n$ is odd. To $(\alpha)$ there corresponds a so much the more intimate contact, i.e.\ the two curves have so many more elements in common, the greater $n$ is; by $(\beta)$ it is decided whether the contact is bound up
with intersection or not. In mechanics the expression
for velocity: $v=\frac{ds}{dt}$, and for force: $\Phi=\frac{d^{2}s}{dt^{2}}$ furnish easily grasped examples.
% the note 1) of p. 86 runs over onto the foot of this page: set in full at its mark

% --- p. 89 ---
\setcounter{footnote}{0}%
\opage{89}From the application of the differential coefficients it is seen further
that the differentials themselves cannot possibly be pure zeros,
but must, as determinate relational moments, be magnitude-dispositions,
and consequently in principle themselves magnitudes. Thus
the condition for
contact ``of the nth order shows that the two
curves have the following points in common:

\begin{quote}
\emph{Abscissas:} \emph{Ordinates:}
\end{quote}
\[ x \quad y, (\eta) \]
\[ x+dx \quad y+dy, (\eta+d\eta) \]
\[ x+2dx \quad y+2dy+d^{2}y, (\eta+d\eta+d^{2}\eta) \]
\[ x+n dx \quad y+n dy+\frac{n(n-1)}{1.2}d^{2}y\ldots d^{n}y, \]
\[ (\eta+n d\eta+\ldots d^{n}\eta) \]

Contact of the nth order will thus say, in other words, that the two curves have $n$ consecutive elements in common.
Mathematics itself makes a distinction between contact of the $n$th order and contact
of the 1st order, and teaches thereby at the same time that
to have $n$ elements in common is, for $n>1$, more than to have one element in common, or in symbols: $n dx>dx$; but when $n dx>dx$, although $n$ is a finite number, $dx$ cannot possibly be $0$.\footnote[1]{In descriptive geometry, in order also to see the matter from this side, it is proved by applying the projective properties of the conic sections that, when a hexagon is inscribed in a conic section, the 3 points in which the opposite sides of this hexagon intersect one another, two and two, lie in one and the same straight line. A pentagon is, in consequence of this, to be regarded as a hexagon with \emph{one infinitely small side}; this infinitely small side is nevertheless more than a point, for it indicates the direction of the tangent to the conic section (and a point cannot indicate any direction; therefore the 3 points in which the opposite sides of the pentagon together with the tangent intersect one another, two and two, likewise lie in one and the same straight line.}

That the differential as such cannot immediately be $0$ is seen, if possible, still more clearly when we, as concerns
% --- p. 90 ---
\setcounter{footnote}{0}%
\opage{90}motion, consider the relation of space and time.
For while $\frac{ds}{dt}=v$ denotes a relation of space and
time and thereby a velocity determinate in its disposition, $\frac{0}{0}$, which
after all expresses an absolute negation of time as well as of
space, and thus falls in pure abstraction outside
time and space, can impossibly indicate a velocity-disposition for a motion
that is to take place in time and space. By one-sidedly
holding fast to the abstract expressions of calculation one misjudges the
essence of the matter, and entangles oneself in a contradiction. In the equation $f(t+dt)=s+ds$ $dt$ is constant, while $ds$ can vary, according as t varies;
how dialectical this constant magnitude
nevertheless is, is revealed among other things by the fact that for
a varying velocity one sets $v=\frac{ds}{dt}$, as if $v$ were constant in the time-element $dt$. If now $dt=0$, then there is here neither velocity nor velocity-disposition; if $dt$, on the other hand, is not $0$, but a real time-element, then the velocity continues
to vary even within the chosen time-element; here there is then
a varying which, inasmuch as it corresponds to disposition and possibility, can correspond to its concept only
by hovering between Something and Nothing, and thus remaining
dialectical.

\begin{center}γ) \emph{Principles of differentiation: the method determined by the relational consequences.}\end{center}

What different impressions the differential calculus introduced by Leibniz
made on his contemporaries is seen
among other things from the fact that while mathematicians such as Jacob and Johan
Bernoulli applied it with success to the most difficult % print: Bernouilli (misprint)
problems, Leibniz's friend Huygens would not abandon the
method of the Ancients, and the Dutchman Nieuwentijt even came forward
polemically, objecting: 1) that Leibniz, no more
than Barrow and Newton, in whom the infinite
magnitudes already occurred, could explain how these
magnitudes differed from nothing or zero; 2) that it was not
% --- p. 91 ---
\setcounter{footnote}{0}%
\opage{91}clear how the differentials of higher order were to be
distinguished from those of the first order, and $3)$ that the differential method
could not be applied to the exponential function.\footnote[1]{Weissenborn, work cited, p. 98.}
The wavering and uncertainty which runs through Leibniz's
explanations and defenses shows clearly enough that one had
entered upon a new path, a way whose direction and
character could be investigated and more closely
tested only by degrees.\footnote[2]{What Leibniz writes to Oldenburg in a letter from Paris, 27 August 1676: \textit{Fateor tamen nondum me quicquid in hoc genere desiderari potest consecutum, quanquam maximi momenti esse sciam,} holds not only for that time, but in part also for the time that followed. Cf.\ Gerhardt: Die Entdeckung der Differentialrechnung durch Leibnitz. Halle 1848. p. 26. ``\textit{M. Maclaurin, l'un des plus grands Géometres de ce siecle} (18th Cent.) \textit{éléva le premier sa voix au milieu des acclamations qu'attiroient les utilités réelles du Calcul des Infiniment-Petits. Il osa avertir du danger qu'on couroit en s'y livrant avec tant de confiance. Il n'y a rien de concevable, dit il, dans la supposition d'un nombre Infiniment-Grand ou Infiniment-Petit; \& le passage du Fini à l'Infini est obscur et incomprehensible.}'' \textit{M. Saverien. Histoire critique du Calcul des Infiniment-Petits (in ``Application de la Géométrie ordinaire'')} p. 7—8. ``Da... Einiges in Newtons Verfahren nicht hinreichend klar und deutlich war, so könnte es nicht fehlen, daß dasselbe mannigfache Angriffe zu erleiden hatte. Gegen dieselben trat Maclaurin auf (``\textit{Treatise of Fluxions}''), um einerseits alle Zweifel, die noch in der Fluxionsrechnung übrig seien, zu zerstreuen, andererseits die Leibnizische Differentialmethode anzugreifen. Der letzteren machte er zwei Dinge zum Vorwurfe, einmal nämlich, daß das Uendlichkleine für den menschlichen Verstand unfaßbar, sodann daß Leibnitz selbst den Begriff des Uendlichkleinen für eine bloße ``Fiction'' erklärt habe.'' Weissenborn. p. 58. Yet another characteristic trait of the above-mentioned \textit{Histoire du Calcul des Infiniments-Petits.} After a short description of the Archimedean method of limits the author continues on p. XIV as follows: ``\textit{Cette méthode fut long-tems en usage. Dans la vûe de la} simplifier, on crut qu'on pouvoit se passer de concevoir des figures circonscrites et inscrites dans des aires curvilignes, comme étant toûjours finies et déterminables. L'expédient qu'on trouva pour cela, fut de substituer à ces figures finies et déterminables des élémens indivisibles ou Infiniment-Petits. Le nombre de ces élémens étant supposé infini, leur somme est évidemment égale à l'aire curviligne. On doit cette idée à Cavallieri. Quelque séduisante qu'elle lui parut par sa simplicité, elle ne le rassura pas sur son entiere certitude. Ce nombre infini lui paroissoit géométriquement indéterminable. Il tâchoit donc d'écarter cette idée de nombre infini, lorsqu' il reconnut bien des difficultés qu'il ne put résoudre. C'en fut assez pour lui rendre suspecte sa nouvelle Géométrie. Afin de l'appuïer, Cavallieri joignit des démonstrations incontestables à celles qu'il avoit déduites de ces principes hypothetiques; mais ce mêlange de vérités \& de suppositions, bien loin de rassurer les esprits, les inquiéta. On vit alors pour la premiere fois (selon la juste remarque de M. Maclaurin, dans l'Introduction de son Traité des Fluxions) des disputes parmi les Géometres.''} But while the analysis of the infinite % print: Uendlichkleine, Uendlichkleinen (misprint)
% --- p. 92 ---
\setcounter{footnote}{0}%
\opage{92}thus certainly gives us in one respect something entirely new,
this new is nevertheless, when the transitional stages
which the history of science displays are taken into consideration,
a natural consequence of the original presuppositions, an
unfolding of the germs present in the old: the development of
magnitude into relational magnitude is decisive here. In order to
see on what that which is peculiar in the principles of differentiation
essentially rests, one should note the procedure with the
single functions, next the treatment of the composite ones,
and, as concerns the method in its totality, the serviceability
of Taylor's series as the fundamental formula of the systems of functions.

As concerns the single functions, the differential coefficient is determined by the presentation corresponding directly to the definition; thus:
\[ \frac{d(a+x)}{dx}=\lim\frac{a+(x+h)-(a+x)}{h}=\lim\frac{h}{h}=1;\quad\frac{d(ax)}{dx} \]
% the note 2) of p. 91 runs over onto the foot of this page: set in full at its mark
% --- p. 93 ---
\setcounter{footnote}{0}%
\opage{93}\[ =\lim\frac{a(x+h)-ax}{h}=\lim\frac{ah}{h}=a;\quad\frac{dx^{n}}{dx}=\lim\frac{(x+h)^{n}-x^{n}}{h} \]
\[ =nx^{n-1} \]
etc. What is peculiar is thus seen from the mutual
relation of the relational magnitudes.

In order to differentiate a composite function one expands
it in a series according to ascending powers of $h$, but since
this series properly receives its significance from the single functions
of which the composite consists, these must be
expanded in a series each by itself, from these in turn a single series corresponding to the composition
of the function is formed, and
the method of indeterminate coefficients is applied. Thus one finds
$\frac{d(y_{1}+y_{2})}{dx}$, since $Y=y_{1}+y_{2}$, i.e.\ $F(x)=f_{1}(x)+f_{2}(x)$, by expanding $F(x+h)=f_{1}(x+h)+f_{2}(x+h)$ in a series; thus:
\[ F(x+h)=F(x)+F'(x)\frac{h}{1}+\ldots. \]
\[ =f_{1}(x)+f_{2}(x)+f_{1}'(x)\frac{h}{1}+f_{2}'(x)\frac{h}{1}+\ldots. \]
; since now, in accordance with the
method of indeterminate coefficients, $F'(x)=f_{1}'(x)+f_{2}'(x)$,
we have $\frac{d.Y}{dx}=\frac{d(y_{1}+y_{2})}{dx}=\frac{dy_{1}}{dx}+\frac{dy_{2}}{dx}$. If the task is to
find $\frac{d(y_{1}\cdot y_{2})}{dx}$, then $F(x+h)=f_{1}(x+h)f_{2}(x+h)$
\[ =F(x)+F'(x)\frac{h}{1}+\ldots=(f_{1}(x)+f_{1}'(x)\frac{h}{1}+\ldots)(f_{2}(x)+f_{2}'(x)\frac{n}{1}+\ldots)=f_{1}(x)\cdot f_{2}(x)+(f_{1}(x)\cdot f_{2}'(x)+f_{2}(x)f_{1}'(x))\frac{n}{1} \] % print: n/1 for h/1, twice (misprint)
$+\ldots$, thus: $F'(x)=f_{1}(x)f_{2}'(x)+f_{2}(x)f_{1}'(x)$, that is to say:
\[ \frac{d(y_{1}y_{2})}{dx}=y_{1}\frac{dy_{2}}{dx}+y_{2}\frac{dy_{1}}{dx}, \]
etc.

In this context it can naturally not be
the task to enumerate the individual ``rules of differentiation''; the development of concepts
turns essentially on the conditions
for the application of Taylor's series. In order to obtain the single
% --- p. 94 ---
\setcounter{footnote}{0}%
\opage{94}differential coefficient $f'(x)$ alone, one must after all let the remaining
terms of the series drop out. This dropping out of the terms
does indeed receive an exact justification by assuming, in the equation:

\[ \lim\frac{f(x+h)-f(x)}{h}=f'(x)+f''(x)\frac{h}{1.2}+f'''(x)\frac{h^{2}}{1.2.3}+\ldots \]

$h=0$; but by making $h$ into $0$, one makes at the same time $f'(x)$ into $\frac{0}{0}$, and the difficulties
pointed out above return. If, on the other hand, one regards $h$ not as $0$, but as $dx$, then one gets, for $f(x+h)=y+k$, indeed a real differential coefficient $f'(x)=\frac{dy}{dx}$ in place of a designation contradicting % print: Differentialcoesficient (misprint)
itself; but then there seems again to belong to
the sum of the remaining terms of the series such a significance
that they, when an exact expression is demanded, cannot
be thrown away.

This difficulty is removed by closer consideration of
the series in question. To $f(x)=(x-a)^{n}$ there corresponds $(x+h-a)^{n}=((x-a)+h)^{n}=(x-a)^{n}+n(x-a)^{n-1}\frac{h}{1}+n(n-1)(x-a)^{n-2}\frac{h^{2}}{1.2}\ldots.+n(n-1)\ldots(n-r+1)(x-a)^{n-r}\frac{h^{r}}{[r]}\ldots..$

If now $1)$ $n$ is a positive integer, then the series terminates with $h^{n}$, and, so far as we still regard it as Taylor's formula, $\frac{d^{n+p}(x-a)^{n}}{dx^{n+p}}=0$, since $\frac{d^{n}(x-a)^{n}}{dx^{n}}=c$. % print: heeel (misprint)
If $2)$ $n$ is negative, it is seen that all the differential coefficients
must contain $(x-a)$ in the denominator, and
thus, when $x=a$, $\frac{d^{n}y}{dx^{n}}=\infty$.
If $3)$ $n$ is fractional, e.g.\ $n+1>r>n$, then it is clear that $\frac{d^{r}(x-a)^{n}}{dx^{r}}=n(n-1)$
% --- p. 95 ---
\setcounter{footnote}{0}%
\opage{95}\[ \ldots.(n-r+1)(x-a)^{n-r}=\frac{n(n-1)\ldots.(n-r+1)}{(x-a)^{r-n}}, \]

and that thus $\frac{d^{r}(x-a)^{n}}{dx^{r}}$ must contain $x-a$ in the denominator.
The differential coefficients of higher than the $r$th order must then thus, for $x=a$, all become $=\infty$.
The result now is
this: Taylor's formula, although generally valid, can for
special values of $x$ become unusable; but, \emph{when} the series is usable, it is, for $\mathrm{Lim}\frac{f^{n+1}(x)}{(n+1)f^{n}(x)}h<1$,
always convergent, and the more convergent, the
smaller $h$ is, in other words: the smaller $h$ is, the
fewer terms are needed to attain a certain degree of
\emph{accuracy.}

That there is nevertheless still something unsatisfactory in this conception
Hegel has pointed out, and not without reason, in the
following way: „Lagrange hat bekanntlich die ursprüngliche
Methode Newtons, die Methode der Reihen, wieder aufgenommen,
um die Schwierigkeiten, welche die Vorstellung
des Uendlich-Kleinen, so wie derjenigen, welche die Methode
der ersten und letzten Verhältnisse und Grenzen mit sich
führt, überhoben zu seyn. Es ist von seinem Functionen-Kalkul,
dessen sonstige Vorzüge in Rücksicht auf Präcision,
Abstraktion und Allgemeinheit anerkannt genug sind, als
hierher gehörig nur dieß anzuführen, daß er auf dem Fundamentalsatze
beruht, daß die Differenz ohne daß sie Null
werde, so klein angenommen werden könne, daß
jedes Glied der Reihe die Summe aller folgenden
an Größe übertreffe. --- Es wird auch in dieser Methode
von den Kategorien vom Zuwachs und von der Differenz
der Function angefangen, deren veränderliche Größe
den Zuwachs erhalte, womit die lästige Reihe hereinkommt,
von der ursprünglichen Function; so wie im Verfolg die
wegzulassenden Glieder der Reihe nur in der Rücksicht, daß
sie eine Summe konstituiren, in Betracht kommen, und der
% --- p. 96 ---
\setcounter{footnote}{0}%
\opage{96}Grund sie wegzulassen, in das Relative ihres Quantums
gesezt wird.“\footnote[1]{Wissensch. der Logik. 1ster Theil. p. 316—317.}

In order to remove this unsatisfactoriness we need not, however, % print: di (misprint)
with Hegel demand so qualitative a closure that we
thereby at bottom give up quantity; in order to convince
oneself that an absolute exactness can be attained, it suffices
to make clear to oneself the qualitative change the series
undergoes when $h$ is transformed into $dx$. In $\frac{f(x+dx)-f(x)}{dx}$
\[ =f'(x)+f''(x)\frac{dx}{1}+f'''(x)\frac{dx^{2}}{1.2}+\ldots. \]
it must be expressly emphasized from the standpoint of
mathematics that $f''(x)dx$ is not merely much smaller, but infinitely many times smaller than $f'(x)$, that they are, in other words, ``incomparable''.
Since now the same repeats itself in the juxtaposition of $f'''(x)\frac{dx^{2}}{1.2}$ with $f''(x)dx$ etc., it is seen that the terms
of the series drop out, not because they are so insignificant that the inexactness
committed by their omission is very slight;
but because, as incomparable terms, as terms which,
according to their nature, belong to another
sphere, they must in the concept drop out.

In some cases, where it really looks as if
one committed an error by omitting a term of higher order,
it turns out that the error is compensated, so that there is nevertheless % print: sva (misprint)
no error committed. Such a ``compensation of error''
can be seen e.g.\ in the following case. For the ordinary
parabola: $y^{2}=px$ the subtangent is, according to the method of differentiation,
determined by $subtg=y\frac{dx}{dy}$; $(y+dy)^{2}=p(x+dx)$ and $dy=\frac{pdx}{2y}-\frac{dy^{2}}{2y}$; this gives, when $\frac{dy^{2}}{2y}$ is thrown away,
% --- p. 97 ---
\setcounter{footnote}{0}%
\opage{97}\[ \frac{dx}{dy}=\frac{2y}{p}. \]
One thus obtains: $subtg.=y\frac{dx}{dy}=\frac{2y^{2}}{p}$, or, since $y^{2}=px$, $subtg.=2x$. By this procedure it certainly seems to be an error that the magnitude $\frac{dy^{2}}{2y}$ was omitted; but by letting the differential of the ordinate of the parabola
at the consecutive point coincide with the differential
of the ordinate of the tangent, one has, as far as the parabola
is concerned, made $dy$ too large from the beginning. According to
a theorem set up by Apollonius, and thus independent of the differential calculus, the subtangent is $=2x$, and, denoting by $u$ the difference between the differential of the ordinate of the tangent and that of the parabola, we have $\frac{2x}{y}=\frac{dx}{dy+u}$, whence
it is seen that $u=\frac{dy^{2}}{2y}$, and that thus the piece that was to be sub-tracted from $dy$ corresponds exactly to the term that, in accordance with
the principle of differentiation, was thrown away\footnote[1]{Weissenborn. Work cited, p. 156—58.}

Understood essentially, in the equation $dy=\frac{pdx}{2y}-\frac{dy^{2}}{2y}$ the last term is not thrown away because $dy^{2}$ is assumed to be
a still smaller part of a line than $dx$ or $dy$, but
because $dy^{2}$ is not homogeneous with a line, but with a line's
limit, a point, a vanishing moment. If
the throwing away of a term means that this, in order to be thrown away,
must have been transformed from abiding magnitude into vanishing moment, then the exact procedure, seen
under reflection, must surely seek its last justifying ground in
a dialectic.
% --- p. 98 ---
\setcounter{footnote}{0}%
\opage{98}
\begin{center}c) The integral: the element and the finite magnitude.\end{center}

What the differential really is can be fully elucidated only
by its opposite, namely the integral. ``In opposition
to the differential as the negative the integral is the
positive. If the differential is a point, then the integral is a
line; if the differential is a line, then the integral is a surface,
if the differential is a surface, then the integral is a solid.
If the differential signifies a moment in the infinite discretion of the line,
then the integral signifies the resolution of the discrete moments
into the finite continuum; if the differential designates the
infinitely small term, then the integral designates the sum of such
terms in infinitely great number;''\footnote[1]{Philos. og Mathem. p. 94, p. 60.} if the differential is determined by
the derived function $f'(x)dx$, then the integral is in turn determined by the original $f(x)$.

Since the infinite summation
is gradually seen to coincide with the task
of leading a derived function back to the corresponding original one\footnote[2]{By defining the integral so that $y=\int f'(x)\,dx$ is the magnitude which, when differentiated with respect to x, gives $f'(x)$, one undeniably obtains by the integration \emph{quite exactly} the sought function, namely $f(x)$; but this exactness, resting on the definition, stands and falls with the question whether the definition itself is adequate, i.e.\ whether one is now really entitled to set $\frac{dy}{dx}=\lim\frac{f(x+h)-f(x)}{h}=f'(x)$, or whether the series for $\frac{dy}{dx}$ is not really after all: $f'(x)+f''(x)\frac{dx}{1.2}+f'''(x)\frac{dx^{2}}{1.2.3}+\ldots$ and whether one has not then committed an inexactness by throwing away from this series all the terms following $f'(x)$. By a given definition, an agreed designation, a confusion of thought can be forestalled; but a problem is not thereby yet solved.}
the development of the concept of the integral will at the same time show
itself to coincide with the solution of the problem of infinity.
% --- p. 99 ---
\setcounter{footnote}{0}%
\opage{99}
\begin{center}α) Integrating summation: moment and element.\end{center}

The mathematical proposition that every finite magnitude
can be assumed to consist of infinitely many infinitely
small parts is in its generality certainly trivial, but
nevertheless always just as significant as the logical principle
that every thing is what it is. The trivial lies
not in the abstraction, but in the empty abstraction
being held fast as something existent, whereby the dialectical is overlooked,
and the development of the concept comes to a halt. As a given whole
the finite magnitude naturally consists of parts; if now
this whole is grasped from the side of unity, continuity is seen; if it is grasped
from the side of the parts, discretion is seen. The doubling of reflection,
according to which the essential whole has its subsistence
in the subsisting parts, magnitude as pure
quantum, however, cannot present; that which is immediate in the category of quantity
is still so predominant that the one of the
sides of reflection vanishes when the other is posited: when
discretion is posited, continuity is broken; when continuity is
posited, discretion vanishes. The separation of discretion
and continuity has in turn occasioned a separation of discrete
and continuous magnitudes, as if they constituted two
species of magnitudes subsisting for themselves\footnote[1]{``Magnitudes can be either discrete (separated) or continuous (connected). The former are mere collections of any homogeneous individual items separated from one another (e.g.\ a human crowd, a herd of animals, a pile of stones), so that the degree or value of the magnitude must be stated by a number. The latter are wholes in which no separation is found between any homogeneous parts (e.g.\ a length, a weight, a period of time). Two homogeneous continuous magnitudes can always be compared with regard to their values. The simplest case is that where the one magnitude is contained exactly a whole number of times, $p$, in the other, where $p$ may be 1, 2, 3... If the first is taken as unit, the other can be said to be equal to the number p''... Ramus. Elementær Geometrie, Kbhvn 1850; Introduction.} That such a
separation is merely provisional is revealed, however, by the fact that a
continuous magnitude can be determined only by positing
another homogeneous with it as unit, and now investigating
how many times this is contained in that; but by such an
investigation the continuous magnitude is precisely
seen under the point of view of discretion. The separation of
discrete and continuous magnitude reduces in the concept to
a separation of finite and infinite discretion.\footnote[1]{The separation of discrete and continuous magnitude has thus here been explained in such a way that in the explanation it cancels itself. A collection of 20 people is a crowd, a collection of 19, of 18, of 17, of 16\dots{}. people likewise: the magnitude: crowd is, thus conceived, a continuous magnitude. If the mark of a discrete magnitude is to be this, that its value ``must be stated by a number'', then such magnitudes as: ``a length, a weight, a period of time'' are sheer discrete magnitudes; for an ell is 4 quarters, a pound is 32 loths, a minute is 60 seconds. As every part of a length is itself a length, so also is every part of the crowd, when we keep to the pure concept of magnitude, and do not intermingle extraneous representations, itself a crowd.} With the % the print marks this note 1) too, and sets it as a second paragraph of the note above
% --- p. 100 ---
\setcounter{footnote}{0}%
\opage{100}finite discretion continuity is broken, with the
infinite discretion the broken continuity returns
restored; in finite discretion the parts, as existing
terms, are still particular magnitudes\footnote[1]{The transition from difference- to differential coefficient is seen, as regards Taylor's series, thus: let $y_{0}, y_{1}, y_{2} \ldots y_{n}$ be $f(x), f(x+h), f(x+2h), \ldots f(x+nh)$; one then has: $y_{1}=y_{0}+\Delta y_{0}$, $y_{2}=y_{0}+2\Delta y_{0}+\Delta^{2}y_{0}$, $y_{3}=y_{0}+3\Delta y_{0}+3\Delta^{2}y_{0}+\Delta^{3}y_{0}, \ldots$ $y_{n}=y_{0}+\frac{n}{1}\Delta y_{0}+\frac{n(n-1)}{1.2}\Delta^{2}y_{0}+\frac{n(n-1)(n-2)}{1.2.3.}\Delta^{3}y_{0}$ $+\ldots$. thus $f(x+nh)=f(x)+\frac{n}{1}\Delta.f(x)+\frac{n(n-1)}{1.2}\Delta^{2}.f(x)$ $+\ldots.$; if $nh=i$ is set, then $n=\frac{i}{h}$. $f(x+i)=f(x)+\frac{i}{1}\Delta\frac{f(x)}{h}+\frac{i(i-h)}{1.2}\Delta^{2}\frac{f(x)}{h^{2}}+\ldots$; if now $h$ is transformed into $dx$, one obtains $f(x+i)=f(x)+\frac{d.f(x)}{dx}\frac{i}{1}+\frac{d^{2}f(x)}{dx^{2}}\frac{i^{2}}{1.2}+\ldots=f(x)+f'(x)\frac{i}{1}+f''(x)\frac{i^{2}}{1.2}+\ldots.$\par The relation between the differential and the integral calculus is already prefigured in the analysis of the finite by the opposition between the direct and the inverse calculus of differences. The calculus of differences consists in ``finding the differences of given functions and then the differences of their differences. Thus $\Delta.x^{n}=x^{n+1}-x^{n}=x^{n}(x-1)$. From this it is seen that if one takes the difference of the terms in a quotient series, a new quotient series is obtained. To $1, x, x^{2}, x^{3}, x^{4} \ldots.$ there correspond $x-1, x(x-1), x^{2}(x-1), x^{3}(x-1)\ldots.$ and $(x-1)^{2}, x(x-1)^{2}, x^{2}(x-1)^{2}$ is in turn formed from the differences in the preceding series.'' A. Steen. Reen Mathem. p. 338.\par ``The calculation by which one finds a function from its difference is called the inverse calculus of differences, and in opposition thereto the proper calculus of differences is called direct. The inverse is also called integration (integrate in opposition to differentiate). From $\Delta.s_{n}=u_{n}$ one gets the integral $s_{n}=\Sigma.u_{n}$, from $\Delta.x^{n}=x^{n}(x-1)$ the integral ``$x^{n}=\Sigma.x^{n}(x-1)$; but since herein $x-1$ is a common factor for all terms under $\Sigma$, which is independent of $n$ (constant), it can be placed outside $\Sigma$, so that $x^{n}=(x-1)\Sigma.x^{n}$ or $\Sigma.x^{n}=\frac{x^{n}}{x-1}$. Work cited, p. 340.}, in infinite discretion
% the notes 1) of p. 99 run over onto the foot of this page: set in full at their marks
% --- p. 101 ---
\setcounter{footnote}{0}%
\opage{101}the parts are vanishing terms, magnitude-moments, dialectical
intermediate determinations of Something and Nothing\footnote[1]{With this it agrees that the concept of function is conceived ``als reale Existenz des discret-continuirlichen Qvantums''. H. Schwarz. Work cited, p. 22.} As moment
the part is in an infinite transition to Nothing, without however
becoming Nothing; as element the part is in an infinite transition % print: Delene (misprint)
to Something, a Something from which the finite is produced.
To the finite discretion, in which the parts maintain themselves as
existing terms, there corresponds a finite summation: the whole
results from a counting together of the parts; to the infinite
discretion, in which the parts are moments, there corresponds an
infinite summation, a transformation of the moment into
% the note 1) of p. 100 runs over onto the foot of this page: set in full at its mark
% --- p. 102 ---
\setcounter{footnote}{0}%
\opage{102}element, an intellectual act which determines the finite
as product of the infinite. The relation between the element
and the finite magnitude can then for the first be elucidated
by summation.

If the magnitude $x$ is divided in such a way that each part $\Delta x=\frac{x}{4}$,
the whole magnitude becomes:
\[ 4\frac{x}{4}=\Sigma\Delta x=x \]
; if, on the other hand, $x$ is thought
divided in such a way that the single part is not a finite
magnitude, but a moment transformed into element, $x$ as finite magnitude cannot arise by addition: $dx+dx+\ldots.$, but only according to the concept in that $x$ is conceived
as an infinite sum of $dx$, thus as $\int_{0}^{x}dx=x$. That the moment of discretion $dx$ is not $0$, but an element, is now seen
directly from the fact that an infinite sum of zeros cannot give
anything but $0$; but if dx is not absolute zero, then $\int_{0}^{x}$ is likewise
not absolute $\infty$. If $\int_{0}^{x}dx=x$ is denoted by $\infty .0$, it must be
expressly emphasized that $\infty$ corresponds to $\frac{x}{dx}$, as $0$ corresponds to $dx$. The difference between $\Sigma_{0}^{x}\Delta x$ and $\int_{0}^{x}dx$ is by no means this,
that only the first expression should indicate something determinate,
the last, on the other hand, dissolve into empty formalism; but
the difference is that $x$ is conceived in the first expression as discrete,
in the last expression, on the other hand, as a continuous magnitude.

This difference assumes a still more articulated form,
when we compare the sums\footnote[1]{By $\Sigma_{0}^{x}f'(x)\Delta x$ we understand here the sum of all the values that $f'(x)\Delta x$ runs through, when x runs through all values between 0 and x with the difference $\Delta x$.} $\Sigma_{0}^{x}f'(x)\Delta x$ and $\int_{0}^{x}f'^{(x)}dx$.
% --- p. 103 ---
\setcounter{footnote}{0}%
\opage{103}In the finite sum $\Sigma_{0}^{x}x\Delta x=(\Delta x+2\Delta x+3\Delta x+\ldots n\Delta x)\Delta x=(1+2+3+\ldots n)(\Delta x)^{2}$, the difference $\Delta x=2\Delta x-\Delta x=3\Delta x-2\Delta x\ldots.$ shows the finite discretion,
a leap, a finite distance between the finite
terms; in the infinite sum $\int_{0}^{x}xdx=(dx+2dx+3dx+\ldots ndx)dx=(1+2+3+\ldots n)dx^{2}$, one sees from the differential $dx=2dx-dx=3dx-2dx\ldots$, that the discretion
is infinite, that here thus, in other words, from $0$ to $x$ there is a
continuous transition. In the equation $\Delta x=\frac{x}{n}$ $n$ is finite; in the equation $dx=\frac{x}{n}$ $n$ is on the other hand infinite; as a consequence
there results then
\[ \Sigma_{0}^{x}x\Delta x=\frac{n(n+1)}{2}\frac{x^{2}}{n^{2}}=\frac{x^{2}}{2}+\frac{x^{2}}{2n} \]
, whereas
\[ \int_{0}^{x}x\,dx \]
on the other hand gives
\[ \mathrm{Lim}\frac{(1+\frac{1}{n})}{2}x^{2}=\frac{x^{2}}{2} \]
. Whereas
\[ \Sigma_{0}^{x}x^{2}\Delta x=[\Delta x^{2}+(2\Delta x)^{2}+(3\Delta x)^{3}+\ldots(n\Delta x)^{2}]\Delta x= \] % print: (3\Delta x)^{3} (misprint)
\[ (1+2^{2}+3^{2}+\ldots n^{2})(\Delta x)^{3}=\frac{n(n+1)(2n+1)}{1.2.3}(\Delta x)^{3} \]
\[ =\frac{n(n+1)(2n+1)}{1.2.3}\frac{x^{3}}{n^{3}} \]
gives
\[ \frac{x^{3}}{3}+\frac{x^{3}}{2n}+\frac{x^{3}}{6n^{2}} \]
on account
of the discretion: gives
\[ \int_{0}^{x}x^{2}dx=[(dx)^{2}+(2dx)^{2}+\ldots.(ndx)^{2}]dx=(1+2^{2}+3^{2}+\ldots n^{2})dx^{3}=\frac{n(n+1)(2n+1)}{1.2.3} \]
\[ \frac{x^{3}}{n}=\mathrm{Lim}\frac{(1+\frac{1}{n})(2+\frac{1}{n})x^{3}}{1.2.3} \] % print: \frac{x^{3}}{n} (misprint)
on the other hand $\frac{x^{3}}{3}$ on account of
the continuity. To
\[ \Sigma_{0}^{x}x^{3}\Delta x=[\Delta x^{3}+(2\Delta x)^{3}+\ldots(n\Delta x)^{3}]\Delta x=(1+2^{3}+\ldots n^{3})(\Delta x)^{4}=(1+2+3+\ldots n)^{2}(\Delta x)^{4} \]
there corresponds
\[ \frac{n^{2}(1+n)^{2}}{2^{2}}\frac{x^{4}}{n^{4}}=\frac{x^{4}}{4}+\frac{x^{4}}{2n}+\frac{x^{4}}{4n^{2}} \]
, since $n$ is
% --- p. 104 ---
\setcounter{footnote}{0}%
\opage{104}finite; to
\[ \int_{0}^{x}x^{3}dx=[dx^{3}+(2dx^{3}+\ldots(ndx)^{3}]dx=(1+2^{3}+\ldots.n^{3})dx^{4}=(1+2+3+\ldots.n)^{2}dx^{4} \]

there corresponds, on the other hand,
\[ \frac{n^{2}(1+n)^{2}}{2^{2}}\frac{x^{4}}{n^{4}}=\mathrm{Lim}\frac{(\frac{1}{n}+1)^{2}}{4}x^{4}=\frac{x^{4}}{4} \]
; etc.

In the infinite summation there is made clear, what has already been
pointed out earlier (II., $\gamma$), the dialectical unity of the finite and of the
infinite. The proposition that a magnitude consists of an infinite number of infinitely small parts
has then in it a disposition to articulating determinations that is not merely just as rich, but even richer,
than the proposition that
the finite magnitude consists of a finite number of finite
parts. The elements $dx, xdx, x^{2}dx$ etc.\ are in relation to
the magnitudes concerned $x, \frac{x^{2}}{2}, \frac{x^{3}}{3}\ldots$ just as peculiarly
determined as any finite part $\Delta x, n\Delta x\ldots$, in relation to the corresponding finite whole. As long as
the understanding seeks all difference and all determinateness in the finite,
it can indeed treat magnitudes as terms in a comprehensive
operational mechanics, but it is not able to
grasp what magnitude really is, not able to cognize
the essence of magnitude. For a cognition of the essence of magnitude
there is required a cognition of its origin, its becoming % print: tens (misprint)
and vanishing, its dissolution into its moment, its
progression from its element.

\begin{center}β) Principles of integration: the relation of derived functions to the original ones.\end{center}

The infinite summation cannot be directly carried out;
in order to furnish analogies, like the above, the
finite summation must already undergo so essential a transformation
that it no longer rests on simple addition. The
% --- p. 105 ---
\setcounter{footnote}{0}%
\opage{105}infinite summation is determined by the relation between element
and magnitude; but this relation is determined in turn, since the
complete concept of magnitude coincides with the concept of function,
by the relation between the derived function and
the corresponding original one. As element the differential in question
is: $d f(x)=dy$ a product of the derived
function, the differential coefficient $f'(x)$, and the moment, $dx=dx$, by which the independent varies. Inasmuch as
magnitude is a function of itself and unity: $f(x)=x$, and $dx=df(x)=f'(x)dx=dy$, there is here between moment
and element only the formal difference that dx as an expression of differentiation
is moment, as the real possibility of integration,
on the other hand, element. In essentially determinate functions
the difference has essential significance; thus in $f(x+dx)=(x+dx)^{2}=y+dy$ the differential dx is only moment, the differential $dy=df(x)=f'(x)dx$ on the other hand the proper
element.

From the difference between moment and element integration
first receives in the concept an all-sided elucidation. If $\int_{0}^{x}df(x)=\int_{0}^{x}f'(x)dx$ is conceived as the sum of the infinitely many
terms that arise as $x$ in $f'(x)dx$ continuously runs through
the values from $0$ to $x$, then this sum will, according to
the nature of $f'(x)dx$, i.e.\ according to the nature of the element,
give, for $f(x)=x$, the sum: $dx+dx+\ldots=x$; for $f(x)=x^{2}$, the sum:
\[ \int_{0}^{x}d(x^{2})=\int_{0}^{x}f'(x)dx=\int_{0}^{x}2xdx \]

\[ =2\int_{0}^{x}xdx=2(1+2+3+\ldots.n)dx^{2}=\mathrm{Lim}\,2\frac{x^{2}}{n^{2}} \]
\[ (1+2+3+\ldots.n)=\mathrm{Lim}\,2\frac{(n+1)}{1.2}x^{2}=\mathrm{Lim}(1+\frac{1}{n}) \]
\[ x^{2}=x^{2} \]
, etc. If it is now evident from this how the
infinite summation transforms itself into integration, inasmuch as
integration, by the emphasis falling on the element $f'(x)dx$, and consequently on $f'(x)$, solves its task in principle
% --- p. 106 ---
\setcounter{footnote}{0}%
\opage{106}of going from the differential back to the integral, from the derived
function back to the original one: then it will now
also be able to show itself in what the difference between finite and
infinite summation is essentially grounded.

If geometric significance is attached to the expression $(x+\Delta x)^{2}=x^{2}+2x\Delta x+(\Delta x)^{2}$, it is seen that all the terms are
comparable, inasmuch as they all denote areas; but if
$\Delta x$ is now changed to $dx$, then $(x+dx)^{2}=x^{2}+2xdx+dx^{2}$ gives

incomparable terms, since $x^{2}$ is an area, $xdx$ the element of the area and $dx^{2}$ the moment of the element. As regards the solid
form, $(x+\Delta x)^{3}=x^{3}+3x^{2}\Delta x+3x(\Delta x)^{2}+(\Delta x)^{3}$ will give comparable terms, namely sheer solids; whereas $(x+dx)^{3}=x^{3}+3x^{2}dx+3xdx^{2}+dx^{3}$ on the other hand gives incomparable
terms, since $x^{3}$ is after all a solid, $x^{2}dx$ a surface, $xdx^{2}$ a line
and dx a point. If now the task is set to determine
the magnitude of the increment that an area, $a^{2}$, receives, by summing
the increments which this area obtains, as $x$ from the
finite value $a$, by repetition of the increment $\Delta x$, has reached the value $b$; then it lies in the nature of the case that the
finite summation must in each addend have all the increment's
terms, so that the sum becomes $\sum_{x=a}^{x=b}[2x\Delta x+(\Delta x)^{2}]$, not $\sum_{x=a}^{x=b}2x\Delta x$.

If, on the other hand, the task is to sum the increments which the area
$a^{2}$ receives, as $x=a$ by increments of $dx=dx$ has grown
to $x=b$; then the infinite summation must exclusively
keep to the element of the area, so that the sum becomes:
\[ \int_{x=a}^{x=b}2xdx \]
, not
\[ \int_{x=a}^{x=b}(2xdx+dx^{2}) \]
.\footnote[1]{$\int_{a}^{b}(2x\,dx+dx^{2})=2a\,n\,dx+n^{2}dx^{2}$, where $n=\frac{b-a}{dx}$, thus $\int_{a}^{b}(2xdx+dx^{2})=2a(b-a)+(b-a)^{2}=b^{2}-a^{2}$. Whereas $\int_{a}^{b}2xdx=2a n\,dx+n^{2}dx^{2}-n\,dx^{2}=b^{2}-a^{2}-(b-a)dx$. If we were now to see this difference under the point of view of ``compensation of errors'' (cf.\ III., b, $\gamma$); then we have, by omitting $dx^{2}$ from dy, really committed an error, but this error is then in turn compensated by our omitting in $\int_{a}^{b}2xdx=b^{2}-a^{2}$ the negative term ÷ (b — a) dx. The essential reason why (b — a) dx is omitted is, however, this, that here by this expression only a line is denoted, not an area.} That the moment of the element
% --- p. 107 ---
\setcounter{footnote}{0}%
\opage{107}drops out is after all directly grounded in the fact that a moment
which is already contained in the element itself cannot acquire
special significance.\footnote[1]{How a thinker such as G. Berkeley, e.g., has sought a determination of concept in which the relation between moment, element and finite magnitude could be held fast is seen among other things from the following: Because there is no Number of Parts so great, but it is possible there may be a Line containing more, the Inchline is said to contain Parts more than any assignable Number; which is true, not of the \emph{Inch} taken absolutely, but only for the Things signified by it. But Men not retaining that Distinction in their Thoughts, slide into a belief that the small particular Line described on Paper contains in it self Parts innumerable. There is no such thing as the ten-thousandth Part of an Inch; but there is of a \emph{Mile} or \emph{Diamether} of the \emph{Earth}, which may be signified by that Inch. When therefore I delineate a Triangle on Paper, and take one side not above an Inch, for Example, in length to be the \emph{Radius}: This I consider as divided into ten thousand or an hundred thousand Parts, or more. For though the ten-thousandth Part of that Line considered in it self, is nothing at all, and consequently may be neglected without any Error or Inconveniency; yet these described Lines being only Marks standing for greater Quantities, whereof it may be the ten-thousandth Part is very considerable, it follows, that to prevent notable Errors in Practice, the \emph{Radius} must be taken of ten thousand Parts, or more. Principles of human knowledge. Lond. 1734. p. 147—148.}
% the note 1) of p. 106 runs over onto the foot of this page: set in full at its mark

% --- p. 108 ---
\setcounter{footnote}{0}%
\opage{108}As the element of the area, $xdx$ is a dialectical intermediate determination
of line and surface; as the moment of the element, $dx^{2}$ is a dialectical intermediate determination of the rudiment of a line
and the rudiment of a surface. In the geometrical conception of $x^{3}$, to which corresponds $\int_{0}^{x}3x^{3}dx$, not $\int_{0}^{x}(3x^{2}dx+3xdx^{2}+dx^{3})$, % print: 3x^{3}dx (misprint)
one sees, as regards $(x+dx)^{3}$, not merely in $x^{3}$ a body, in
$x^{2}dx$ a rudiment of a body: the proper surface, in $xdx^{2}$
a rudiment of a corporeal rudiment: the corporeal line, but
even in $dx^{3}$ a rudiment of the rudiment of the corporeal rudiment: the
corporeal point.

The old dialectical dispute concerning the mutual relation of the point, the line,
the surface and the body: whether
it is namely the point that is the original, so that the point
generates the line, the line the surface, the surface the body; or
whether perhaps conversely it is the body that is the original, so
that the surface is derived from the body by abstraction, the line from
the surface, the point from the line\footnote[1]{In his Metaphysics (Metaphys. III. 5) Aristotle has treated the question whether numbers, bodies, surfaces and points are substances (οὐσίαι τινές), or not. ``In comparison with the things to which one would most readily ascribe substantiality, such as water, earth and air, of which the composite bodies consist, heat and cold and similar qualities (πάθη, affections) are surely not substances; only the body, which has these affections, remains as a being, a substance (ὡς ὄν τι καὶ οὐσία τις οὖσα). And yet the body seems to have less substantiality than the surface, this less than the line, and the line less than the unit and the point. For by these the body is limited, and it seems that these might well exist without the body, but the body, on the contrary, not without them. If one grants, however, that the lines and points have more substantiality than the bodies, one falls into perplexity; for these determinations of limit rest, after all, evidently on divisions of the body (διαιρέσεις ὄντα τοῦ σώματος), partly in breadth, partly in depth, partly in length. Moreover, either every form must be contained in the body, or none: if, therefore, Hermes is not in the stone, and the half of the cube is not found delimited in the cube itself, then the surface too will not be found in it; for if any surface were found in the cube, then the one by which the cube is halved would have to be found there too. The same holds of the line, the point and the unit. To the absurdities mentioned (ἄλογα) there is added yet the one that concerns coming-to-be and passing-away; for at one time two surfaces become one, namely when two bodies touch each other, at another time, when a body is divided, there are again two surfaces where before there was none: an existing surface can thus vanish, and one not previously existing come forth, according as the bodies are united or separated. It seems then, accordingly, as though points, lines and surfaces, if they were substances, would have to alternate between existing and not existing; but points, lines and surfaces can neither come to be nor pass away, although at one time they are, at another not.'' This whole aporia, in which Aristotle does not so much clarify the matter as put it on the rack, has its nodal point here not in the relation of the intellectual to the material, but in the question: what is the original, that which is valid in itself, the abstract or the concrete?}, --- this problem finds
its solution only when the principles of integration are clarified on all sides.

% --- p. 109 ---
\setcounter{footnote}{0}%
\opage{109}That the mathematical point cannot by any motion
generate a line, or the mathematical line a
surface, or the mathematical surface a body, is seen from this:
that the element of the line, $dx$, in order to correspond to the mathematical
point, must be $0$; but when $dx=0$ cannot be integrated
into a line, then neither can $xdx=x.0=0$ be integrated into a
surface, nor $x^{2}dx=x^{2}.0=0$ into a body.
In opposition to the mathematical point, therefore, the
finite bodily form is the original. But this bodily form
must nevertheless, when we look at what is original in the essential
sense, be conceived as the derived,
for everything that rests on becoming and coming-to-be
is indeed a derived. The finite bodily form has, as a derived, $x^{3}$, a coming-to-be, in that it arises by integration
of the corporeal surface, $x^{2}dx$, this by integration
% the note 1) of p. 108 runs over onto the foot of this page: set in full at its mark
% --- p. 110 ---
\setcounter{footnote}{0}%
\opage{110}of the corporeal line, $xdx^{2}$, this by integration of the
corporeal point, $dx^{3}$. This corporeal point, the element of the elements, is thus here the not-derived, that which is original in itself,
the principle of corporeality, not an empty abstraction,
but a rudiment of totality\footnote[1]{When $dx^{3}$, held fast in pure mathematical form, nevertheless presents itself as an empty abstraction, in which one cannot seek a principle, the reason is this: that the point so determined, although it is the principial rudiment of extensive magnitude, cannot acquire reality without at the same time signifying a rudiment of intensive magnitude; but the doctrine of the intensive, inasmuch as intensive magnitude is always more than magnitude, falls outside the abstractly mathematical.}, the point that \textit{in concreto}
corresponds to its concept.

As regards the particular functions in general,
one can now recognize the correctness of the integration by this: that
by differentiating the integral found one obtains the
original differential, just as, as regards subtraction,
one can recognize the difference by this: that added to the
subtrahend it gives the minuend. Thus, e.g.,
\[ \int x^{m}dx=\frac{x^{m+1}}{m+1}+c \]
since
\[ d\left(\frac{x^{m+1}}{m+1}+c\right)=x^{m}dx \]
; here in the
different principles there is essentially only one principle. But whereas
from a given original function one can always find
the derived one, it is on the other hand not always possible, when a
function is given as a differential coefficient, to find
a corresponding original one. To be sure, one can, to give
an example, transform $\int\frac{x^{m}dx}{l(x)^{n}}$ into
\[ -\frac{x^{m+1}}{(n-1)(lx)^{n-1}}+\frac{m+1}{n-1}\int\frac{x^{m}dx}{lx^{n-1}}\ldots \]
, but by continued application of this
formula one stops at last at $\int\frac{x^{m}dx}{l.x}$, which in general
% --- p. 111 ---
\setcounter{footnote}{0}%
\opage{111}($m \lessgtr -1$) cannot be integrated, because there is
no known function whose differential coefficient is $\frac{x^{m}}{l.x}$\footnote[1]{Cf. Jürgensen. Begyndelsesgr. p. 75—77.}.

\begin{center}γ) The definite integral: the solution of the problem of infinity.\end{center}

Although one cannot integrate every function
in finite form, one can nevertheless find a differential for every function,
and consequently regard every finite magnitude
as an integral. For since the integral is a function of
the variables contained in it, $\int_{0}^{x}\Phi(x)dx=f(x)$, one can in general
put: $f_{1}(x)=2, f_{2}(x)=3, f_{3}(x)=4\ldots$, and
determine $x$ in such a way that each of these equations, even if
one can give the numerical value of $x$ only by approximation, is
satisfied, and in this way obtain: $\int_{0}^{x}\Phi_{1}(x)dx=2$,
\[ \int_{0}^{x}\Phi_{2}(x)dx=3, \int_{0}^{x}\Phi_{3}(x)dx=4\ldots \]

The articulating reciprocal relation of the infinite with the finite
shines forth from the manner in which the definite integral
is stated; for although the finite magnitude, as regards its value,
is equal to itself, by whatever way we may arrive at
it, yet it makes an infinite difference in the concept whether
such a magnitude is had as immediately given, or
arises by integration.

In integration the magnitude has its becoming; in order that
this becoming may end in being, the integral must
be ``taken between definite limits.'' The demand for definite limits is signified immediately by the fact that to
the general integral of the differential function there is added ``an arbitrary
constant.'' Inasmuch as such a constant is
% --- p. 112 ---
\setcounter{footnote}{0}%
\opage{112}simply added, one can, when it is only a matter of showing
what is permissible in this, indeed appeal to the circumstance
that, since $d.c=0$ when $c$ is constant, and
consequently $d(f(x)+c)=f'(x)dx$, one will, when the integration
is carried out, obtain $\int f'(x)dx=f(x)+c$, and thus
get the constant along; but the main point is that one see
not merely what is permissible, but what is necessary in
the adding of the constant.

In the analysis of the finite, where every magnitude, however
it may otherwise arise, can be taken by itself, is conceived
as existing, and inasmuch as it is ascribed an arbitrary
value, a special introduction of arbitrary constants
can play no principal role;\footnote[1]{In the calculus of differences, where the magnitude varying by finite intervals breaks the way for the opposition to the constant, the integral will already require, with regard to this, the addition of a constant. ``Every integral contains an indeterminateness, as long as it is not removed by closer conditions. This lies in the fact that $\Delta.(s_{n}+c)=\Delta.s_{n}=u_{n}$, c being an arbitrary constant, and so one can have $s_{n}+c=\Sigma.u_{n}$. To every integral there can therefore be added an arbitrary constant, whose determination depends on conditions that must be specified more closely. Therefore $\Sigma.x^{n}=\frac{x^{n}}{x-1}+c$. If by $\Sigma.x^{n}$ one wishes to understand the sum of the terms $1, x, x^{2}:..x^{n-1}$, then to $n=0$ there must correspond $\Sigma.x^{0}=0$, thus: $0=\frac{1}{x-1}+c$, whence $c=-\frac{1}{x-1}$, consequently $\Sigma.x^{n}=\frac{x^{n}-1}{x-1}$, which is well known, and agrees with the formula for the sum of the terms of a geometric series.'' A. Steen. Reen Mathem. p. 340—341. Just because the difference removes the constant, the integral must demand it back; just because $s_{n+1}+c-(s_{n}+c)=s_{n+1}-s_{n}$, $\Sigma u_{n}=\Sigma\Delta.s_{n}$ must expressly be put, not as $s_{n}$, but as $s_{n}+c$; for c, which, since $c-c=0$, does not concern $\Delta.s_{n}$, concerns, on the other hand, since $s_{n}+c$, taken by itself, also depends on c, the integral corresponding to $\Delta.s_{n}$.} in the analysis of the infinite
the arbitrary constant, on the contrary, forms a characteristic
opposition to the infinitely variable magnitude. If the constant velocity, $v$, is expressed
% --- p. 113 ---
\setcounter{footnote}{0}%
\opage{113}e.g. by $\frac{s}{t}=a$, then
naturally $s=vt=at$. Since $t$ itself in this expression is conceived as a finite arbitrary magnitude, the addition of
an arbitrary constant would be not merely superfluous, but even
impermissible; for if $\frac{s}{t}=a$, there can impossibly follow from it

\[ s=at+c \]
unless one, by putting $c=0$, makes
the addition a sham fight. Otherwise when $s$ is to be determined by $\int adt=at$; for in this presentation $t$ is not
arbitrary, but infinitely variable, has as yet no finally
determinate value, is, in other words, not in being, but
in becoming. To set a limit to this becoming by stating
a closing point, without at the same time stating a corresponding
beginning point, is as insufficient as it is, in stating
how far someone has travelled, insufficient to say
that he has travelled to Copenhagen, if one has forgotten
to say where his journey started from; the starting point of the
becoming, of the change,
in our example: of the motion, is stated
by the constant. In order to determine $s$ by $\int adt$, the expression must
be: $s=\int adt=at+c$, for by putting
$t=0$ one then obtains in $c$ the limit from which the motion
is to be reckoned. If $c$ is finite, then here there is a finite
distance between a previously given point and the beginning point
fixed for the motion; if $c=0$, these two points
coincide.

The limit determinations necessary for the determining integration
form, moreover, a characteristic opposition
to the limit determinations that arise by differentiation.
For $f(x)=x^{2}$, $d.f(x)=d.x^{2}=f'(x)dx=2xdx$ is an essential limit; for, geometrically viewed, $xdx$ as moment is a line, and the line is the limit of the surface: the
limit to which differentiation leads is, as we have
seen, in relation to what is limited a moment, a rudiment,
% --- p. 114 ---
\setcounter{footnote}{0}%
\opage{114}an element. The limit by which the integral is determined,
on the other hand, when it is not held fast in pure abstraction
as $0$ or $\infty$, is never a mere moment or a mere
possibility, but always an existing finite magnitude, and in
this its finitude coordinate with the integral determined by it.
Thus, since $\int dx=x+c$, $\int 2x\,dx=x^{2}+c$, $\int 3x^{2}dx=x^{3}+c$, $c$ is a line when $x$ is
a line, an area when $x^{2}$ is an area, a volume when $x^{3}$ is a volume. The added constant, which, as regards
value, is arbitrary, is therefore, as regards the kind and nature
of the magnitude, by no means arbitrary. That
the constant must be homogeneous with the integral is seen expressly
when it is conceived as a function of the variable of the integral,
belonging to the determinateness of the integral, and thereby
shows itself as an essential part of the integral itself. If $\int 2x\,dx$ is denoted by $\int f'(x)dx+c=f(x=b)-f(x=a)$
\[ =\int_{x=a}^{x=b}2xdx=b^{2}-a^{2} \]
; then $c=-f(x=a)=-f(a)$
has given up its immediacy, and, as a limit of the function, is transformed
into a reflected determination of limit.

If now the determination of limit is reflected in the function, then
it is also seen from this that, instead of putting
\[ \int_{x=a}^{x=b}2\,x\,dx \]

\[ =\int 2\,b\,dx-\int 2a\,dx \]
, which here, since $x$
is variable, would
be altogether meaningless, one must first carry out the integration,
in order afterwards in the expression $\int 2x\,dx=x^{2}$ to be able to ascribe to $x$

the arbitrary values $b$ and $a$. Inasmuch as the finitely
determined magnitude is then presented in this way as a
definite integral, the reciprocal determination of finite and infinite is
articulated in objective reflection. In and with every finite
magnitude the infinite is therefore posited; for every finite
% --- p. 115 ---
\setcounter{footnote}{0}%
\opage{115}magnitude, when seen in the concept, is a determination
of the infinite, a definitive expression for $\infty. 0.$

When the integral is determined in such a way that the limits
reflected in it must themselves be regarded as integrals, the
problem of infinity may be said to be solved; for whereas the differential
shows the origin of the infinite from the finite,
the integral conversely shows the origin of the finite from the
infinite. Only the analysis of the infinite, as mathematics
carries it through, can therefore lead to the solution of the problem;
but the solution itself is, essentially viewed, not
mathematical: it is not a difficult calculation, not an intricate
mechanics of operations, but a dialectical development of concepts,
on which in the last instance everything depends.

\begin{center}Conclusion.\end{center}

If the solution of the problem of infinity given in this work
can to any degree stand its scientific
test, then it will not be difficult either to demonstrate
that a dialectic into which the exact elements of analysis enter
is of essentially formal, and just thereby
of true and real, interest for philosophy.

Already in the older Greek philosophy, particularly
among the Eleatics, there arise such dialectical problems of quantity
as, after having been worked upon through centuries even
by the most eminent thinkers, nevertheless find their complete solution
only in the present time, namely by being seen from the standpoint of the higher analysis. In
emphasizing in this connection in particular the problem of division and
the problem of motion, we may conclude by recalling
how the various considerations belonging here all in the end
fall under the one comprehensive problem, decisive for philosophy first
% --- p. 116 ---
\setcounter{footnote}{0}%
\opage{116}and last, namely: \emph{the problem of ideality.}

\begin{center}a) The problem of division.\end{center}

The well-known objections that the Eleatic Zeno raised against
the possibility that a plurality in being could be
thought have been summarized and presented by Zeller\footnote[1]{Die Philosophie der Griechen. 1ster Th. Tübingen 1856. p. 425—28.} in the following
turns of phrase:

``If being were a multiple, it would have
to be at one and the same time both infinitely small and infinitely
great. Infinitely small, for since every plurality is a number
of units, but only the indivisible is a real unit,
each of the many must either itself be an indivisible
unit, or consist of such units. But what
is indivisible cannot have any magnitude, for everything that
has a magnitude is divisible to infinity. The individual
parts of which the multiple consists have therefore
no magnitude. There is therefore nothing that can become greater
by the addition of such parts to it, and nothing that can become
smaller by their subtraction from it. But what, when
it is added to another, does not enlarge it, and by being
subtracted from it does not diminish it, is nothing. The multiple
is therefore infinitely small, for each of its constituents
is so small that it is a nothing. But on the other
hand these parts must also be \emph{infinitely great.} For since
that which has no magnitude is not, the many, in order
to be, must have a magnitude, their parts must therefore
be separated from one another. Likewise that which
removes them from one another must also have a magnitude, and that which
in turn removes this from them likewise, and so on
to infinity, so that in consequence of this we obtain infinitely many
magnitudes or an infinity of magnitudes.'' ``By a
similar way of considering, Zeno shows that the multiple, as regards number, must be both limited and
% --- p. 117 ---
\setcounter{footnote}{0}%
\opage{117}unlimited\dots{}. As in the first proof the determination
of the infinite magnitude, so here the determination
of the infinite number is obtained by the multiplicity
being conceived as a plurality of separated magnitudes, and by a third
separating element being interposed
between every pair of separated ones.
For this reason the ancients were accustomed to designate this part of both
proofs as the proof of the dichotomy.''

The point is here grasped so sharply and the knot tied
so firmly that neither a Plato's dialectic nor an Aristotle's
analytic was able to solve the problem. What among
the Greeks must have made the solution of the problem of infinity difficult even from the point of view of general logic, was the well-known
separation, running through the whole of Greek philosophy, of
the limited (πέρας, πεπερασμένον) and the unlimited
(ἄπειρον), a separation according to which the limited was
one with the rational, the unlimited on the contrary with the
irrational, the indeterminable, that which is in itself unthinkable (ἄλογον).\footnote[1]{The Aristotelian view, which (cf. e.g. Phys. III. V. VI.) on the whole amounts to this, that the infinite is not in actuality, but only according to possibility, is a denial of the reality of the infinite: Ὅτι δὲ εἰ μὴ ἔστιν ἄπειρον ἁπλῶς, πολλὰ ἀδύνατα συμβαίνει, δῆλον. τοῦ τε γὰρ χρόνου ἔσται τις ἀρχὴ καὶ τελευτή. καὶ τὰ μεγέθη οὐ διαιρετὰ εἰς μέγεθος, καὶ ἀριθμὸς οὐκ ἔσται ἄπειρος .... καὶ δῆλον ὅτι πῶς μὲν ἐστι, πῶς δ’ οὔ. Λέγεται δὲ τὸ εἶναι, τὸ μὲν δυνάμει, τὸ δὲ ἐντελεχείᾳ... Λείπεται οὖν δυνάμει εἶναι τὸ ἄπειρον. Phys. III. 6. Since the infinite is in this way held in mere possibility, that is, in a possibility that can never become actuality, it becomes an impossible. The unlimited, which among the Greeks signified the lowest, matter, was transformed in medieval speculation into the highest, the infinite, which then in turn became, if not in the oriental sense, an expression of perfect, divine being. As the one who has brought the absolutely infinite unity into relation to number and the multiple, we must here mention the Cardinal Nicolaus of Cusa († 1464); in his three books \textit{De docta ignorantia} he has stated the following view: ``The truly infinite being cannot have any other being outside itself that would be capable of forming a real opposition, for thereby the perfect being would become limited; the infinite being is the absolute self-sameness (absoluta identitas), the form of all forms, the actuality of all possibilities. Since now the absolutely greatest is one, that which is capable of a more or a less can only be a plurality, and thus only be in number. Thus the infinite unity, as the absolutely smallest, is the ground and beginning of all numbers, and, as the absolutely greatest, the end and limit of all numbers.'' Maximum, quo majus esse nequit, simpliciter et absolute, cum majus sit, quam comprehendi per nos possit, quia est veritas infinita: non aliter quam incomprehensibiliter attingimus. Nam cum non sit de natura eorum, quæ excedens admittunt et excessum, super omne id est, quod per nos concipi potest: omnia enim quæcunque sensu, ratione aut intellectu apprehenduntur, inter se et ad invicem taliter differunt, quod nulla est æqualitas præcisa inter illa. Excedit igitur maxima æqualitas, quæ a nullo est alia aut diversa, omnem intellectum: quare maximum absolute, cum sit omne id quod esse potest, est penitus in actu, et sicut non potest majus esse, eadem ratione nec minus, cum sit omne id, quod esse potest. Minimum autem est, quo minus esse non potest. Et quoniam maximum est hujusmodi, manifestum est minimum maximo coincidere. Et hoc tibi clarius fit, si ad quantitatem maximum et minimum contrahis. Maxima enim quantitas est maxime magna. Minima quantitas est maxime parva. Absolve igitur a quantitate maximum et minimum, subtrahendo intellectualiter magnum et parvum: \& clare conspicis, maximum et minimum coincidere. D. Nicolai de Cusa Cardinalis opera. Basiliæ. De docta ignoratia. Lib. I. cap. IV. This identity of a limited-unlimited, an infinite and yet absolute maximum, an infinite and yet absolute minimum, was developed further in the 16th century by G. Bruno.}

% --- p. 118 ---
\setcounter{footnote}{0}%
\opage{118}With Newton's method of fluxions and Leibniz's differential
calculus the modern age makes a considerable advance over
antiquity, in that regard for the infinite,
however one may otherwise conceive it, from now
on shows itself ever more pressing. When the problem of division
passed over into the higher analysis, the dialectical began
to stir in the exact science; how significant in
% the note 1) of p. 117 runs over onto the foot of this page: set in full at its mark
% --- p. 119 ---
\setcounter{footnote}{0}%
\opage{119}this respect are the words of Saverien, cited by Maclaurin:
``\textit{on vit alors pour la premiere fois des disputes parmi
les Géometres}''! The mathematician drives the exact analyses
to their utmost limits, and then sees to his unpleasant
surprise dialectic break forth.\footnote[1]{Nous supposons que la soudivision d'un Etre peut se continuer a l'infini, \& dès que cette soudivision devient telle, que nous ne pouvons plus la saisir, nous regardons comme nul, le reste de cet Etre infiniment petit, quoique nous voulions qu' il y ait encore dans lui une division infinie, \& que nous égalons néanmoins à zero. Cela est admirable: nous voulons concevoir l'infini, \& connoître cet infini par le fini même. Cette contrarieté vient de ce que nous ne savons au juste ce que c'est que l'infini. Dire qu'une chose est infinie, c'est dire que nous n'y connoissons pas de terme; et dès-lors nous devons rejetter tout ce qui sembleroit supposer en nous cette connoissance. Saverien Histoire du Calc. p. XVII.—XIX. Even Lacroix, who in his \textit{Traité du Calcul différentiel et intégral. Tom. I. Paris 1810} does not enter further into what he calls \textit{le changement de métaphysique,} nevertheless cannot refrain from a side-blow at this brilliant concept, this little dialectical expression, \textit{cet être infiniment petit!} Thus in the work cited, p. 242. En partageant avec toute l'Europe, le respect attaché au nom et aux travaux de M. Lagrange, j'oserai néanmoins n'être pas de l'avis de cet homme si justement célèbre, sur les motifs qui paraissent le porter à introduire une nouvelle manière d'écrire les résultats du calcul différentiel; car je crois qu'il est facile de prouver que l'ancienne n'a point en elle même l'inconvénient de rappeler continuellement l'idée fausse des infiniment petits.} The 17th and in part
the 18th century's philosophy was meanwhile, inasmuch as
it was not altogether skeptical, so fixed in dogmatism that
the consciousness of the dialectical, even in thinkers such as
Spinoza and Leibniz, was altogether pushed back. No wonder
then that the concepts of a Spinoza crystallized in exact
definitions\footnote[2]{Cf. Benedict Spinoza, by Dr. H. Brøchner. Kbhn. 1857. p. 44—58 and elsewhere.} no wonder that a Leibniz with all his genius
% --- p. 120 ---
\setcounter{footnote}{0}%
\opage{120}was not able to comprehend as philosopher what he himself
discovered as mathematician; no wonder that a Wolf could
hail the mathematical method as a model for the philosophical,
while a Berkeley\footnote[1]{A. Treatise concerning the Principles of Human Knowledge. London 1734. p. 146. He whose Understanding is prepossest with the Doctrine of \emph{abstract} general Ideas may be persuaded, that (whatever be thought of the Ideas of Sense) Extension in abstract is infinitely divisible. And one who thinks the Objects of Sense exist without the Mind, will perhaps in virtue therefore be brought to admit, that a Line but an Inch long may contain innumerable Parts really existing, though too small to be discerned. These Errors are grafted as well in the Minds of Geometricians, as of other Men, and have a like influence on their Reasonings .... p. 150: Of late the Speculations about Infinites have run so high, and grown to such strange Notions, as have occasioned no small Scruples and Disputes among the Geometers of the present Age. If it be said that several Theoremes undoubtedly true, are discovered by Methods in which Infinitesimals are made use of, which could never have been, if their Existence included a Contradiction in it. I answer, that upon a thorough Examination it will not be found, that in any Instance it is necessary to make use of or conceive infinitesimal Parts of finit Lines, or even Quantities less than the Minimum Sensibile: Nay, it will be evident this is never done, it being impossible. ... Particularly, \emph{Matter or the absolute Existence of Corporeal Objects,} hath been shewn to be that wherein the most avowed and pernicious Enemies of all Knowlegde, whetever humane or divine, have ever placed their chief Strength and Confidence. p. 152 ff.} precisely in the contradiction that
lies in infinite divisibility thought he saw a bottomless
absurdity, a means of explaining away the reality of matter.

Although the dogmatic and skeptical philosophy toward
the close of the 18th century was more and more supplanted
by the critical, as good as nothing was gained with regard to the solution
of the problem, since thinking, in spite of its
altered position toward being, still remained standing at the
% --- p. 121 ---
\setcounter{footnote}{0}%
\opage{121}formal principle of contradiction. Among Kant's cosmological
antinomies the problem of division naturally also plays its
role as insoluble; just as it is, as regards these antinomies,
characteristic on the whole that the insight which Kant
as a mathematician undeniably has into the relation between the
finite and the infinite is not at all present to him,
and thus does not come to his benefit when he philosophizes.
Only in the dialectic prepared by Fichte and developed by Hegel
does a new light begin to dawn upon the infinite\footnote[1]{See H. Schwarz. Work cited.}

As to the problem of division in particular, we must
presumably, in order to grasp its complete solution, preferably take it up from
two main sides, namely the merely formal and the substantial
side; for the fusion of these two points of view,
as it is already found in Zeno, has surely
contributed not a little to confuse the consideration and make the solution
of the task difficult.

Zeno's presupposition that being, if the
multiple is, must have a magnitude, or that ``that which has
no magnitude is not,'' demands in the first place
an abstracting from the substantial, a disregarding of the reality of
being. To be is to be magnitude, and
the contradiction in being is a contradiction in
being magnitude: here is the problem in its abstract form.

The contradiction in the concept of magnitude is divisibility.
This contradiction dissolves in the insight that divisibility,
by being continued to infinity, entails absolutely no
annihilation for any part whatever of the divided
magnitude. Just as $1$ loses nothing by being presented
as $\frac{1}{2}+\frac{1}{2}$ or as $2.\frac{1}{2}$, so a line also loses
not the very least of its magnitude as a whole by being
thought divided into two equal parts, each of these parts again
into two parts, etc. If we let the series $1+\frac{1}{2}+\left(\frac{1}{2}\right)^{2}....$
% --- p. 122 ---
\setcounter{footnote}{0}%
\opage{122}$+\left(\frac{1}{2}\right)^{n}+\ldots$ represent a dichotomy
proceeding to infinity, then the expression $\mathrm{Lim}\left(\frac{1}{2}\right)^{n}$ will denote an element,
an infinitely small magnitude, but not the absolute zero. By
a division continued to infinity the absolute zero cannot
arise; if a continued division were to lead to zero, it would have
to entail both a continued annihilation and at the same time
stop at a last limit, but both assumptions are in
equal degree contrary to infinite divisibility and to the
essence of magnitude.\footnote[1]{That zero cannot be reached by continued division follows also from this, that every new division presupposes that a new limit is set; but the new limit cannot be set immediately together with the one already set; if there is not here an interval, however small, between the limits, they coincide into one limit, and no new division takes place. When the limits move closer and closer to each other, what is limited must undeniably become smaller and smaller; but the limits approaching each other nevertheless do not annihilate the parts, they only move over the parts. If, thus, $y=f(x)$ is the equation of a curve, and $z=\Phi(x)$ the expression for an area bounded by the axis of abscissas, a part of the curve and two ordinates, of which the farthest corresponds to the abscissa $x$; then $\Phi(x+h)-\Phi(x)$ is smaller than the rectangle that corresponds to the greatest ordinate and the increment of the abscissa $h$, but greater than the rectangle that corresponds to the smallest ordinate and $h$. Thus: $h f(x) \lessgtr \Phi(x+h)-\Phi(x) \lessgtr h.f(x+h)$; $h f(x) \lessgtr \Phi'(x) h+\Phi''(x)\frac{h^{2}}{1.2}+\ldots \lessgtr h f(x)+h^{2} f'(x)+\ldots$, and $f(x) \lessgtr \Phi'(x)+\Phi''(x)\frac{h}{1.2}+\ldots \lessgtr f(x)+f'(x)\frac{h}{1}+\ldots$ If one now goes to the limit by putting $h=0$, i.e. lets the two ordinates approach each other until they coincide, the above limits become one with the magnitude lying between them, i.e. $\Phi'(x)=f(x)$. But while the limits coincide, and the area enclosed by them shrinks to $0$, nothing is nevertheless lost, no part of the magnitude is annihilated; only this change takes place, that the parts of the area which before fell inside the separated limits now fall outside the coinciding limits. A change of limits indeed always brings about a change in magnitudes and in the relations of magnitudes; but no change of limits, no displacement, expansion, approach or coinciding of limits of magnitude is able to make any already existing magnitude into nothing.}

% --- p. 123 ---
\setcounter{footnote}{0}%
\opage{123}The giddiness that so easily seizes one at the notion
of a division continuing to infinity
does not properly concern the matter; it is a confusion of mind
which arises from this, that instead of thinking, one
sets about fantasizing, and, fantasizing, wants to represent to oneself that which
cannot be represented. The continued division is, when we
hold fast to thought, only a repetition of the first division:
in the first division lies the whole problem. If a finite
line $x$ is thought divided by a single point into two lines, then the element of the line
$dx$ is at the given point at the same time
divided into two elements: to divide a line without dividing one of
its infinitely small parts is impossible. If now the line $x$ can be divided without anything being lost or annihilated, then the same
holds of $dx$; but if $dx$ is divided without anything being
annihilated, then, eo ipso, $\frac{dx}{2},\ \frac{dx}{3}\ldots\frac{dx}{n}$ is also divided
without any part being annihilated.
To divide the line $x$ is therefore, in other words, to divide the element $dx$, but to divide
the element $dx$ is in turn to divide an element of the element, to
divide $\frac{dx}{n}$; of the divided $\frac{dx}{n}$, $\mathrm{Lim}\,\pm\frac{dx}{n}$ will then denote the moments
which on both sides border on the limit, and
thus by infinite approximation pass over into the $+0$ and $-0$, from which the pure $0$ dialectically distinguishes itself.
% the note 1) of p. 122 runs over onto the foot of this page: set in full at its mark

% --- p. 124 ---
\setcounter{footnote}{0}%
\opage{124}If, after the formal difficulties have been
removed, particular emphasis is laid on the reality of the
existing magnitude, then the problem shows itself from the substantial
side. Zeno's words: ``Only the indivisible is a real
unit, each of the many must therefore either itself be
an indivisible unit or consist of such units,''
do not properly concern magnitude merely as magnitude, but the
real in magnitude, substance. The real unit
then does not suffer first of all from the contradiction of being divisible
to infinity; but the contradiction is that the proper substance
is to be both divisible and indivisible. Inasmuch as
magnitude is substantial, its unity is indivisible; inasmuch as
substance is magnitude, its unity is divisible:
this contradiction stood clear before Leibniz when, in order to
avoid it, he devised the monadology\footnote[1]{Log. et Metaph. p. 5. Necesse autem est dari substantias simplices (vel monades, nam monas non est nisi substantiá simplex), quia dantur composita. Cf. ``En Fremstilling af Leibniz's Monadelære.'' C. F. Nielsen. Kbhvn 1856.} For what is
a monad? A single substance, a substantial magnitude,
a magnitude that is more than a magnitude, an original
unity that cannot be divided. To distinguish the monads from
the mathematical and physical units Leibniz chose to
call them substantial units, metaphysical points. ``\textit{On
les pourroit appeller points métaphysiques\dots{}; ainsi les
points physiques ne sont indivisibles qu'en apparence:
les points mathématiques sont exacts, mais ce ne sont
que des modalités: il n'y a que les points métaphysiques
ou de substance, qui soient exacts et réels; et sans eux
il n'y auroit rien de réel, puisque sans les véritables
unités il n'y auroit point de multitude}''. \textit{Log. et Metph.
p. 53.}

In the monads, therefore, Leibniz means to have the existing
units, which in spite of the force inherent in them
% --- p. 125 ---
\setcounter{footnote}{0}%
\opage{125}bring it about that physical and divisible magnitudes form themselves, nevertheless in
their substantial simplicity are absolutely
indivisible. The question how that which in itself
has no magnitude can be the principle of real magnitudes, how that which is itself without extension can
produce the extended, is then answered to the effect that the monads
as single substances are individual, so that every
monad, in virtue of its individual limitation, has an original
active \textit{(force active primitive)} and an original
passive force \textit{(force passive primitive)}. As the principle of
matter, the passive force has a tendency toward extension,
and thereby causes the original matter \textit{(materia
prima)}, which does not consist \textit{``in extensione''}, but \textit{``in extensionis
exigentia''}, to become an extended and thereby infinitely
divisible magnitude.\footnote[1]{Il est vrai (selon mon système) qu'il n'y a point de portion de la matiere, ou il n'y ait une infinité de corps organiques .... Il n'y a point de partie de la matiére, qui ne soit divisée actuellement. Log. et Metaph. p. 39, 44 and elsewhere.}

What is unsatisfactory in this explanation, so long as
the dialectical is overlooked, is too conspicuous for the contradiction
not to have found its corresponding expression in the critical
philosophy. Under ``the Antinomies of Pure Reason''
Kant also presents, as ``the second conflict of the transcendental ideas,''
the conflict about the indivisible-divisible. Thesis:
\emph{Every composite substance in the world consists
of simple parts, and there exists nowhere anything
except the simple, or that which is composed of it.}
Antithesis: \emph{No composite thing in the world
consists of simple parts, and there exists nowhere
anything simple in the world.}

``If one assumes,'' the thesis is thus proved, ``that the
composite substances did not consist of simple parts, then
when all composition is thought away, there would
% --- p. 126 ---
\setcounter{footnote}{0}%
\opage{126}remain no composite part, and (since there are no simple parts)
no simple one either: one would then be left in the end with only the
pure nothing, and consequently no substance could be given.''
The proof of the antithesis rests on the following fundamental thought:
``The composite thing (substance) must be in space;
but space does not consist of simple parts, but of
smaller spaces; therefore every part of the composite must
occupy a space. In everything composite the properly first
parts are simple; therefore the simple occupies a space. Since
now every real thing that occupies a space comprises a multiplicity
of parts existing outside one another, and is consequently
composite, and, as a real composite, not composed of accidents
(for the accidents, in order to be outside one another, must presuppose
substance), but of substances: the
simple, inasmuch as it is to be in space, would have to be a substantial
composite; which is a contradiction.''\footnote[1]{That Kant in this Antithetic merely conceives the simple as the presupposition for and opposite of the composite, and so far sees the Leibnizian single substance only from its one side, he himself expressly emphasizes in his ``Anmerkung zur zweiten Antinomie.'' (Krit. d. r. V. p. 356) ``Ich rede übrigens hier nur von dem Einfachen, so fern es nothwendig im Zusammengesetzten gegeben ist, indem dieses darin, als in seine Bestandtheile aufgelöset werden kann. Die eigentliche Bedeutung des Wortes \emph{Monas} (nach Leibnizens Gebrauch) sollte wohl nur auf das Einfache gehen, welches \emph{unmittelbar} als einfache Substanz gegeben ist (z. B. im Selbstbewußtseyn) und nicht als Element des Zusammengesetzten, welches man besser den Atomus nennen könnte. Und da ich nur in Ansehung des Zusammengesetzten die einfachen Substanzen, als deren Elemente beweisen will, so könnte ich die Antithese der zweiten Antinomie die transcendentale \emph{Atomistik} nennen. Weil aber dieses Wort schon vorlängst zur Bezeichnung einer besonderen Erklärungsart körperlicher Erscheinungen (molecularum) gebraucht worden, und also empirische Begriffe voraussetzt, so mag er der dialektische Grundsatz der \emph{Monadologie} heißen.''}

% --- p. 127 ---
\setcounter{footnote}{0}%
\opage{127}The question of the relation of the simple to the composite
refers, whether one uses of the simple
the expression ``atomus'' or ``monas,'' originally to
the question of the relation of matter to force, and thereby
to the question of the relation of extensive magnitude
to intensive. It is rational mechanics which, by
the application of the analysis of the infinite, from its exact standpoint
spreads a higher light over the fundamental relations in question,
and thereby supplies a presupposition indispensable for the solution of the problem;
but the solution itself is dialectical.

By holding fast the critical limit, exact
science everywhere lets such opposite concepts as matter
and force, the intensive and the extensive, the simple and
the composite, exclude one another; what is matter
is therefore not force, the extensive is not intensive, the simple
not composite. That force is essentially matter and
matter force, that the extensive in its reality nevertheless
is intensive, and the intensive extensive, that the simple
at bottom can be a composite, and the composite a
simple: such propositions rational mechanics must reject in
the most determinate manner, precisely because it is an exact, not a
dialectical science.

But precisely by its exact analyses rational
mechanics can in a quite peculiar way be said to prepare
and call forth a scientific insight into the dialectical.
For rational mechanics does not stop
at the fact that finite masses attract one another in a finite
manner, but interprets the attraction in such a way
that every smallest part in every mass must, in that it attracts
it is attracted, be at one and the same time both point of departure
and point of application \textit{(point d'application)} for an acting force.
Inasmuch as what is at issue here is not this or that actual thing,
which is always subject to these or those particular conditions,
but the essence of the phenomena, the principle of the actual,
the finite is consistently dissolved into the infinite
% --- p. 128 ---
\setcounter{footnote}{0}%
\opage{128}as into its ground, and this infinite in turn conceived as
the real possibility, which is always realized under given conditions
in a corresponding finite expression.\footnote[1]{Il est évident, ... que la division de la masse en élémens infiniment petits, et la supposition d'une densité de chaque élément, qui ne varie pas dans les corps homogènes, ou qui varie par degrés insensibles dans les corps hétérogènes, ne conviennent point aux corps naturels; mais cela n'empêche pas qu'on ne puisse faire usage des formules fondées sur cette considération ... Poisson Traité de Méchanique. Tom. I. Paris 1833. p. 175. That Poisson by \textit{``corps naturels''} thinks only of things in their finitude, not of the essence of things, is also seen from an expression like this: \emph{Un infiniment petit} est une grandeur moindre que toute grandeur donnée de la même nature. Work cited, p. 14.}

If, to refer to an example, the material
volume $\varrho v = M$ shrinks to an extensive element $dM$, then
the force (weight) $P$ shrinks at the same time to an intensive
element $dP$; a finite force (weight) is therefore only
possible with a corresponding finite volume of the mass. If $p$ denotes the weight of a unit of volume, then the weight of a
volume element, expressed in polar coordinates, becomes $d.P = p r^{2}\,dr\,\sin\alpha\,d\alpha\,d\Theta$. If this equation is now integrated, finite
weight is obtained through finite volume. If it is e.g. a
homogeneous sphere, so that $p$ is thus constant, one obtains,
for $r=a$, by integration:
\[ P=p\int_{0}^{a}\int_{0}^{\pi}\int_{0}^{2\pi} r^{2}\,dr\,\sin\alpha\,d\alpha\,d\Theta \]
\[ =2p\pi\int_{0}^{a}\int_{0}^{\pi} r^{2}\,dr\,\sin\alpha\,d\alpha=4p\pi\int_{0}^{\alpha} r^{2}\,dr=\frac{4}{3}p\pi a^{3}. \]

For any variation of $a$ whatever, even the very smallest, since $p$ is constant, a corresponding variation must take
place in $P$, and conversely; i.e. any change whatever, even the very slightest,
in the material, the extensive, is at the same time
a change in force, the intensive.

% --- p. 129 ---
\setcounter{footnote}{0}%
\opage{129}From the analysis here indicated there follows, as an incontrovertible result, that matter presupposes force in the infinite, just as force in turn presupposes matter in the infinite; that matter therefore can no more be without force than force without matter; that the existence of matter and of force is essentially one and the same being. Since force is possible only by means of matter, it is impossible to derive matter from force; since matter, conversely, is possible only by means of force, it is impossible to derive force from matter; the \emph{dynamic} view, which makes force the principle, is as untenable as the \emph{materialistic} view, which makes matter the principle: the principle is neither force as force nor matter as matter, but the dialectical unity of matter and force.

In the dialectical unity of matter and force is seen at the same time the dialectical unity of the extensive magnitude with the intensive. When the extensive magnitude $M$ shrinks to $dM$ and the intensive magnitude $P$ to $dP$, the two are identical in their disposition; when $dM$ is integrated to $M$ and $dP$ to $P$, the two are identical in their actuality. Here, in the real sphere, as far as the relation between its mass and its weight is concerned, there is not an intensive magnitude by itself, combined with an extensive one, but an intensive magnitude whose reality consists in the extensive, an extensive magnitude whose reality is posited in the intensive. ``One can,'' it is said of the center of gravity,\footnote[1]{Naturlærens mechaniske Deel af H. C. Ørsted. p. 28.}
``imagine things as if the whole weight of the body were united in the center of gravity, and all the remaining parts had no weight.'' Yes, one can certainly ``\emph{picture}'' things in this way; but one cannot \emph{think} things in this way.

In the dialectical unity of the extensive magnitude with the intensive is seen the unity of the indivisible with the divisible, of the simple with the composite. Insofar as the intensive is set over against the extensive in representation, the former will
% --- p. 130 ---
\setcounter{footnote}{0}%
\opage{130}naturally represent the indivisible, the latter the divisible, the former the simple, the latter the composite. That it is just such a represented opposition to which Leibniz and Kant, each in his own way, hold, is easy to show. What Leibniz demands of his ``points metaphysiques,'' namely that they be exact like the mathematical points and real like the physical ones, fits exclusively the points that are in actuality intensive; such an intensive point is now precisely, to stay with the example chosen, the center of gravity of a sphere. When, on the other hand, Kant, after having in the ``Thesis'' made the same demand upon the simple as Leibniz, sets over against it the assertion that no composite thing consists of simple parts, that on the whole nothing simple exists: then he makes in this his ``Antithesis'' a very one-sided use of the extensive magnitude. The misunderstanding in Leibniz and Kant rests, despite the great difference otherwise in their philosophical outlook, essentially on the same thing, as far as the substantial element in the problem of division is concerned. The composite must, on this both are agreed, necessarily presuppose the absolutely simple, the divisible the absolutely indivisible; but this assumption is in accord neither with rational mechanics nor with the nature of things.

In the intensive unity consists real simplicity; but since the intensive, as has been shown above, has its reality in the extensive, and since every finite force is a resultant of components whose material points of application, by being side by side with one another, yield the extensive: the atom itself, insofar as it occupies a space, will have parts, and these parts constitute material points of application for the components whose resultant is the finite force of the atom, its intensive unity. If, however, the atom is nevertheless assumed to be indivisible in itself, an absolutely simple thing: then it is not simple because it has no parts, but because so great an intensity is attributed to the cohesive force of the parts that no
% --- p. 131 ---
\setcounter{footnote}{0}%
\opage{131}power in existence is able to separate them; in other words: the atom is then an absolute indivisible, insofar as the union of its parts is absolute, the atom is an absolutely simple thing, insofar as it is an absolutely composite thing; but an absolutely composite thing is a contradiction, and an absolute atom a fiction. The notion that an absolutely simple thing must necessarily lie at the ground of the composite is therefore also a pure imagination; for what, then, is the absolutely simple? That which, according to its concept, does not enter, and cannot enter, into any finite composition! That rational mechanics can decompose force in the infinite just as exactly as it divides mass in the infinite, while it is nevertheless the intensive in force that gives simplicity its reality, shows, what moreover lies in the nature of the case, that, just as every finite composition is relative, so every corresponding simplicity must also be relative. On the relativity of the simple rests the reality of the composite; but the relativity of the simple rests in turn, when it is seen from the standpoint of magnitude and in particular of extensive magnitude, on its infinite divisibility. The composite has therefore, in accordance with the natural transitions, its reality in the simple, in that the simple in turn has its reality in the composite: the infinite dialectic of the simple and the composite is the solution of the problem of division.

\begin{center}b) The Problem of Motion.\end{center}

How intimately this problem is connected with the foregoing is seen from the objections of the Eleatics pertaining to it. Against the possibility of a motion in space Zeno set up four proofs, all of which are grounded on the infinite divisibility of space. ``If space, or the matter that fills space, is infinite, so that no part is the smallest, then motion is impossible, because it presupposes that a body advances from one place to another, and thus traverses a space. But this is impossible, since every space
% --- p. 132 ---
\setcounter{footnote}{0}%
\opage{132}is infinite, and consequently cannot be traversed in a given finite time. Or, more briefly: the infinity of the space that is to be traversed is in conflict with the finitude of time, without which no motion is possible.''\footnote[1]{Tennemann. Geschichte der Philos. vol. 1. Leipzig 1798. p. 197.}

The care with which Aristotle in his physical writings seeks to disarm Zeno's objections gives us not only a testimony to his acuteness, but also reminds us of the sting that is hidden in the Zenonian arguments. Zeno's first proof, which has already been taken into account above (p. 35), can be presented thus: ``When a body moves, it traverses a line from one point to the other: before it can arrive at the end point, it must therefore first have arrived at the midpoint. Now since every space is divisible in the infinite, every one of the parts of the line consists just as well of infinitely many parts as the whole line does; the body can therefore in its motion never reach the midpoint of the line, let alone its end point, i.e., motion cannot be thought at all.''

Insofar as the nerve of the proof is placed in this, that a finite space, which by being divided in the infinite becomes infinite, would, in order to be traversed, require not a finite but an infinite time, Aristotle has aptly remarked that time is just as divisible in the infinite as space, and that Zeno has erred precisely in having regarded time as limited otherwise than space\footnote[2]{Ζήνων οὖν παραλογιζόμενος, ἄλλως τὸ τοῦ χρόνου πεπερασμένον λαμβάνει, καὶ ἄλλως τὸ τοῦ σταδίου . . . Cf. Tennemann. Op. cit., vol. 3, p. 151.} In this, however, and certainly with full right, the objection has been raised that, ``even if Zeno has erred in regarding time in another way than space, Aristotle, by exposing this formal weakness of the proof, has nevertheless in reality not shaken the proof itself; the difficulty that a motion in
% --- p. 133 ---
\setcounter{footnote}{0}%
\opage{133}space can be thought only as a transition from point to point losing itself in the infinite, that the body that was to move would therefore first have to be finished with an infinity of infinitely small motions,\footnote[1]{„Weil aber diese Theile unendlich sind, so muß er sich eine endlose Zeit bewegen, ehe er sie durchläuft, er kommt also nie aus den letzten.'' Tennemann. Op. cit., vol. 3, p. 152.} before a finite motion could be attained, still stands, even if one thinks of time as divided in the same manner as space.'' This difficulty, which is moreover not removed by a Kantian separation ``zwischen Ding an sich und Erscheinung,'' belongs among the striking examples that show the necessity of letting exact analysis furnish the presuppositions of the development of concepts. The dazzling element in Zeno's proof corresponds to an illusion that runs through the Kantian\footnote[2]{„Denn die Bewegung des Räumlichen ist so gut als eine Theilung, weil sie nicht anders kann gedacht werden, als wie ein Fortrücken von einer Stelle zur andern. Und da nun Aristoteles das Wesen jeder stetigen Größe darin setzt, daß jeder Theil derselben wieder aus Theilen bestehet (nicht etwa von uns blos so vorgestellt wird), so muß das Bewegliche unendliche Theile durchlaufen, ehe es von einer Stelle zur andern kommen kann; das ist aber so viel als, es kann nie von einer Stelle zur andern übergehen. Und gerade das ist es, was Zeno wollte.'' Tennemann. Op. cit., vol. 3, p. 153.}
as well as through the Aristotelian thinking, the illusion namely that the infinite should fall alongside the finite, be its negation without being its affirmation, that one should therefore, insofar as one was in an infinity, be situated outside finitude. That this presupposition is utterly false has been demonstrated by the analysis of the infinite, in that it has demonstrated that the finite magnitude, essentially conceived as a definite integral, is precisely a determination of $0 . \infty$, that one therefore, by moving in the definite infinity, moves in the midst of a finitude. Insofar as the finite is measured by the finite, motion is an uninterrupted succession, a
% --- p. 134 ---
\setcounter{footnote}{0}%
\opage{134}steady transition from a lesser to a greater: first the lesser, then the greater, first an inch, then a quarter of an ell, first half the way, then the whole; but such a succession, which in its continuity is possible only by means of the dialectical unity of the finite with the infinite, does not hold of the relation between infinite and finite with one another: motion does not require first an infinite, for in the positing of an infinite the finite is already presupposed\footnote[1]{``In what sense the finite can be said to be simultaneous with the infinite, rational mechanics makes clear from many places. ``D'Alembert's principle is applied not only to continuously acting forces, comparable with gravity, but also to the so-called \emph{instantaneous} forces, which produce sudden changes in velocity, e.g., in the collision of solid bodies. This second kind of forces is, strictly speaking, not different from the first; for in nature no change in velocity will occur suddenly, but the velocity must, even if in an ever so small and imperceptible interval of time, continuously traverse all values from the original one, which was present when the force began to act, up to the final velocity, which is attained at the same time as the force ceases. The measure of these forces is consequently, on account of the same principle by which e.g. gravity is measured, namely according to the formula $g=\frac{v}{t}$, which also holds for a variable force $\Phi$ with regard to its mean value in a certain time $t$, namely $\frac{\int_{0}^{t}\Phi dt}{t}=\frac{v}{t}$.'' C. Ramus. Analytisk Mechanik. Kbhvn. 1852. p. 211. We come then, according to this way of viewing the matter, to distinguish between: $\alpha$) an absolute Now, the abstract-logical moment, the limit between past and future, a zero that cannot be integrated, $\beta$) a relative Now, the mathematical moment, the time element $dt$, which can be integrated, and $\gamma$) an empirical moment, the real Now of the impact, the imperceptibly small portion of time. In the absolute Now is nothing, and more nothings. In the relative Now the infinite is posited as the beginning of the finite, its element and disposition, and at the same time also as the beginning of an infinity: the beginning of the finite coincides therefore with, is in the mathematical moment simultaneous with, the beginning of the infinite. In the empirical moment the finite is posited as completed; but in and with this completion of the finite an infinity itself is completed: the completion of the finite coincides therefore with, is in every real moment simultaneous with, the completion of an infinity.} Had motion, as Zeno % print: nemlg (misprint)
% --- p. 135 ---
\setcounter{footnote}{0}%
\opage{135}pretends, exhaust the infinite before it could reach the finite: yes, it would take an eternal time, it would be an infinite dawdling! but when motion begins just as originally with the finite as with the infinite, when the finite in motion is always, not after or before, but \emph{simultaneous} with the infinite, because by presupposing the infinite it is itself the presupposition of the infinite: then the dam is broken, then motion has free course, and Zeno is disarmed.

Against Zeno's second proof (Achilles) Aristotle in essence, apart from a slight alteration, brings forward the same objection that he had already used against the first. The proof, which is to establish that Achilles cannot reach the tortoise, i.e., that a body assumed to move with great speed cannot overtake another body that is thought to move at some distance with even a very slight speed, has, like the preceding one, the advantage that it keeps to the punctual. Suppose, e.g., that two bodies, $A$ and $B$, move on the same line, in such a way that $A$, which has a speed ten times greater than $B$, traverses the first half of the line while $B$ traverses the tenth part of the second half; before $A$ can reach $B$, it must then also traverse this tenth part, but while $A$ traverses the said tenth part, $B$ in turn traverses a tenth part of the following tenth part, etc. Although $B.s$ lead becomes smaller and smaller each time, it nevertheless gains by its motion a small lead; this reasoning can now be continued in the infinite, and in the infinite give us a tenth part of a tenth part of a tenth part:
% the note 1) of p. 134 runs over onto the foot of this page: set in full at its mark
% --- p. 136 ---
\setcounter{footnote}{0}%
\opage{136}fancy is at work, representation is confused, and thought grows dizzy. Against so punctual a presentation general-logical remarks are without effect: to say a few fitting words about continuity and discretion, about steady diminution, etc., does not suffice; if the dazzle is to be dispelled, the punctual must be set against the punctual, and the exact analysis brought in. If we choose, then, as unit of time the time that $A$ will take to traverse the distance $l$, the distance between the starting point of $A$ and of $B$, and call the starting point of $A$ $u$, the following series are formed:

At the time $0$ $A$ has a distance from $u=0$, $B$ a distance from $u=l$; at the time $1$ the distance of $A$ from $u=l$, the distance of $B$ from $u=(1+\frac{1}{10})l$; at the time $1+\frac{1}{10}$ the distance is for $A=(1+\frac{1}{10})l$, for $B=(1+\frac{1}{10}+(\frac{1}{10})^{2})l$, at the time $1+\frac{1}{10}+(\frac{1}{10})^{2}$ for $A=(1+\frac{1}{10}+(\frac{1}{10})^{2})l$, for $B=(1+\frac{1}{10}+(\frac{1}{10})^{2}+(\frac{1}{10})^{3})l$; at the time $1+\frac{1}{10}\ldots+(\frac{1}{10})^{n}$ for $A=(1+\frac{1}{10}\ldots+(\frac{1}{10})^{n})l$, for $B=(1+\frac{1}{10}\ldots+(\frac{1}{10})^{n}+(\frac{1}{10})^{n+1})l$. If it is now asked: when will $A$ reach $B$?
then it follows of itself that this will be able to happen only when the following equation is satisfied: $(1+\frac{1}{10}+(\frac{1}{10})^{2}\ldots\ldots+(\frac{1}{10})^{n})l=(1+\frac{1}{10}+(\frac{1}{10})^{2}\ldots\ldots+(\frac{1}{10})^{n}+(\frac{1}{10})^{n+1})l$, that is, when $(\frac{1}{10})^{n+1}=0$, thus when $n=\infty$. % print: tilstillet (misprint)
The time corresponding to this: $1+\frac{1}{10}+(\frac{1}{10})^{2}\ldots+\mathrm{Lim}(\frac{1}{10})^{n}=\frac{1}{1-\frac{1}{10}}=\frac{10}{9}=1\frac{1}{9}$ is, as one sees, finite, and $A$ therefore reaches $B$ in a finite time. If the chosen unit of time is e.g. $1$ sec., $A$ will need $1\frac{1}{9}$ sec. to reach $B$; the infinitely many smaller and smaller intervals of space, which correspond to the infinitely many smaller and smaller portions of time, will then in the course of $1\frac{1}{9}$ sec., in that they result point by point from the finite as its elements, give the finitely determinate result.

The cognition of the manner in which the infinite here is absorbed in the finite rests on the proposition that a series with innumerable terms can yield a finite sum, and is
% --- p. 137 ---
\setcounter{footnote}{0}%
\opage{137}to that extent conditioned by insight into the convergence of infinite series. That the limit unattainable in abstracto is attained precisely by this, that an infinite series is carried to its end not merely in thought but in actuality, that an infinity is actually traversed in and with the finite, is therefore the dialectic with which the illusion ceases.

The third proof (the flying arrow) is held in a more general form, and can be set forth thus\footnote[1]{Tennemann. Op. cit., vol. 1, pp. 198—199.}: ``Were a body to move, it would follow from this that it must at one and the same time be in motion and at rest; for every moving body must at every moment be in a space that is exactly as large as itself. But in this way the body is at rest; for what, even if only for a moment, is situated in a corresponding equal space, is at rest with respect to this space. If we imagine an arrow that with the greatest swiftness flies through a space, then it must at every moment be in a different space, and insofar as it occupies this at every moment be at rest. Such a flying arrow would then have to be thought to be at one and the same time both at rest and in motion, which is a contradiction.''

What Aristotle has developed on this occasion concerning the continuous magnitudes, space and time and the infinite divisibility of both\footnote[2]{Physic. \textit{VI, 1.}}, concerning the difference between the limit of time, the Now, and the element of time\footnote[3]{Physic. \textit{VI, 3.} „Die Unmöglichkeit aber, daß in dem Augenblicke als solchem etwas sich bewege, oder, als Gegensatz der Bewegung, ruhe, ergiebt sich aus den Begriffen der Schnelligkeit und Langsamkeit, welcher Gegensatz als wesentlich jeder Bewegung zukommend bereits vorhin gesetzt worden war; sodann aber auch aus der Definition des Augenblicks selbst, in welchem, dafern er zugleich der Vergangenheit und der Zukunft angehört, so oft die Ruhe in Bewegung, und umgekehrt, übergeht, Ruhe, und Bewegung in Bezug auf dasselbe Ding zugleich statt finden müßte, was widersinnig wäre.'' C. H. Weisse. Aristoteles Physik. Leipzig 1829. Anmerkung zum 6ten Buche. p. 593.}, concerning
% --- p. 138 ---
\setcounter{footnote}{0}%
\opage{138}the difference between a rest that excludes motion and that infinite hindrance of rest which constitutes motion itself\footnote[1]{Characteristic above all is the manner in which the 8th chapter (Phys. VI) ends: ``But since everything that moves moves in a time, and passes from something into something (ἐκ τινὸς εἰς τὶ μεταβάλλει), that which moves cannot, in the first place, be at rest in any respect as regards the time in which it moves wholly (ἐν ᾧ χρόνῳ κινεῖται καθ᾽ αὐτό) and not in part (καὶ μὴ τῶν ἐν ἐκείνου τινί). For rest consists in this, that both the thing itself and each of its parts is for some time in one and the same thing. For only then do we speak of rest, when it can truly be said that the thing together with its parts is at different moments one after another (ἐν ἄλλῳ καὶ ἄλλῳ τῷ νῦν) in one and the same place. But when rest consists in this, that which changes cannot be wholly and entirely in anything as regards the first time. For all time is divisible: so that one could truly say that the changing thing was in the different parts of time in the same place, itself as well as its parts. If, on the other hand, this is not so, and that which moves is only in a single Now (ἐν ἑνὶ μόνῳ τῷ νῦν) here or there: then it will not as regards any time whatever, but only as regards a limit of time, be in something. In every Now it will no doubt always find itself in something, but not in such a way that it rests (ἐν δὲ τῷ νῦν ἔστι μὲν ἀεὶ κατά τι μένον οὐ μέντοι ἠρεμεῖ). In the pure Now neither motion nor rest can take place, but instead non-motion and non-rest. But in time that which moves cannot behave in the same way as that which rests; for then it would follow that what moves in space also rested.''}, certainly belongs among the most significant things that have so far been achieved in the philosophy of nature; just as it must indeed follow from the nature of the case that the category that in Aristotle is the fundamental category of physics, namely motion (κίνησις), would in so many-sided and persevering a thinker find a many-sided and, as far as possible, exhaustive treatment.

But although the Aristotelian acuteness penetrates into the matter from many sides, a satisfactory
% the note 3) of p. 137 runs over onto the foot of this page: set in full at its mark
% --- p. 139 ---
\setcounter{footnote}{0}%
\opage{139}solution of the knot has nevertheless been impossible, and for a twofold reason. In the first place the lack of exact analyses is striking; what is adduced here and there to make clear the relation between change of place and change of time, the faster and the slower, and the like, cannot, since reflection is not supported by a corresponding operation, get beyond logical generality, and the relations are therefore also not articulated to such a degree that a punctual agreement of % print: punklig (misprint)
concept and actuality could show itself; the distinction between the element of time and the pure Now is like a flight of the reasoning, a subtle turn, a flowery phrase left in obscurity\footnote[1]{„Es ist schwer sich einen bestimmten Begriff von dem, was Aristoteles eigentlich unter diesen Augenblicken gedacht hat, zu bilden. Keine Theile der Zeit und doch Grenzpuncte? Sind es nur Ruhepuncte für das Vorstellende, wenn es eine Reihe von Bewegungen durchläuft, und sie vorstellt. So scheint es. Wenn er aber behauptet, daß die Grenze nur dem angehöre, dessen Grenze es ist (τὰ μὲν γὰρ πέρατα ἐκείνου μόνον ἐστίν, οὐ ἐστι πέρατα. \textit{Physic. IV. c. 11}), so muß man annehmen, daß auch der Augenblick etwas in der Zeit sey. Aber dieses läßt sich selbst nicht mit dem Begriff des Aristoteles von der Zeit vereinigen, daß sie ein Ganzes, dessen Theil jedesmal ein Zeittheil und ins Unendliche in Zeittheile getheilt werde, eine Folgerung, die selbst dann noch gültig ist, wenn er unter den Augenblicken bloß imaginäre Abschnitte sich gedacht hätte. Tennemann. Gescht. d. Phil. 3rd vol. pp. 154—155.}. % print: Physie. (misprint)

The second, equally essential hindrance is the lack of thoroughgoing dialectic. Aristotle is indeed too much a thinker and too much a Greek not to come upon the dialectical everywhere in the analysis of concepts\footnote[2]{„Uebrigens sind die Ergebnisse dieses Capitels'' (the reference is to \textit{Physic. VI, 6)} „ein Beispiel, wie die scharf und kräftig durchgeführte aristotelische Verfahrweise der abstract-verständigen Sonderung und Scheidung zuletzt auf ächt speculative Begriffbestimmungen führt: wie denn dieser Gegensatz von Stetigkeit der Veränderung und Plötzlichkeit — (Platon soll, wie Simplicius be- richtet, den Ausdruck ἐξαίφνης für den Begriff des Augenblicks gebraucht haben) — des durch die Veränderung hervorgerufenen seienden Zustandes, eine Knotenlinie von zeitlich-räumlichen oder andern zwischen Qualität und Quantität oscillirenden Maßverhältnissen, ganz unstreitig eine ächt speculative und dialektische, über die gemeine Reflexion weit hinausgehende und an den wichtigsten Folgen reiche Bestimmung ist.'' C. H. Weisse. Op. cit. p. 602.}; but the matter is,
% --- p. 140 ---
\setcounter{footnote}{0}%
\opage{140}that Aristotle, instead of holding fast to Plato's divine thought that philosophy is scientifically dialectic, and carrying the thought further, on the contrary, as far as this point is concerned, falls away from the Idea, reduces dialectic to something subordinate, and must for this reason, where it pinches, help himself by finite reasoning back and forth. Nowhere in the whole Aristotelian physics is the lack of the assistance of an exact analysis perhaps more palpable than precisely in those sections where Aristotle makes use of the ``moment,'' and aims at the Zenonian sphinx; and nowhere is the lack of a thoroughgoing dialectic more palpable than in the reflection with which the dogmatic thinker evades contradiction, and so goes to meet ``the flying arrow.''

Although motion is by no means, as Zeno has presented the matter, the contradiction that that which moves should at one and the same time be in a motion excluding all rest and a rest excluding all motion, a contradiction that in its insolubility would be an unmistakable sign of nonsense; Zeno's proof nevertheless demonstrates that in every motion there is, however much Aristotle may seek to conceal it and to wriggle out of it, a contradiction that is posited and resolved in the infinite; that the category motion, in other words, like the categories becoming\footnote[1]{``Becoming is namely the relation that the Absolute itself is. Therefore it alternates between being and non-being, and is itself the continued series of the succession of both; for what becomes or comes to be belongs partly to what is, insofar as it already is something, partly to non-being, insofar as it has a future that is not yet fulfilled; and both moments form a continued series, because their alternation is continuous, for only when the consideration halts at some determinate point does that which becomes show itself as a separation of the moments, in that only then the one part of that which becomes is comprehended under that which is, and the other under non-being.'' J. L. Heiberg. Perseus. No. 2. p. 23.} and
% the note 2) of p. 139 runs over onto the foot of this page: set in full at its mark
% --- p. 141 ---
\setcounter{footnote}{0}%
\opage{141}change, is infinitely dialectical. In order, however, that this dialectic, by being detached from the punctual, shall not take on the appearance of a play of thought, in which the real motion is transformed into empty motion of words, we shall secure the development by exact presuppositions, and for this reason first take along the Aristotelian treatment of Zeno's fourth proof.

The fourth proof, by which Zeno attacks motion, is, if its transmitted form is otherwise authentic, so conspicuous a sophism, or perhaps rather: paralogism, that Aristotle did not so much need to refute it as to dismiss it, which he has also done\footnote[1]{``But the fourth is the one about the equal masses that move in the track toward one another from opposite sides, the one from the end of the track, the other from its middle: from which it is supposed to follow that the double time becomes equal to the half. But the fallacy (ὁ παραλογισμός) lies in the demand that the one, in relation to that which is at rest, should in equal time with equal speed traverse an equal distance. This is false (τοῦτο δὲ ψεῦδος).'' Physic. \textit{VI, 9.}} The tendency in the proof is nonetheless decisive enough. If it could really be proved that every difference of speed vanishes when motion is thought, then the unthinkability of motion would have been proved at the same time. The contradiction involved here Aristotle has hardly become aware of: he presupposes without further ado a difference of fast and slow, in that he asserts motion, and then presupposes motion in turn, in that he asserts the difference between fast and slow\footnote[2]{Cf. \textit{Phys. VI, 1, 2, 3 and VI, 9.}}

Only through an insight into the exact analysis does it become possible to see what lies in the problem. Let e.g. the motion on two tracks, $A$ and $B$, take place in such a way that
% the note 1) of p. 140 runs over onto the foot of this page: set in full at its mark
% --- p. 142 ---
\setcounter{footnote}{0}%
\opage{142}while on $A$ one mile is covered in an hour, on $B$ eight miles are covered in the same unit of time. A slight acquaintance with the first propositions of the differential and integral calculus is then enough to understand that when on the track $A$ in finite time the velocity
\[ v=\frac{s}{t}=1 \]
, then also the instantaneous velocity
\[ \frac{ds}{dt}=1 \]
, and that on the track $B$, where the velocity

\[ v_{1}=\frac{s_{1}}{t}=8 \]
in finite time, the instantaneous velocity must likewise be
\[ v_{1}=\frac{ds_{1}}{dt}=8 \]
. But as easy as it now is, after the analysis of the infinite has once been introduced, to find one's way with the operational mechanics belonging here --- it goes without headaches with a little routine, as of itself ---, so impossible has it been, and still is, for the reflection unacquainted with analysis to master the task and to defend itself against the sophistic that, by using the ``moment'' to level every difference in velocity, will abolish velocity itself, and make motion meaningless.

In the infinitely small portion of time, if one keeps to a general reasoning, the stretch of way that is traversed on $B$ is just as infinitely small as that which is traversed on $A$, $ds_{1}$ thus just as infinitely small as $ds$. ``But if $ds<ds_{1}$, how can $ds_{1}$ then be infinitely small? How is it possible, how is it conceivable, that something should be able to be still smaller than that which is already infinitely small? How is it possible, how is it conceivable, that something should be greater than that which is already infinitely great? The mathematician says so! Yes, but the mathematician is not in earnest about the infinite; he posits in thought either an empty negation or else only the finite. The mathematician designates the infinite; but he does not think the thought of infinity: it is on the whole not so much concepts as
% --- p. 143 ---
\setcounter{footnote}{0}%
\opage{143}designations that he is concerned with, he thinks only of getting his designations so arranged that the calculation will come out right.''

Once reflection has been bewitched into the fixed notion that the infinite must absolutely be equal to itself: the opposition between velocity in finite time and instantaneous velocity will easily become so hard that it can appear as a Kantian antinomy. Thesis: There \emph{must necessarily} be different instantaneous velocities, since in actuality there are different velocities. Antithesis: Since every difference in the pure instantaneous velocity is unthinkable, the so-called actual differences of velocity are impossible. In order to prove the Thesis one needs only to start from the finite difference between
\[ \frac{s}{t} \]
and
\[ \frac{s_{1}}{t} \]
, and
infer from it a corresponding difference between
\[ \frac{ds}{dt} \]
and
\[ \frac{ds_{1}}{dt} \]
. In order to prove the Antithesis one begins by positing
\[ ds=ds_{1} \]
thus
\[ \frac{ds}{dt}=\frac{ds_{1}}{dt} \]
, and then draws from this the conclusion: % print: Slntningen (misprint)
\[ \frac{s}{t}=\frac{s_{1}}{t} \]
. That the one infinitely small magnitude can be, not merely greater, but even infinitely many times greater than the other, the one infinitely great magnitude, not merely smaller, but even infinitely many times smaller than the other, is indeed a paradox so hard and offensive that the so-called common sense would soon be done with it, were it not that exact analysis from so many sides found itself called upon to impress it, and thereby came to uphold its standing. But ``the sound understanding,'' which is so agreed with itself in rejecting the view that there could be a difference in the infinitely small and the infinitely great, and that the infinite on this account could not enter under the concept: magnitude, % print: Meneskeforstand (misprint)
% --- p. 144 ---
\setcounter{footnote}{0}%
\opage{144}is at the same time agreed with itself in rejecting the assertion that there should not be a real difference between slower and faster motion. That the reasoning understanding, by being thus agreed with itself, is precisely at variance with itself, it has in its immediate complacency of course no notion; but that the matter nevertheless stands thus is easy to show. The difference between slow and fast motion must evidently rest on this, that the first has more, the other less rest in it, that the first is therefore nearer to rest than the other; a motion with infinitely small velocity is indeed so infinitely near to rest that it can be regarded as rest. Nonetheless the opposition of rest and motion is qualitative: what rests does not move, what moves does not rest. If we now cast away the infinite magnitudes, the difference of velocity can be explained only by this, that the pure motion is finitely broken off by rest. That the body on $A$ in the unit of time $t$ traverses only $s$, while the body on $B$ traverses $s_{1}=8s$, must therefore be understood in this way, that the former body moves only $i \frac{1}{8} t$, but on the other hand rests $i \frac{7}{8} t$. So long as the unit of time in which such an alternation takes place is 1 hour, the alternation of rest with motion is conspicuous; if the unit of time is 1 second, the alternation is less striking, and looks like a jerky motion; in a unit of time of $\frac{1}{100}$ second the alternation will not be noticed at all, but the phenomenon will show itself as a smooth motion. That on $A$ only $1$ mile is traversed, while on $B$ $8$ miles are traversed, will then, when the alternation of rest with motion occurs in every hundredth of a second or in the time $\tau$, look to intuition as if the body on $A$, in spite of its lesser velocity, were just as much in smooth and uninterrupted motion as the body on $B$; while the explanation nevertheless is that the motion on $A$ is incessantly interrupted
% --- p. 145 ---
\setcounter{footnote}{0}%
\opage{145}by rest, in that here to each motion in the portion of time $\frac{1}{8}\tau$ there corresponds a rest in the portion of time $\frac{7}{8}\tau$. If the motion on $A$ is to be continuous, time and space must be divisible in the infinite; on this point Aristotle has seen rightly. But what Aristotle has not seen and could not see, because he lacked the point of support that only the analysis of the infinite can give, is that here, by this infinite divisibility of time and space, there precisely form themselves infinitely small magnitudes with such a mutual difference that they enter into determinate finite relations with one another.

Zeno's objection, that the body which was to be in motion would have to be at one and the same time both at rest and in motion, thus has the truth in it that every motion presupposes a rest that is incessantly on the point of setting in, and that is constantly sublated. It is evident in this connection that Hegel is right when he lists ``sublation'' among the dialectical categories. Rest is, to be sure, also a sublated motion, but this sublation is undialectical insofar as rest, seen finitely, does not have the sublated motion in itself, but outside itself; in motion, on the other hand, the sublation of rest, which takes place in every element of time, is dialectical, for the sublated rest is present in the motion itself, so that the difference between greater and lesser velocity rests precisely on the greater or lesser energy with which rest is sublated: the extensive is absorbed here too in the intensive. With the dialectical sublation of rest, Zeno's objections against motion are then also sublated.

Should anyone, however, wish to go further, and demand the problem solved in such a way that motion in space were grasped and held fast in a concept that excluded intuition: such a demand might well be said to rest on a pure misunderstanding. Space and time are, to recall Kant, not pure concepts of the understanding, but a priori intuitions. To wish to
% --- p. 146 ---
\setcounter{footnote}{0}%
\opage{146}transform intuition into an abstract concept of the understanding, to wish to grasp that which is only for intuition in a concept that annihilates intuition, is of course meaningless; such questions as: can space be thought? can time be thought? can motion be thought? must then in this sense be answered straight out with No! Had Zeno without further ado set up as a demand that whoever wished to think motion should think it without intuition, and outside intuition: then the contradiction lying in such a demand would surely have been too conspicuous to furnish a problem running through centuries. No, Zeno presupposes that space, time and motion are given as intuitions, and demands then that these intuitions shall at the same time let themselves be grasped in corresponding concepts. As long as concept and intuition % print: Auskuelse (misprint)
are in conflict, so that that which presents itself to intuition as an actuality shows itself to thought as a contradiction and thereby as an impossibility, the problem stands unresolved; when, on the other hand, it has been shown from all sides that that in motion which presents itself immediately to intuition also, insofar as it is seen from the standpoint of reflection, satisfies thought, the problem has found its solution.

\begin{center}c) The Problem of Ideality.\end{center}

The dialectic of motion is a dialectic of concept and intuition, and precisely thereby a dialectic of the infinite and the finite. To resolve the finite into the infinite in such a way that it emerges again from it, is, in other words, to demonstrate the ideality of the finite. Since such sciences as astronomy and physics, whose ``ideal principle,'' to use an expression of J. L. Heiberg\footnote[1]{See ``\emph{Perseus},'' Journal for den speculative Idee. No. 1. Til Læseren.},
``is mathematics,'' must necessarily furnish the foundation when the cosmological idea with its real content\footnote[2]{Cf. Phil. Propæd. p. 184.} is to be developed and presented, philosophy must by many-sided
% --- p. 147 ---
\setcounter{footnote}{0}%
\opage{147}contact with these sciences find itself called upon in manifold ways to make use of the higher analysis. The doctrine of the Idea cannot, namely, gain real content by having the empty forms of abstraction filled out, not to say stuffed, with examples and theorems borrowed from astronomy and physics; on the other hand even the most concrete phenomena can be thoroughly illuminated by the light of ideality, when infinity % print: Ueudeligheden (misprint)
gains such power with the finite that a combustion of matter can take place in the fiery air of dialectic
(πῦρ ἀεὶ ζῶον ἁπτόμενον μέτρα καὶ ἀποσβεννύμενον μέτρα).

An idealism that properly concerns itself only with what is real in the finite and material sense, insofar as it everywhere aims at dissolving this real as a configuration of the universal concept, as a moment vanishing in the pure Idea, has at bottom no essential need to engage with the exact analyses. In Hegel's Logic, accordingly, the treatment of the mathematical, relegated to a couple of remarks, stands in a rather loose connection with the dialectical development of the categories of quantity; that the Hegelian method must consequently lead away from the mathematical, in that it leads away from the physical analyses of actuality, is shown by his philosophy of nature, and in particular by his explanation of the law of falling bodies.\footnote[1]{Cf. Phil. Propæd. p. 162—163.}

The philosophy, on the other hand, which in order to know the essence of things and to steel the validity of the Idea must come to grips with facts, and by unraveling the knots of actuality seek the ideal threads, has an irrefutable need not merely of exact analyses, but of real concepts, which originally belong to sciences ``whose ideal principle is mathematics.'' But if the need is essential, then the demand too is irrefutable; to what extent the mathematical is to be taken along into philosophy, and the exact concepts drawn into the stream of dialectic, will therefore not rest on the
% --- p. 148 ---
\setcounter{footnote}{0}%
\opage{148}arbitrariness of the one who philosophizes, but on the course of cognition, on the disposition of totality, on the necessary method of the development of the subject matter\footnote[1]{Cf. Phil. Propæd. p. 111—117.}.

``The history of science,'' says Ramus\footnote[2]{Differential- og Integralregning. Preface.}, ``presents the fact that the greatest discoveries in physics and astronomy, which have intervened so powerfully in the life of mankind, both spiritual and social, have been partly conditioned by, partly most precisely connected with, the discoveries of the most abstract truths of mathematical analysis. If the study of higher mathematics is long and difficult, its results in turn yield a so much finer reward. Fruitless, on the other hand, are the endeavors of those philosophers who would shorten the way by mystical and metaphysical speculations.'' But although mathematics is now certainly opposed to ``the endeavors of the philosophers'' who by ``mystical speculations'' would make both the exact analyses and the observations of empirical science superfluous: mathematics at its present stage is nevertheless the science that, in particular from one main side, is eminently suited to correct ``the speculations,'' and to further scientific metaphysics. To be sure, the atomic theory of chemical physics finds in a certain way its confirmation and the conditions of its consequent development in mathematics; but the notion of eternal atoms hardened in themselves (\textit{indivisibilia
absoluta}) can impossibly be brought into harmony with an analysis that divides matter in the infinite, and transforms the parts of matter into points of function. The proposition that every element is essentially a moment is a principal proposition that proves to what degree the exact analyses must be able to contribute to giving the idealism of actuality its proper a priori foundation, when the concepts that are originally developed in mathematics are grasped and worked up by dialectic.

% --- p. (back matter) ---
% ===== BACK MATTER: Indhold (PDF 163-164) and Rettelser (PDF 165), read at the image; no folios printed =====
\clearpage
\begin{center}{\large Contents.}\end{center}
{\small\noindent [The page numbers are those of the 1859 print, given in the margin.]\par}
{\small\noindent\hfill Page.\\
\par\medskip\centerline{I.}\par\noindent
\par\medskip\centerline{The Concept of Function.}\par\noindent
a) The original requisites of the concept of function \dotfill 4--12.\\
$\alpha$) A constant and a variable magnitude. $\beta$) Two variable\\
magnitudes, a dependent and an independent one. $\gamma$) Fundamental relation\\
between the two variables.\\
b) The function and the law \dotfill 12--20.\\
$\alpha$) The case. $\beta$) The rule. $\gamma$) The law.\\
c) The distinctive presentation of the concept of function \dotfill 20--27.\\
$\alpha$) Language. $\beta$) Thought. $\gamma$) Language and thought.
\par\medskip\centerline{II.}\par\noindent
\par\medskip\centerline{The Problem of Infinity.}\par\noindent
a) The finite in its immediate difference from the\\
infinite \dotfill 28--33.\\
$\alpha$) The infinite outside the finite. $\beta$) The infinite\\
inside the finite. $\gamma$) The dialectical unity of the finite and\\
the infinite.\\
b) The transition of the finite into the infinite: approximation. \dotfill 34--38.\\
$\alpha$) An attainable: approximation of finite to finite.\\
$\beta$) An unattainable: approximation of finite\\
to infinite. $\gamma$) The infinite approximation.\\
c) Approximation in punctual development \dotfill 38--56.\\
$\alpha$) Mathematical-logical transitions. $\alpha\alpha$) The exact application of analysis and\\
synthesis: Operation. $\beta\beta$) Analysis\\
on a synthetic basis: Transformation.\\
$\gamma\gamma$) Synthesis on an analytic basis: Combination.\\
$\beta$) Systems of functions: the inner infinity.\\
$\alpha\alpha$) The fundamental formula of systems of functions: Reduction.\\
$\beta\beta$) Particular systems of functions: Deduction. $\gamma\gamma$) Increments\\
and remainders: Articulation.
\par\medskip\centerline{III.}\par\noindent
\par\medskip\centerline{The Analysis of the Infinite.}\par\noindent
a) The limit of the function: the dialectical contact of the finite\\
with the infinite \dotfill 57--76.\\
$\alpha$) Breaking off of infinite series: the limit as unattainable.\\
$\beta$) The essential closure of series:\\
the limit as attainable. $\gamma$) Infinite magnitudes:\\
the dialectic of limitation.\\
b) The differential: the dialectic of the moment \dotfill 78--97.\\
$\alpha$) Differential and differential coefficient: relational magnitude.\\
$\beta$) The different orders of the differential coefficients: system\\
of relational magnitudes. $\gamma$) Principles of differentiation:\\
the method determined by the consequences of the relations.\\
c) The integral: the element and the finite magnitude \dotfill 98--115.\\
$\alpha$) Integrating summation: moment and element.\\
$\beta$) Principles of integration: the relation of derived functions\\
to the original ones. $\gamma$) The definite integral:\\
the solution of the problem of infinity.
\par\medskip\centerline{Conclusion.}\par\noindent
a) The problem of division \dotfill 116--131.\\
b) The problem of motion \dotfill 131--146.\\
c) The problem of ideality \dotfill 146--148\\
}
\clearpage
\begin{center}{\large Errata.}\end{center}
{\small\noindent
p. 38. line 9 and 14 from top Assymptote, read: Asymptote.\\
- 40. --- 13 --- analythisk --- analytisk.\\
- 43. --- 23 from bottom. retvinklene --- retvinklede.\\
- 45. --- 2 from bottom. $A$ --- $A_1$ and\\
$\frac{1}{a}$ --- $\frac{1}{A_1}$.\\
- 46. note $l(a)^{n-1}$ --- $(l.a)^{n-1}$.\\
- 50. note line 1 --- impiriske --- empiriske.\\
- 55. line 1 from top. Da nemlig --- Er nemlig.\\
- 59. --- 5 --- Reduction --- Reduction,\\
6 --- hianden --- hinanden.\\
- 67. --- 13 --- nte Led --- mte Led.\\
- 75. --- 12 --- $(a - x)\,l^{\frac{1}{a-x}}$ --- $(a - x)\,e^{\frac{1}{a-x}}$.\\
- 81. --- 1 --- bive --- blive.\\
- 90. --- 9 from bottom. Benouilli --- Bernoulli.\\
- 94. --- 9 from top. Differentialcoesficient, read: Differentialcoefficient.\\
- 94. --- 7 from bottom. heeel, read: heel.\\
- 96. --- 9 and 14 from top. $\frac{dx}{1}, \frac{dx^2}{1.2}$ --- $\frac{dx}{1.2}, \frac{dx^2}{1.2.3}$ and\\
10 --- $dx$ --- $\frac{dx}{1.2}$.\\
- 101. --- from top. er Delene --- er Delen.\\
- 105. --- 5 and 4 from bottom. --- Lim $\left(1+\frac{1}{n}\right)^{x^2}$\\
$= x^2$.\\
- 110. --- 3 --- $\int \frac{'x^m\,dx}{l(x)^n}$ --- $\int \frac{x^m\,dx}{(l.x)^n}$.\\
2 --- $\int \frac{'x^m\,dx}{l\,x^{n-1}}$ --- $\int \frac{'x^m\,dx}{(l.x)^{n-1}}$.\\
last --- $\int \frac{'x^m\,dx}{l.x}$ --- $\int \frac{'x^m\,dx}{l.x}$, for n a positive integer,\\
}

\end{document}
